\section{Kato-type vanishing of Selmer groups}\label{sec:pf:kato}

\providecommand{\katoQlb}{\overline{\Q}_\ell}
\providecommand{\katoHf}{H^1_{f}}
\providecommand{\katoO}{\mathcal{O}}
\providecommand{\katoZ}{Z^1}
\providecommand{\katoB}{B^1}
\providecommand{\katoht}{\mathrm{ht}}
\providecommand{\katoadj}{\operatorname{adj}}
\providecommand{\katoeq}[1]{\equiv_{#1}}
\providecommand{\katorep}[1]{[#1]}
\providecommand{\katoN}{\mathrm{N}}

This section proves \cref{thm:kato} and \cref{lem:window:heights}. \Cref{thm:kato} asserts that the
package $\sfname{HypothesesH'}$ of published results (\cref{sec:window:arch}) implies the proposition
$\sfname{KatoTwoW}$ of \cref{def:window:kato}: if $W$ is a two-dimensional $\ell$-adic representation of
$G_\Q$, ordinary at $\ell$ with Hodge--Tate weights $\pm e$, with determinant the quadratic character $\theta$
of an imaginary quadratic field, irreducible on $G_{K(\mu_{\ell^\infty})}$, and such that $W\otimes\varepsilon^{-e}$
is attached to a newform of weight $2e+1$, and if a window condition holds, then the Bloch--Kato Selmer groups
of $W$ and of $W\otimes\theta$ vanish. As explained after \cref{thm:kato}, this is in substance a special case
of \cite[Thm.~14.2]{Kato2004}. The formal proof \cite{LeanSource2026}, which we follow, does not assume Kato's
theorem; it re-derives it by Kato's method. It has three layers.

\emph{Reduction} (\cref{pf:kato:sec:reduction,pf:kato:sec:sph}). The representations $W$ of
$\sfname{KatoTwoW}$, and their twists $W\otimes\theta=W\otimes\det W$, are \emph{admitted}
(\cref{pf:kato:def:admitted}). A local--global compatibility statement at $\ell$, derived from
$\sfname{Row0b}$ (\cref{def:row0b}) and quadratic base change, shows that the window condition forces the
cuspidal representations attached to twists of $W$ to be unramified at $\ell$; this produces an
$\ell$-ordinary newform of level prime to $\ell$ attached to a twist of $W$ (a \emph{source}). It then
suffices to prove the vanishing statement \textup{(V)} of \cref{pf:kato:def:sph}: admitted representations with a
source have vanishing Bloch--Kato Selmer group.

\emph{From an Euler system to vanishing} (\cref{pf:kato:sec:nf,pf:kato:sec:alg,pf:kato:sec:local,pf:kato:sec:global,pf:kato:sec:heights,pf:kato:sec:upgrade}).
After conjugation, $W$ is in a normal form: integral, unramified outside a finite set $S$, upper triangular
at a place above $\ell$. An Euler system for $W$ consists of a bottom class $y$, unramified away from $\ell$ and
non-trivial in the quotient line at $\ell$, together with, for every auxiliary prime $w\equiv1\bmod\ell^n$ and
every cyclic character $\psi$ of conductor $w$ and order $\ell^n$, a class over the field cut out by $\psi$
whose corestriction is the Euler factor at $w$ times $y$. A cochain-level version of Kolyvagin's argument,
using one auxiliary prime at a time, shows that a ``height'' built from $y$ and a strict Selmer class $z$ of
$W^*=W\otimes\varepsilon(\det W)^{-1}$ vanishes on commutators (\cref{lem:window:heights}); a commutator trick then
kills the strict Selmer group of $W^*$. Local duality for characters upgrades this to the Greenberg Selmer
group of $W^*$, the Greenberg--Wiles formula transfers the vanishing to $W$, and Bloch--Kato classes of $W$ are
Greenberg classes because the quotient character of $W$ at $\ell$ has negative Hodge--Tate weight. This part is
proved here from the statements of Galois cohomology contained in $\sfname{HypothesesH'}$
(\cref{pf:kato:hyp:litcoh}); only some auxiliary statements (a bound on the torsion of integral classes, the
counting argument based on the Greenberg--Wiles formula, and a lemma on homomorphisms from compact subgroups of
$\GL_2$ to $\Z_\ell$) are cited from the formal development.

\emph{Construction of the Euler system} (\cref{pf:kato:sec:zeta,pf:kato:sec:katozeta}). The Euler system is built from Kato's zeta
elements: Kummer classes of Siegel units on a tower of modular curves, realised through the Kummer theory of the
field of modular functions of all levels instead of \'etale cohomology, projected to the part attached to the
source newform; an explicit reciprocity law computes their Coates--Wiles derivatives along
$\Q_\ell(\mu_{\ell^\infty})$; and the non-vanishing of the bottom class at $\ell$ is deduced, through Coleman
power series, from the non-vanishing of an $\ell$-adic $L$-value. We prove the existence statement from its main
ingredients. Those which are forms of published theorems (of Kato, Shimura, Coleman, Jacquet--Shalika) are quoted from the
literature, with a precise description of the formal statement and of the differences; their formal proofs (about $120{,}000$ of
the $143{,}000$ lines of the formal proof) are not reproduced.

The main theorem is restated and proved in \cref{pf:kato:sec:proof} (\cref{pf:kato:main}); the two parts of
\cref{lem:window:heights} are \cref{pf:kato:heights,pf:kato:strict}.

% ---------------------------------------------------------------------------
\subsection{Conventions}\label{pf:kato:sec:conv}

Throughout, $\ell$ is a prime and $\katoQlb$ is an algebraic closure of $\Q_\ell$ (not completed), with the
unique extension $|\cdot|$ of the $\ell$-adic absolute value. A \emph{coefficient ring} is a
$\Z_\ell$-subalgebra $\katoO\subset\katoQlb$ which is a discrete valuation ring and finite over $\Z_\ell$, i.e.\
the ring of integers of a finite extension of $\Q_\ell$ in $\katoQlb$; its elements satisfy $|x|\le1$. For
$a,b\in\katoQlb$ and $n\in\N$ we write
\[
  a\katoeq{n}b\quad\Longleftrightarrow\quad|a-b|\le|\ell^n|,
\]
and $a\katoeq n b$ for vectors if it holds coordinatewise. By the ultrametric inequality $\katoeq n$ is an
equivalence relation compatible with sums, and $a\katoeq nb$ implies $ca\katoeq ncb$ when $|c|\le1$. We use
these two facts without comment. For $x\in\Z/\ell^n$ we let $\katorep x\in\{0,\dots,\ell^n-1\}$ be its
representative.

$\Omega$ is an algebraic closure of $\Q$ and $G_\Q=\Aut(\Omega/\Q)$, with the Krull topology; for a number
field $F\subset\Omega$, $G_F\subset G_\Q$ is the subgroup fixing $F$. A \emph{place over} a finite prime $v$ of
$F$ is a valuation subring $A\subset\Omega$ such that $x\in\OO_F$ has valuation $<1$ at $A$ if and only if
$x\in v$. Its decomposition group $D_A$ is the stabiliser of $A$ in $G_F$, and its inertia group $I_A$ consists
of the $\sigma\in D_A$ with $\sigma x-x$ in the maximal ideal of $A$ for all $x\in A$. An element $\sigma$ is a
\emph{Frobenius at $A$ for $q$} if $\sigma\in D_A$ and $\sigma x-x^q$ lies in the maximal ideal of $A$ for all $x\in A$;
we write $\katoN v=\#\OO_F/v$, and for $F=\Q$ we identify $v$ with the prime $p$, so $\katoN v=p$. All places
over a given prime are conjugate under $G_F$, and $D_{gA}=gD_Ag^{-1}$, $I_{gA}=gI_Ag^{-1}$. For a finite set $S$
of primes, $N_S\subset G_\Q$ is the closure of the normal subgroup generated by the inertia groups of all
places over primes not in $S$, and $G_S=G_\Q/N_S$.

$\varepsilon\colon G_\Q\to\Z_\ell^\times\subset\katoQlb^\times$ is the $\ell$-adic cyclotomic character, and
$\varepsilon_r\colon G_\Q\to(\Z/\ell^r)^\times$ its reduction. A homomorphism $\chi$ on $G_\Q$ \emph{kills
$\mu_M$} if $\chi(s)=1$ for every $s$ fixing the $M$-th roots of unity, i.e.\ if $\chi$ factors through
$\Gal(\Q(\mu_M)/\Q)$. For $r\colon G_\Q\to\GL_n(\katoQlb)$ and a character $\chi$, $r\otimes\chi$ is
$s\mapsto\chi(s)r(s)$; for $n=2$ we write $\det W$ for the determinant character of $W$. The representation $W$
acts on column vectors in $\katoQlb^2$, with coordinates indexed by $0,1$, and $W(s)_{ij}$ denotes its matrix
entries. A representation $r$ is \emph{irreducible} on a subgroup $H$ if $\katoQlb^n$ has no $r(H)$-stable
subspace other than $0$ and $\katoQlb^n$; it has the \emph{flag} $(\kappa_0,\kappa_1)$ on $H$ (characters of $H$)
if some conjugate $gr(s)g^{-1}$, $s\in H$, is upper triangular with diagonal $(\kappa_0(s),\kappa_1(s))$; two
representations are \emph{isomorphic} if they are conjugate by one matrix. A representation $r$ of $G_\Q$ has
\emph{level $M$ away from $\ell$} if for every prime $w\ne\ell$ there are a place $B$ over $w$ and characters
$\chi_0,\chi_1$ killing $\mu_M$ such that some conjugate of $r$ is $\mathrm{diag}(\chi_0,\chi_1)$ on $I_B$.

For a topological group $H$ acting through $\rho$ on a topological module $V$, $\katoZ(H,\rho)$ is the module of
\emph{continuous} maps $c\colon H\to V$ with $c(st)=c(s)+\rho(s)c(t)$, and $\katoB(H,\rho)$ the module of the
coboundaries $s\mapsto\rho(s)x-x$, $x\in V$. Cocycles of $G_\Q$ are continuous for the Krull topology and the
$\ell$-adic topology of $\katoQlb^2$.

\begin{definition}[Bloch--Kato classes]\label{pf:kato:def:bk}
Let $\rho$ be a representation of $G_F$ on a $\katoQlb$-vector space $V$. A cocycle $c\in\katoZ(G_F,\rho)$ is a
\emph{Bloch--Kato class} if
\begin{enumerate}[label=(\roman*)]
\item for every finite prime $w\nmid\ell$ of $F$ there are a place $A$ over $w$ and $x\in V$ with
  $c(s)=\rho(s)x-x$ for $s\in I_A$; and
\item for every embedding $\tau\colon F\to\katoQlb$ there is an embedding $j\colon\Omega\to\katoQlb$ extending
  $\tau$ such that, on the group $D_j=\{(\sigma,\sigma')\in G_F\times\Aut(\katoQlb/\Q_\ell):\sigma'\circ j=j\circ\sigma\}$,
  the cocycle $(\sigma,\sigma')\mapsto c(\sigma)$ is \emph{crystalline}: there is $x\in V\otimes_{\Q_\ell}B_{\mathrm{cris}}$
  with $c(\sigma)\otimes1=(\rho(\sigma)\otimes\sigma')x-x$ for all $(\sigma,\sigma')\in D_j$.
\end{enumerate}
We write $\katoHf(F,\rho)=0$ (or $\katoHf(F,W)=0$ when $\rho$ is the action of a matrix representation $W$ on
column vectors) if every Bloch--Kato class is a coboundary.
\end{definition}

\begin{remark}\label{pf:kato:rem:bcris}
Here $B_{\mathrm{cris}}$ is the ring of a period ring package fixed once and for all in the development; it is
characterized by a list of axioms, not constructed as Fontaine's ring. As noted at the end of
\cref{sec:window:kato}, the crystalline condition is used only through Bloch--Kato's local dimension formula,
which is part of the hypothesis package for every structure satisfying these axioms
(\cref{pf:kato:hyp:litcoh}(BK)); with Fontaine's ring \cite{Fontaine1994}, $\katoHf(F,\rho)=0$ is the
vanishing of the Bloch--Kato Selmer group \cite{BlochKato1990} with $\katoQlb$-coefficients and no condition
at the archimedean places. Since $\katoQlb$ is not complete, the statement concerns rational, not integral,
Selmer groups.
\end{remark}

\begin{definition}[Newforms and attached representations]\label{pf:kato:def:newform}
A \emph{newform of weight $k$} is a triple $(M,f,\chi)$ where $M\ge1$, $f$ is a cusp form of weight $k$ for
$\Gamma_1(M)$ and $\chi$ is a Dirichlet character modulo $M$, such that $f|_k\gamma=\chi(d)f$ for every
$\gamma=\left(\begin{smallmatrix}a&b\\c&d\end{smallmatrix}\right)\in\Gamma_0(M)$; the $q$-expansion coefficients
$a_n(f)$ satisfy $a_1(f)=1$ and, for every prime $p\nmid M$ and $n\ge1$,
$a_{pn}(f)+\chi(p)p^{k-1}a_{n/p}(f)=a_p(f)a_n(f)$ (with $a_{n/p}=0$ if $p\nmid n$); and for every proper
divisor $M'$ of $M$, every cusp form $g$ of weight $k$ for $\Gamma_1(M')$ with a character $\chi'$ modulo $M'$
which satisfies the same Hecke relations with the eigenvalues $a_p(f)$ for all but finitely many $p$ is zero.
Its level is $M$, and $H_p(X)=X^2-a_p(f)X+\chi(p)p^{k-1}$. Let $\iota\colon\katoQlb\to\C$ be a field
isomorphism and $V\colon G_\Q\to\GL_2(\katoQlb)$ a homomorphism. The newform is \emph{attached to $V$ via
$\iota$} if $V$ is semisimple and, for all but finitely many primes $p$ and every place $A$ over $p$, $V$ is
trivial on $I_A$ and, for every Frobenius $\sigma$ at $A$ for $p$, the characteristic polynomial of
$V(\sigma^{-1})$, transported by $\iota$, is $H_p(X)$.
\end{definition}

% ---------------------------------------------------------------------------
\subsection{The statement}\label{pf:kato:sec:statement}

\begin{definition}[The data $\sfname{Used}$]\label{pf:kato:def:used}
Let $\ell$ be a prime, $K\subset\Omega$ a number field, $e,N\in\N$, $\theta,\psi\colon G_\Q\to\katoQlb^\times$
homomorphisms and $W\colon G_\Q\to\GL_2(\katoQlb)$ a homomorphism. $\sfname{Used}(\ell,K,e,N,\theta,\psi,W)$ is
the conjunction of:
\begin{enumerate}[label=(U\arabic*),leftmargin=*]
\item\label{pf:kato:U1} $\ell\ne2$, $e\ge1$, and $\ell-1\nmid2e$;
\item\label{pf:kato:U2} $[K:\Q]=2$, $K$ is totally complex, and $\ell\nmid|d_K|\,h_K$;
\item\label{pf:kato:U3} $N\ge1$, $N$ is even, and $\gcd(|d_K|,N)=1$;
\item\label{pf:kato:U4} $\ker\theta=G_K$ and $\theta^2=1$;
\item\label{pf:kato:U5} $\psi$ kills $\mu_N$; $\psi^m=1$ for some $m$ prime to $\ell$; and $\psi^2=1$ on $I_A$ for
  every place $A$ over $\ell$;
\item\label{pf:kato:U6} there is a complex conjugation $c$ (i.e.\ $\phi\circ c=\overline{\phi}$ for some embedding
  $\phi\colon\Omega\to\C$) with $\varepsilon(c)^e\psi(c)=1$;
\item\label{pf:kato:U7} $W$ is continuous, has level $M=N|d_K|$ away from $\ell$, and has level $M$ and weights
  $(e,-e)$ at $\ell$: at the prime $\ell$ there are a place $A$ over $\ell$, characters $\chi_0,\chi_1$ killing
  $\mu_M$ and a flag $(\kappa_0,\kappa_1)$ of $W$ on $D_A$ with $\kappa_0=\varepsilon^e\chi_0$,
  $\kappa_1=\varepsilon^{-e}\chi_1$ on $I_A$, such that either $\kappa_i\ne\kappa_j\varepsilon$ for all $i,j$ or $W$ is
  potentially crystalline at $A$;
\item\label{pf:kato:U8} for some place $A$ over $\ell$ and some $N'\ge1$, every $s\in I_A$ fixing $\mu_{N'}$ has
  characteristic polynomial $(X-\varepsilon(s)^e)(X-\varepsilon(s)^{-e})$;
\item\label{pf:kato:U9} $\det W=\theta$;
\item\label{pf:kato:U10} for every place $A$ over $\ell$ and $s\in I_A$,
  $\tr W(s)=\varepsilon(s)^e\psi(s)+(\varepsilon(s)^e\psi(s))^{-1}$;
\item\label{pf:kato:U11} for all $g\in G_\Q$,
  $\bigl|\tr W(g)-\theta(g)\varepsilon(g)^e\psi(g)-(\varepsilon(g)^e\psi(g))^{-1}\bigr|<1$;
\item\label{pf:kato:U12} there are a place $A$ over $\ell$, a character $\chi$ killing $\mu_\ell$ and a flag
  $(\kappa_0,\kappa_1)$ of $W$ on $D_A$ such that $I_A\subset\ker\theta$ and $\kappa_0=\varepsilon^e\chi$,
  $\kappa_1=\varepsilon^{-e}\chi$ on $I_A$;
\item\label{pf:kato:U13} $W$ is irreducible on $G_K\cap\ker\varepsilon$.
\end{enumerate}
\end{definition}

\begin{definition}[Window witnesses]\label{pf:kato:def:wit}
Let $F\subset\Omega$ be a CM field, $\iota\colon\katoQlb\to\C$ an isomorphism and $r\colon G_F\to\GL_n(\katoQlb)$.
At a prime $v\mid\ell$ of $F$, $r$ is \emph{unlinked} if at every place $A$ over $v$ it has a flag $\kappa$ on
$D_A$ with $\kappa_i\ne\kappa_j\varepsilon$ for all $i,j$, and \emph{potentially unlinked} if at every place $A$
over $v$ it has a flag $\kappa$ on $D_A$ with $\kappa_i^m\ne(\kappa_j\varepsilon)^m$ for all $i,j$ and $m\ge1$.
$\mathrm{Win}(\iota,r)$ means: at every $v\mid\ell$, unlinked implies potentially unlinked; and for every
regular algebraic cuspidal automorphic representation $\pi$ of $\GL_n(\A_F)$ attached to $r$ via $\iota$ there is
$t\in\R$ such that at every finite prime $v$ which does not divide $\ell$, or at which $r$ is potentially
unlinked, $\pi_v$ has a supercuspidal support made of characters $\xi$ with
$\katoN v^{t-1/2}<|\xi(\varpi_v)|<\katoN v^{t+1/2}$. For $r\colon G_\Q\to\GL_n(\katoQlb)$,
$\sfname{WitAQ}(\iota,r)$ means that $\mathrm{Win}(\iota,r|_{G_{K'}})$ holds for every CM field $K'\subset\Omega$
of degree $2$ such that $r|_{G_{K'}}$ is irreducible.
\end{definition}

\begin{definition}[$\sfname{KatoTwoW}$, $\sfname{KatoTwoW}^{\mathrm{sph}}$]\label{pf:kato:def:katotwow}
$\sfname{KatoTwoW}$ is the proposition: for all $\ell,\Omega,K,e,N,\theta,\psi,W$ with
$\sfname{Used}(\ell,K,e,N,\theta,\psi,W)$, every isomorphism $\iota\colon\katoQlb\to\C$ and every newform $f$ of
weight $2e+1$ attached to $W\otimes\varepsilon^{-e}$ via $\iota$, if $\sfname{WitAQ}(\iota,W)$ holds then
\[
  \katoHf(\Q,W\otimes1)=0\qquad\text{and}\qquad\katoHf(\Q,W\otimes\theta)=0.
\]
$\sfname{KatoTwoW}^{\mathrm{sph}}$ is the same proposition with the hypothesis $\sfname{WitAQ}(\iota,W)$ replaced by the
spherical condition $\mathrm{Sph}(\iota,e,W)$ of \cref{pf:kato:def:sph}.
\end{definition}

\begin{remark}[Hypotheses which play no role]\label{pf:kato:rem:unused}
The proof below uses only the following parts of $\sfname{Used}$: $\ell\ne2$ and $e\ge1$; the continuity of $W$
and its level away from $\ell$ in \ref{pf:kato:U7}; $N\ge1$; \ref{pf:kato:U4}; the total complexity of $K$;
\ref{pf:kato:U9}, \ref{pf:kato:U12} and \ref{pf:kato:U13} (see \cref{pf:kato:usedadm}). The conditions
$\ell-1\nmid2e$, $[K:\Q]=2$, $\ell\nmid d_Kh_K$, $N$ even, $\gcd(d_K,N)=1$, all conditions on $\psi$
(\ref{pf:kato:U5}, \ref{pf:kato:U6}, \ref{pf:kato:U10}, \ref{pf:kato:U11}), the condition at $\ell$ in
\ref{pf:kato:U7} and \ref{pf:kato:U8} are not used. In particular the residual congruence \ref{pf:kato:U11}, which is
how $\sfname{KatoTwoW}$ is used for Eisenstein congruences (\cref{thm:window:countI}), plays no role in its proof.
The newform $f$ of the statement is used only to show that $W$ is modular (\cref{pf:kato:def:admitted}(f)), and
the hypothesis $\sfname{WitAQ}$ only to show that the automorphic representations attached to twists of $W$ are
unramified at $\ell$ (\cref{pf:kato:sphofwit}).
\end{remark}

The hypothesis package $\sfname{HypothesesH'}$ is described in \cref{sec:window:arch}. We use two parts of it.

\begin{hypothesis}[Automorphic statements]\label{pf:kato:hyp:aut}
\textup{(H)} is the sub-package of $\sfname{HypothesesH'}$ consisting of statements of published theorems on
automorphic forms and Galois representations (\cref{sec:literature:aut,sec:literature:cft}). Through the results
cited below, the proof uses from it: the existence of regular algebraic cuspidal automorphic representations of
$\GL_2(\A_\Q)$ attached to newforms and to the Galois representations attached to newforms, with Satake parameters
given by the Hecke polynomials; Galois representations attached to regular algebraic polarized cuspidal
representations over CM fields, i.e.\ $\sfname{Row0b}$ (\cref{def:row0b}); Artin reciprocity and the
correspondence between characters of finite order and algebraic Hecke characters; strong multiplicity one; and
the statements of Galois cohomology listed in \cref{pf:kato:hyp:litcoh}. The semisimplicity of spaces of cusp
forms, which is used once (in \cref{pf:kato:srcs}), is a further field of $\sfname{HypothesesH'}$.
\formal{UpperTheory.V.Af.Hypotheses'}
\end{hypothesis}

\begin{hypothesis}[Galois cohomology]\label{pf:kato:hyp:litcoh}
Let $\ell$ be odd. The following statements of published theorems are part of \textup{(H)}.
\begin{enumerate}[label=(\textup{\arabic*}),leftmargin=*]
\item[(PT)] (Poitou--Tate, \cite{Tate1963,Milne2006,NeukirchSchmidtWingberg2008}.) Let $S$ be a finite set of
  primes containing $\ell$, $n\ge1$ and $M$ a finite $\Z/\ell^n$-module with a continuous action of $G_S$. Then
  $\#\Sha^2(G_S,M)=\#\Sha^1(G_S,M^\vee(1))$, where $\Sha^i$ is the kernel of the restriction to the decomposition
  groups at all $v\in S$ and $M^\vee(1)=\Hom(M,\mu_{\ell^n})$; and for $M=\mu_{\ell^n}$ the local invariant maps
  $\mathrm{inv}_v\colon H^2(D_v,\mu_{\ell^n})\to\Z/\ell^n$ are isomorphisms whose sum vanishes on the image of
  $H^2(G_S,\mu_{\ell^n})$.
\item[(LF, LT, LE)] For finite extensions $K$ of $\Q_p$ and finite $\Z/\ell^n$-modules $M$ of $G_K$: finiteness of
  $H^i(G_K,M)$, local Tate duality and the local Euler characteristic formula \cite{Tate1963,Milne2006}.
\item[(GF, GE, UO)] Finiteness of $H^i(G_S,M)$, Tate's global Euler characteristic formula, and the orthogonality
  of unramified classes under local duality.
\item[(GW)] The Greenberg--Wiles formula \cite[Prop.~1.6]{Wiles1995}, \cite{DarmonDiamondTaylor1997}
  for finite $\Z/\ell^n$-modules with local conditions.
\item[(BK)] For every period ring package satisfying the axioms of \cref{pf:kato:rem:bcris} and every finite
  extension $K$ of $\Q_\ell$: Bloch--Kato's local dimension formula for $\katoHf$ and its duality
  \cite[\S3]{BlochKato1990}. In the form used here: if $\kappa$ is a continuous character of $G_{\Q_\ell}$ with
  integral values in a finite extension of $\Q_\ell$ which agrees with $\varepsilon^{w}$ on an open subgroup of
  inertia, for some $w<0$, then every crystalline cocycle of $\kappa$ is a coboundary.
\end{enumerate}
\formal{GlobalTheory.Mem.LiteralAutomorphy.LitCoh}
\end{hypothesis}

\begin{theorem}[\cref{thm:kato}]\label{pf:kato:main}
$\sfname{HypothesesH'}$ implies $\sfname{KatoTwoW}$. That is: assume $\sfname{HypothesesH'}$; let
$\sfname{Used}(\ell,K,e,N,\theta,\psi,W)$ hold, let $\iota\colon\katoQlb\to\C$ be an isomorphism and $f$ a newform of
weight $2e+1$ attached to $W\otimes\varepsilon^{-e}$ via $\iota$, and assume $\sfname{WitAQ}(\iota,W)$. Then every
Bloch--Kato class of $W$, and every Bloch--Kato class of $W\otimes\theta$, is a coboundary.
\formallabel{UpperTheory.V.Z.Kato.WeilSphericalRecords.katoTwoW_of_printsH}
\end{theorem}

The proof is given in \cref{pf:kato:sec:proof}.

% ---------------------------------------------------------------------------
\subsection{Reduction to admitted representations}\label{pf:kato:sec:reduction}

For a homomorphism $W\colon G_\Q\to\GL_2(\katoQlb)$ put
\[
  H_0(W)=\ker\varepsilon\cap\ker\det W,\qquad N_0(W)=\ker\varepsilon\cap\ker W,
\]
two closed normal subgroups of $G_\Q$ with $N_0(W)\subset H_0(W)$.

\begin{definition}[Admitted representations]\label{pf:kato:def:admitted}
Let $\iota\colon\katoQlb\to\C$ be an isomorphism and $e\in\N$. A homomorphism $W\colon G_\Q\to\GL_2(\katoQlb)$ is
\emph{admitted} (for $\iota,e$) if:
(a) $W$ is continuous; (b) $(\det W)^2=1$; (c) $\det W(c)=-1$ for every complex conjugation $c$; (d) $W$ is
irreducible on $H_0(W)$; (e) $W$ has level $M$ away from $\ell$ for some $M\ge1$; (f) there are a newform $f$ of
weight $2e+1$, an integer $M\ge1$ and a character $\varphi$ killing $\mu_M$ such that $f$ is attached to
$W\otimes\varepsilon^{-e}\varphi$ via $\iota$; (g) there are a place $A$ over $\ell$, a character $\chi$ killing
$\mu_\ell$ and a flag $(\kappa_0,\kappa_1)$ of $W$ on $D_A$ such that $I_A\subset\ker\det W$ and
$\kappa_0=\varepsilon^e\chi$, $\kappa_1=\varepsilon^{-e}\chi$ on $I_A$.
\formal{GlobalTheory.Ord2.admitted}
\end{definition}

\begin{lemma}\label{pf:kato:usedadm}
If $\sfname{Used}(\ell,K,e,N,\theta,\psi,W)$ holds and a newform $f$ of weight $2e+1$ is attached to
$W\otimes\varepsilon^{-e}$ via $\iota$, then $\det W=\theta$ and $W$ is admitted.
\formallabel{UpperTheory.V.Z.Kato.Used.adm}
\end{lemma}

\begin{proof}
$\det W=\theta$ is \ref{pf:kato:U9}, and $\theta^2=1$ by \ref{pf:kato:U4}; this gives (b). Let $c$ be a complex
conjugation, with $\phi\circ c=\overline\phi$ for an embedding $\phi\colon\Omega\to\C$. As $\theta(c)^2=1$,
$\theta(c)=\pm1$. If $\theta(c)=1$, then $c\in\ker\theta=G_K$ fixes $K$, so the complex embedding
$\phi|_K$ satisfies $\overline{\phi|_K}=\phi|_K$ and is real, contradicting the total complexity of $K$; hence
$\theta(c)=-1$, which is (c). The group $G_K\cap\ker\varepsilon$ is contained in
$H_0(W)=\ker\varepsilon\cap\ker\theta$, so a subspace stable under $W(H_0(W))$ is stable under
$W(G_K\cap\ker\varepsilon)$; by \ref{pf:kato:U13}, $W$ is irreducible on $H_0(W)$, which is (d). Continuity and
the level $N|d_K|\ge1$ away from $\ell$ are part of \ref{pf:kato:U7} (and $N\ge1$, $d_K\ne0$). For (f) take
$M=1$ and $\varphi=1$. Finally (g) is \ref{pf:kato:U12}, rewritten with $\theta=\det W$.
\end{proof}

\begin{lemma}\label{pf:kato:keropen}
If $W$ is continuous and $(\det W)^2=1$, then $\ker\det W$ is open in $G_\Q$, and $\det W$ kills $\mu_M$ for
some $M\ge1$.
\formallabel{GlobalTheory.Ord2.TwDet.ker_open}
\end{lemma}

\begin{proof}
The values of $\det W$ are $\pm1$, and $1\ne-1$ in $\katoQlb$; so $\ker\det W$ is the complement of the closed
set $(\det W)^{-1}(-1)$. The kernel of a homomorphism to an abelian group contains the commutator subgroup. By
the Kronecker--Weber theorem in the form of \lref{UpperTheory.V.CF.kwAbs} (every open subgroup of $G_\Q$
containing the commutator subgroup contains the subgroup fixing $\mu_M$ for some $M\ge1$; see
\cite[Ch.~14]{Washington1997}), $\det W$ kills some $\mu_M$.
\end{proof}

\begin{lemma}\label{pf:kato:twdet}
If $W$ is admitted, so is $W\otimes\det W$.
\formallabel{GlobalTheory.Ord2.admissible_twist_by_det}
\end{lemma}

\begin{proof}
Put $\delta=\det W$; by \cref{pf:kato:keropen} it is continuous and kills some $\mu_N$, $N\ge1$. Then
$W\otimes\delta$ is continuous and $\det(W\otimes\delta)=\delta^2\det W=\det W$, which gives (b), (c) and
$H_0(W\otimes\delta)=H_0(W)$; on this group $W\otimes\delta=W$, which gives (d). If $W$ is conjugate to
$\mathrm{diag}(\chi_0,\chi_1)$ on $I_B$ with $\chi_i$ killing $\mu_M$, then $\delta=\chi_0\chi_1$ on $I_B$ and $W\otimes\delta$
is conjugate to $\mathrm{diag}(\chi_0\chi_0\chi_1,\chi_1\chi_0\chi_1)$ there, which gives (e). If $f$ is attached to
$W\otimes\varepsilon^{-e}\varphi$ with $\varphi$ killing $\mu_M$, then
$W\otimes\varepsilon^{-e}\varphi=(W\otimes\delta)\otimes\varepsilon^{-e}(\varphi\delta)$ because $\delta^2=1$, and $\varphi\delta$ kills
$\mu_{MN}$; this is (f). For (g), the flag $(\kappa_0\delta,\kappa_1\delta)$ of $W\otimes\delta$ on $D_A$ has the same
restriction to $I_A$ as $(\kappa_0,\kappa_1)$, because $I_A\subset\ker\delta$.
\end{proof}

\begin{definition}[Spherical condition, sources]\label{pf:kato:def:sph}
Let $\iota$, $e$ and $W$ be as above.
\begin{enumerate}[label=(\roman*)]
\item $W$ is \emph{spherical at $\ell$}, written $\mathrm{Sph}(\iota,e,W)$, if: for every $M\ge1$, every character
  $\varphi$ killing $\mu_M$, every place $A$ over $\ell$, every flag $\kappa$ of $W\otimes\varphi$ on $D_A$ with
  $\kappa_0=\varepsilon^e$ and $\kappa_1=\varepsilon^{-e}$ on $I_A$, and every regular algebraic cuspidal automorphic representation
  $\pi$ of $\GL_2(\A_\Q)$ attached to $W\otimes\varphi$ via $\iota$ (at almost all primes, $\pi$ is unramified with Satake
  parameters matching the eigenvalues of Frobenius), $\pi$ has Satake parameters at $\ell$, i.e.\ $\pi_\ell$ is
  unramified.
\item $W$ \emph{has an ordinary source}, written $\mathrm{Src}(\iota,e,W)$, if: (1) there are a newform $g$ of weight
  $2e+1$, a character $\chi$ killing $\mu_\ell$, a place $A$ over $\ell$ and a flag $\kappa$ of $W$ on $D_A$ such that
  $\ell$ does not divide the level of $g$, $|\iota^{-1}(a_\ell(g))|=1$, $g$ is attached to
  $W\otimes\varepsilon^{-e}(\det W)^{-1}\chi$ via $\iota$, and $\kappa_1\varepsilon^e(\det W)^{-1}\chi=1$ on $I_A$; and (2) for
  every place $A$ over $\ell$ and every flag $\kappa$ of $W$ on $D_A$ there is $s\in D_A$ with
  $\kappa_0(s)\ne\varepsilon(s)$.
\item \textup{(S)} is the proposition: for $\ell\ne2$ and $e\ge1$, every admitted $W$ with
  $\mathrm{Sph}(\iota,e,W)$ satisfies $\mathrm{Src}(\iota,e,W)$.
\item \textup{(V)} is the proposition: for every prime $\ell\ne2$, every algebraic closure $\Omega$ of
  $\Q$, every $\iota$, every $e\ge1$ and every $W$ which is admitted and satisfies $\mathrm{Src}(\iota,e,W)$,
  $\katoHf(\Q,W)=0$.
\end{enumerate}
\formallabel{UpperTheory.V.Z.Kato.SphAt}
\formallabel{GlobalTheory.Ord2.SrcL}
\formallabel{GlobalTheory.Ord2.VanishOrdL}
\end{definition}

In $\mathrm{Src}$, the newform $g$ is $\ell$-ordinary of level prime to $\ell$, and the last condition of (1)
says that the quotient character of the flag is $\varepsilon^{-e}\det W\chi^{-1}$ on $I_A$. Condition (2) excludes
an ``exceptional'' sub-character equal to $\varepsilon$.

\begin{lemma}\label{pf:kato:sphtw}
If $\mathrm{Sph}(\iota,e,W)$ holds and $\theta$ kills $\mu_N$ with $N\ge1$, then $\mathrm{Sph}(\iota,e,W\otimes\theta)$
holds. \formallabel{sphAt_tw}
\end{lemma}

\begin{proof}
For $\varphi$ killing $\mu_M$, $(W\otimes\theta)\otimes\varphi=W\otimes(\varphi\theta)$, and $\varphi\theta$ kills $\mu_{MN}$; apply
$\mathrm{Sph}(\iota,e,W)$ to $\varphi\theta$.
\end{proof}

\begin{theorem}\label{pf:kato:twos}
Assume the statements of \cref{pf:kato:hyp:litcoh} and \textup{(S)}. Then $\sfname{KatoTwoW}^{\mathrm{sph}}$ holds.
\formallabel{UpperTheory.V.Z.Kato.WeilSphericalRecords.twoS}
\end{theorem}

\begin{proof}
By \cref{pf:kato:vanishordl} and \cref{pf:kato:zetaur}, \textup{(V)} holds. Let the data of the spherical variant be
given: the conditions $\sfname{Used}(\ell,K,e,N,\theta,\psi,W)$ hold, $f$ is a newform attached to $W\otimes\varepsilon^{-e}$, and
$\mathrm{Sph}(\iota,e,W)$ holds. By \cref{pf:kato:usedadm}, $\det W=\theta$ and $W$ is
admitted; also $\ell\ne2$ and $e\ge1$. Since $W\otimes1=W$, the first vanishing is \textup{(V)} applied to $W$, with
$\mathrm{Src}(\iota,e,W)$ supplied by \textup{(S)}. For the second, $W\otimes\theta=W\otimes\det W$ is admitted by
\cref{pf:kato:twdet}; $\det W$ kills some $\mu_M$ (\cref{pf:kato:keropen}), so the representation $W\otimes\det W$ is
spherical at $\ell$ by \cref{pf:kato:sphtw}; then \textup{(S)} shows that it has an ordinary source, and \textup{(V)} gives
$\katoHf(\Q,W\otimes\theta)=0$.
\end{proof}

% ---------------------------------------------------------------------------
\subsection{Local--global compatibility at \texorpdfstring{$\ell$}{l} and sources}\label{pf:kato:sec:sph}

This subsection supplies \textup{(S)} and the spherical condition. Its arguments assemble statements about
automorphic representations proved elsewhere in the development; we state precisely how they are combined.
For a place $A$ over $\ell$ and a flag $\kappa$ of $r\colon G_F\to\GL_2(\katoQlb)$ on $D_A$ (for $F=\Q$ or a CM
field $F$), $\kappa$ is \emph{exact} (for $e$) if $\kappa_0=\varepsilon^{e}$ and $\kappa_1=\varepsilon^{-e}$ on $I_A$. For a prime
power $q$ and a multiset $B$ of Satake parameters, $\mathrm{Fm}(\iota,e,q,A,\kappa,B)$ means: for every
Frobenius $\sigma\in D_A$ for $q$ there is $b\in\{q^{1/2}\beta:\beta\in B\}$ with
$q^{e}\,\iota\bigl(\kappa_0(\sigma)\varepsilon(\sigma)^{-e}\bigr)\,b=1$.

\begin{definition}\label{pf:kato:def:locq}
We consider the following statements.
\begin{enumerate}[label=(\roman*)]
\item \textup{(LG)} (local--global compatibility at $\ell$): for $\ell\ne2$, $e\ge1$, every continuous irreducible
  $V\colon G_\Q\to\GL_2(\katoQlb)$ whose determinant kills some $\mu_M$, every regular algebraic cuspidal $\pi$ of
  $\GL_2(\A_\Q)$ attached to $V$ via $\iota$, every place $A$ over $\ell$ and every exact flag $\kappa$ of $V$ on $D_A$: if
  $\sfname{WitAQ}(\iota,V)$ holds or $\pi$ has Satake parameters at $\ell$, then $\pi$ has Satake parameters at $\ell$, and
  $\mathrm{Fm}(\iota,e,\ell,A,\kappa,B)$ holds for all Satake parameters $B$ of $\pi$ at $\ell$.
\item \textup{(TW)} (twisting window witnesses): $\sfname{WitAQ}(\iota,W)$ implies $\sfname{WitAQ}(\iota,W\otimes\varphi)$ for
  every character $\varphi$ killing some $\mu_M$, $M\ge1$.
\item \textup{(SW)} (window witnesses give sphericity): for $\ell\ne2$, $e\ge1$ and continuous irreducible $W$ whose
  determinant kills some $\mu_M$, $\sfname{WitAQ}(\iota,W)$ implies $\mathrm{Sph}(\iota,e,W)$.
\item \textup{(FR)} (Frobenius roots): for $\ell\ne2$, $e\ge1$, continuous irreducible $V$ whose determinant kills some
  $\mu_M$, a place $A$ over $\ell$, an exact flag $\kappa$ of $V$ on $D_A$, and a newform $g$ of weight $2e+1$
  attached to $V\otimes\varepsilon^{-e}$ via $\iota$ whose level is prime to $\ell$: for every Frobenius
  $\sigma\in D_A$ for $\ell$ there is a root $\alpha$ of $H_\ell(X)$ with $\alpha=u$ or $\alpha^{-1}=u$, where
  $u=\iota(\kappa_0(\sigma)\varepsilon(\sigma)^{-e})$.
\item \textup{(TN)} (twisted newforms): for $\ell\ne2$, $e\ge1$, admitted $W$ and every character $\varphi$ killing some
  $\mu_M$, $M\ge1$, some newform of weight $2e+1$ is attached to $(W\otimes\varphi)\otimes\varepsilon^{-e}$ via $\iota$.
\end{enumerate}
We write \textup{(QBC)} for quadratic base change for $\GL_2$: for a totally real field $F$, a totally complex quadratic
extension $E/F$ and a regular algebraic cuspidal $\pi$ of $\GL_2(\A_F)$, either $\pi\cong\pi\otimes\eta_{E/F}$ or $\pi$ has a
cuspidal base change to $E$ \cite{Langlands1980}.
\formallabel{UpperTheory.V.Z.Kato.WeilSphericalRecords.LocQ}
\end{definition}

A representation $P$ of $\GL_2$ over a CM field is \emph{plain} at a prime $w$ if every local component of $P$ at $w$ has
a supercuspidal support consisting of characters, no two of which are linked (i.e.\ differ by $|\cdot|^{\pm1}$).

\begin{theorem}\label{pf:kato:locq}
The statements \textup{(H)}, $\sfname{Row0b}$ for the local parameter maps of \cref{sec:galois} and \textup{(QBC)} imply
\textup{(LG)}.
\formallabel{UpperTheory.V.Z.Kato.WeilSphericalRecords.locQ_of_alt}
\end{theorem}

\begin{proof}
The proof combines the following statements.
\begin{enumerate}[label=(\alph*)]
\item (Split base change; \lref{UpperTheory.V.Z.Kato.WeilSphericalRecords.base_change_split_of_printed_quad_bc},
  from \textup{(QBC)} and \textup{(H)}.) For $V$, $\pi$ as in \textup{(LG)} and a place $A$ over $\ell$, there are an
  imaginary quadratic field $K'\subset\Omega$ with $V|_{G_{K'}}$ irreducible, a regular algebraic polarized cuspidal
  $P$ of $\GL_2(\A_{K'})$ attached to $V|_{G_{K'}}$ via $\iota$, and a prime $w\mid\ell$ of $K'$ with $A$ over $w$,
  $\katoN w=\ell$ and $D_A\subset G_{K'}$, such that for every multiset $B$, $\pi$ has Satake parameters $B$ at $\ell$ if
  and only if $P$ has Satake parameters $B$ at $w$.
\item (Compatibility over CM fields;
  \lref{UpperTheory.V.Z.Kato.WeilSphericalRecords.local_k_of_printed_rec_compat_sph_frob}, from \textup{(H)},
  $\sfname{Row0b}$ and the local parameter maps.) For a CM field $K'$, a continuous $r\colon G_{K'}\to\GL_2(\katoQlb)$, a
  regular algebraic polarized cuspidal $P$ attached to $r$, a prime $w\mid\ell$ of $K'$, a place $A$ over $w$ and an exact
  flag $\kappa$ of $r$ on $D_A$: if $P$ is plain at $w$, then $P$ has Satake parameters at $w$ and
  $\mathrm{Fm}(\iota,e,\katoN w,A,\kappa,B)$ holds for all of them.
\item (\cref{pf:kato:potof} below.) A representation with an exact flag at a place over $w$ is potentially unlinked
  at $w$.
\item (Plainness.) If $\mathrm{Win}(\iota,r)$ holds, $P$ is regular algebraic and attached to $r$, and $r$ is potentially
  unlinked at $w\mid\ell$, then $P$ is plain at $w$: by \lref{UpperTheory.V.Z.Win.guard_of} and
  \lref{UpperTheory.V.Z.Win.plainRec}, the window bound of $\mathrm{Win}(\iota,r)$ at $w$ forces every local component
  of $P$ at $w$ to be plain, for every local parameter map of \lref{UpperTheory.Zero.LLC.recTExists} (which
  provides such maps at every place). If $P$ has Satake parameters at $w$, then $P$ is plain at $w$
  (\lref{UpperTheory.V.Z.Win.plainOf}: unramified and window-bounded local components are plain).
\end{enumerate}
Let $V$, $\pi$, $A$, $\kappa$ be as in \textup{(LG)}, and choose $K'$, $P$, $w$ by (a). As $D_A\subset G_{K'}$, the
restriction of $\kappa$ is an exact flag of $V|_{G_{K'}}$ on the decomposition group of $A$ in $G_{K'}$. If
$\sfname{WitAQ}(\iota,V)$ holds, then $\mathrm{Win}(\iota,V|_{G_{K'}})$ holds since $V|_{G_{K'}}$ is irreducible,
$V|_{G_{K'}}$ is potentially unlinked at $w$ by (c), and $P$ is plain at $w$ by (d). If $\pi$ has Satake
parameters $B$ at $\ell$, then $P$ has Satake parameters $B$ at $w$ by (a), and is plain at $w$ by (d). In both
cases (b) shows that $P$, hence $\pi$, has Satake parameters at $w$, resp.\ $\ell$, and that
$\mathrm{Fm}(\iota,e,\ell,A,\kappa,B)$ holds for the restricted flag and all such $B$ (recall $\katoN w=\ell$).
Since $D_A$ fixes $K'$, a Frobenius at $A$ for $\ell$ in $G_\Q$ is a Frobenius at $A$ in $G_{K'}$, so
$\mathrm{Fm}$ holds for $\kappa$ itself.
\end{proof}

\begin{lemma}\label{pf:kato:potof}
Let $F$ be a number field, $e\ge1$, $r\colon G_F\to\GL_2(\katoQlb)$, $w$ a prime of $F$ above $\ell$, $A$ a place
over $w$ and $\kappa$ an exact flag of $r$ on $D_A$. Then $r$ is potentially unlinked at $w$.
\formallabel{UpperTheory.V.Z.Kato.WeilSphericalRecords.potOf}
\end{lemma}

\begin{proof}
Let $A'$ be a place over $w$, and $\sigma\in G_F$ with $A'=\sigma A$. If $g$ triangularizes $r$ on $D_A$, then
$g\,r(\sigma)^{-1}$ triangularizes $r$ on $D_{A'}=\sigma D_A\sigma^{-1}$, with diagonal characters
$\kappa'_i(s)=\kappa_i(\sigma^{-1}s\sigma)$. Write $e_0=e$, $e_1=-e$, so that $\kappa_i=\varepsilon^{e_i}$ on $I_A$. For all
$i,j$ the integer $e_i-e_j-1$ lies in $\{-1,2e-1,-2e-1\}$ and is non-zero. Suppose
$\kappa_i'^m=(\kappa'_j\varepsilon)^m$ for some $m\ge1$. Choose $g_0\in I_A$ with
$\varepsilon(g_0)^{m(e_i-e_j-1)}\ne1$ ($\varepsilon$ has infinite order on $I_A$). Evaluating at $\sigma g_0\sigma^{-1}\in I_{A'}$,
and using $\varepsilon(\sigma g_0\sigma^{-1})=\varepsilon(g_0)$, gives $\varepsilon(g_0)^{me_i}=\varepsilon(g_0)^{m(e_j+1)}$, a contradiction.
\end{proof}

\begin{theorem}\label{pf:kato:sphofwit}
\textup{(TW)} and \textup{(LG)} imply \textup{(SW)}.
\formallabel{UpperTheory.V.Z.Kato.WeilSphericalRecords.sphOfWit}
\end{theorem}

\begin{proof}
Let $W$ be continuous and irreducible, with $\det W$ killing $\mu_N$, and assume $\sfname{WitAQ}(\iota,W)$. Let
$\varphi$ kill $\mu_M$, $A$ be over $\ell$, $\kappa$ an exact flag of $W\otimes\varphi$ on $D_A$ and $\pi$ attached to
$W\otimes\varphi$. The twist $W\otimes\varphi$ is continuous ($\ker\varphi$ is open), irreducible (twisting preserves stable
subspaces), and $\det(W\otimes\varphi)=\varphi^2\det W$ kills $\mu_{NM}$. By \textup{(TW)},
$\sfname{WitAQ}(\iota,W\otimes\varphi)$ holds, and the first alternative of \textup{(LG)} shows that $\pi$ has
Satake parameters at $\ell$.
\end{proof}

\begin{theorem}\label{pf:kato:twwit}
\textup{(H)} implies \textup{(TW)}.
\formallabel{UpperTheory.V.Z.Kato.WeilSphericalRecords.twWit_alt}
\end{theorem}

\begin{proof}
We use the twisting property of window conditions (\lref{UpperTheory.V.Z.Win.aIsoTwOf_alt}, assuming
\textup{(H)}): for a CM field $F$, if $\mathrm{Win}(\iota,r)$ holds then $\mathrm{Win}(\iota,r')$ holds for every $r'$
isomorphic to $r$, and $\mathrm{Win}(\iota,r\otimes\mu)$ holds for every algebraic character $\mu$ of $G_F$ (a character
corresponding to an algebraic Hecke character). Let $K'$ be imaginary quadratic with $(W\otimes\varphi)|_{G_{K'}}$
irreducible. Then $W|_{G_{K'}}$ is irreducible, so $\mathrm{Win}(\iota,W|_{G_{K'}})$ holds. The restriction of $\varphi$ to
$G_{K'}$ has open kernel, hence is algebraic (\lref{UpperTheory.V.Sd.alg_cyclotomic_of_printed_results}: by the
Artin reciprocity statement of \textup{(H)}, characters of finite order are algebraic), and the twisting property gives
$\mathrm{Win}(\iota,(W\otimes\varphi)|_{G_{K'}})$.
\end{proof}

\begin{theorem}\label{pf:kato:frobroot}
\textup{(H)} and \textup{(LG)} imply \textup{(FR)}.
\formallabel{UpperTheory.V.Z.Kato.WeilSphericalRecords.frobRoot_alt}
\end{theorem}

\begin{proof}
Let $V,A,\kappa,g,\sigma$ be as in \textup{(FR)}. By
\lref{GlobalTheory.T.exists_cuspidal_attached_satake_eq_hecke_poly_of_not_dvd} (from \textup{(H)}),
there is a regular algebraic cuspidal $\pi$ of $\GL_2(\A_\Q)$ attached to $V$ via $\iota$ which has Satake parameters
$B$ at $\ell$ with $\prod_{b}(X-\ell^{e}b)=H_\ell(X)$, the product over $b\in\{\ell^{1/2}\beta:\beta\in B\}$. The second
alternative of \textup{(LG)} applies and gives such a $b$ with $\ell^e\,u\,b=1$. Then $\alpha=\ell^eb$ is a root
of $H_\ell$ and $\alpha^{-1}=u$.
\end{proof}

\begin{theorem}\label{pf:kato:twnew}
$\sfname{HypothesesH'}$ implies \textup{(TN)}.
\formallabel{UpperTheory.V.Z.Kato.WeilSphericalRecords.twNew_alt}
\end{theorem}

\begin{proof}
Let $W$ be admitted and $\varphi$ kill $\mu_M$. By \cref{pf:kato:def:admitted}(f) a newform $f$ of weight $2e+1$ is
attached to $V=W\otimes\varepsilon^{-e}\varphi'$, with $\varphi'$ killing $\mu_{M'}$. The character $\kappa=\varphi\varphi'^{-1}$
kills $\mu_{MM'}$. By \lref{UpperTheory.V.Z.Kato.WeilSphericalRecords.TwistNewform.exists_newform_attached_twist_of_kills_mu}
(from $\sfname{HypothesesH'}$; for a newform of weight $k\ge2$ attached to $V$ and a character $\kappa$ killing some
$\mu_{N}$, $N\ge1$, some newform of weight $k$ is attached to $V\otimes\kappa$; compare \cite{AtkinLi1978}) a newform
$g$ of weight $2e+1$ is attached to $V\otimes\kappa=(W\otimes\varphi)\otimes\varepsilon^{-e}$.
\end{proof}

\begin{lemma}\label{pf:kato:unitroot}
Let $\ell$ be odd, $e\ge1$, $g$ a newform of weight $2e+1$ whose level is prime to $\ell$, and $\alpha\in\C$ a root of $H_\ell(X)$ with
$|\iota^{-1}(\alpha)|=1$. Then $|\iota^{-1}(a_\ell(g))|=1$, and $\alpha^m\ne1$ for every $m\ge1$.
\formallabel{GlobalTheory.T.newform_unit_root_not_root_of_unity_odd_weight}
\end{lemma}

\begin{proof}
Let $c=\chi(\ell)$, where $\chi$ is the character of $g$; it is a root of unity, as $\chi$ has finite order and $\ell$ is prime to the
level. As $H_\ell(\alpha)=0$ and $H_\ell(0)=c\ell^{2e}\ne0$, we have $\alpha\ne0$ and $a_\ell(g)=\alpha+c\ell^{2e}/\alpha$. Applying $\iota^{-1}$, the
second term has norm $|\ell|^{2e}<1=|\iota^{-1}(\alpha)|$, so $|\iota^{-1}(a_\ell(g))|=1$ by the ultrametric inequality. If $\alpha^m=1$, then
$|\alpha|=1$ in $\C$, and $|a_\ell(g)|\ge\ell^{2e}-1$. But $|a_p(g)|\le p^{k/2}+p^{k/2-1}$ for every prime $p$ not dividing the level of a
newform of weight $k$ (\lref{GlobalTheory.T.newform_norm_coeff_prime_le_of_not_dvd_level}; this is weaker than Deligne's bound), and
$\ell^{2e}-1>\ell^{e+1/2}+\ell^{e-1/2}$ for $\ell\ge3$ and $e\ge1$; this is a contradiction.
\end{proof}

\begin{theorem}\label{pf:kato:srcs}
$\sfname{HypothesesH'}$, \textup{(TN)} and \textup{(FR)} imply \textup{(S)}.
\formallabel{UpperTheory.V.Z.Kato.WeilSphericalRecords.srcS_alt}
\end{theorem}

\begin{proof}
Let $\ell\ne2$, $e\ge1$, and $W$ admitted with $\mathrm{Sph}(\iota,e,W)$. Choose $N\ge1$ with $\det W$ killing
$\mu_N$ (\cref{pf:kato:keropen}) and $A,\chi,\kappa$ as in \cref{pf:kato:def:admitted}(g), and put
$\varphi=(\det W)^{-1}\chi^{-1}$, which kills $\mu_{N\ell}$. By \textup{(TN)} there is a newform $g$ of weight
$2e+1$ attached to $V\otimes\varepsilon^{-e}$, where $V=W\otimes\varphi$. The representation $V$ is continuous and
irreducible, has finite level away from $\ell$ (twist the characters of \cref{pf:kato:def:admitted}(e) by $\varphi$),
its determinant $\varphi^2\det W$ kills $\mu_{N\ell}$, and the flag $\kappa\varphi$ of $V$ on $D_A$ is exact: on $I_A$,
$\kappa_0\varphi=\varepsilon^e\chi\chi^{-1}=\varepsilon^e$ and $\kappa_1\varphi=\varepsilon^{-e}$, because $I_A\subset\ker\det W$.

By \lref{GlobalTheory.Ord2.SrcForm.exists_rep_alt} (from \textup{(H)}) there is a regular algebraic
cuspidal $\pi$ attached to $V$ via $\iota$ whose Satake parameters at almost all primes $p$ are the roots of
$H_p(g)$ (suitably normalised). By $\mathrm{Sph}(\iota,e,W)$, applied to $\varphi$, $A$, $\kappa\varphi$ and $\pi$, the
representation $\pi$ has Satake parameters at $\ell$; by
\lref{GlobalTheory.NewLevel.not_dvd_level_of_satake_hecke_poly_cofinite} (from \textup{(H)} and the
semisimplicity of the space of cusp forms contained in $\sfname{HypothesesH'}$), $\ell$ does not divide the level
of $g$. Let $\sigma\in D_A$ be a Frobenius for $\ell$. By \textup{(FR)} there is a root $\alpha$ of $H_\ell(X)$ with
$\alpha^{\pm1}=\iota(u)$, $u=\kappa_0(\sigma)\varphi(\sigma)\varepsilon(\sigma)^{-e}$. Diagonal characters of a flag of a continuous
representation of the compact group $D_A$ take unit values, so $|u|=1$ and $|\iota^{-1}(\alpha)|=1$. By \cref{pf:kato:unitroot},
$g$ is $\ell$-ordinary and $\alpha$ is not a root of unity.

Clause (1) of $\mathrm{Src}$ holds with $\chi^{-1}$: $g$ is attached to
$(W\otimes\varphi)\otimes\varepsilon^{-e}=W\otimes\varepsilon^{-e}(\det W)^{-1}\chi^{-1}$, and on $I_A$,
$\kappa_1\varepsilon^e(\det W)^{-1}\chi^{-1}=\varepsilon^{-e}\chi\varepsilon^e\chi^{-1}=1$. For clause (2), suppose that some flag $\kappa'$ of
$W$ on $D_{A'}$, $A'$ over $\ell$, has $\kappa'_0=\varepsilon$. Transporting $\kappa'$ to $D_A$ by an element of $G_\Q$ carrying $A'$
to $A$, and comparing the diagonal characters of two flags of $W$ on $D_A$, we get $\varepsilon=\kappa_0$ or $\varepsilon=\kappa_1$ on
$D_A$. Let $m\ge1$ with $\chi^m=1$, and $t\in I_A$ such that $\varepsilon(t)$ is not a root of unity. If $\varepsilon=\kappa_1$, then
$\varepsilon^{e+1}=\chi$ on $I_A$ and $\varepsilon(t)^{(e+1)m}=1$, a contradiction. If $\varepsilon=\kappa_0$, then $\varepsilon^{e-1}\chi=1$ on $I_A$;
if $e\ge2$ this gives $\varepsilon(t)^{(e-1)m}=1$, a contradiction, so $e=1$; then $u=\kappa_0(\sigma)\varphi(\sigma)\varepsilon(\sigma)^{-1}=\varphi(\sigma)$
has finite order, so $\alpha$ would be a root of unity, again a contradiction.
\end{proof}

% ---------------------------------------------------------------------------
\subsection{Normal forms, Selmer conditions and Euler systems}\label{pf:kato:sec:nf}

From now on $W\colon G_\Q\to\GL_2(\katoQlb)$ is a homomorphism and $\delta=\det W$. Put
\[
  W^*=W\otimes\varepsilon\delta^{-1},\qquad\text{so } W^*(s)=\varepsilon(s)\delta(s)W(s)\text{ if }\delta^2=1;
\]
since $W$ has dimension two, $W^*\cong W^\vee(1)$. Let $\omega(x,y)=x_0y_1-x_1y_0$ on $\katoQlb^2$. As
$\omega(Mx,My)=\det(M)\,\omega(x,y)$ for every $2\times2$ matrix $M$,
\begin{equation}\label{pf:kato:eq:omega}
  \omega\bigl(W(s)x,\,W^*(s)y\bigr)=\varepsilon(s)\,\omega(x,y)\qquad(s\in G_\Q,\ x,y\in\katoQlb^2).
\end{equation}
For maps $a,b\colon G_\Q\to\katoQlb^2$ and $u\colon G_\Q\to\katoQlb$ define the \emph{cup product} and the
\emph{height}
\begin{equation}\label{pf:kato:eq:cupht}
\begin{gathered}
  (a\cup b)(s,t)=\omega\bigl(a(s),W^*(s)b(t)\bigr),\\
  \katoht_{a,b,u}(s)=\det\bigl(W(s)-1\bigr)u(s)-\omega\bigl(\katoadj(W(s)-1)\,a(s),\,b(s)\bigr),
\end{gathered}
\end{equation}
where $\katoadj$ is the adjugate matrix. If $a$ is a cocycle of $W$ and $b$ of $W^*$, then $a\cup b$ is a
$2$-cocycle with values in $\katoQlb(1)$: $\varepsilon(r)(a\cup b)(s,t)+(a\cup b)(r,st)=(a\cup b)(rs,t)+(a\cup b)(r,s)$, by
\eqref{pf:kato:eq:omega}.

\begin{definition}[Normal form]\label{pf:kato:def:nf}
Let $e\in\N$, $S$ a finite set of primes and $A\subset\Omega$ a valuation subring. $W$ is in \emph{normal form} with
respect to $(e,S,A)$, written $\mathrm{NF}(e,S,A,W)$, if: $\ell\ne2$; $e\ge1$; $W$ is continuous with entries in some coefficient ring; $\ell\in S$; $W$ is trivial on
$I_B$ for every place $B$ over a prime not in $S$; $\delta^2=1$; $\delta(c)=-1$ for every complex conjugation $c$;
$W$ is irreducible on $H_0(W)$; $A$ is a place over $\ell$; $W(s)_{10}=0$ for $s\in D_A$; there is $m\ge1$ with
$W(s)_{00}^m=\varepsilon(s)^{em}$ for $s\in I_A$; and $W(s)_{00}\ne\varepsilon(s)$ for some $s\in D_A$.
\formallabel{GlobalTheory.Ord2.Nf}
\end{definition}

Under $\mathrm{NF}$, $W$ and $W^*$ are upper triangular on $D_A$, and $W$ is unramified outside $S$; so are
$W^*$ and $\varepsilon$ at primes $\ne\ell$, hence $W$, $W^*$ and $\varepsilon$ are trivial on $N_S$.

\begin{definition}[Selmer conditions]\label{pf:kato:def:selmer}
Let $A$ be a place over $\ell$ and $c\colon G_\Q\to\katoQlb^2$.
\begin{itemize}
\item $c$ is \emph{unramified away from $\ell$}, written $\mathrm{Unr}(W,c)$, if for every prime $q\ne\ell$ there are a
  place $B$ over $q$ and $x\in\katoQlb^2$ with $c(s)=W(s)x-x$ for $s\in I_B$.
\item $c$ satisfies the \emph{Greenberg condition} at $A$, written $\mathrm{Gr}(W,A,c)$, if there is $x\in\katoQlb$ with
  $c(s)_1=W(s)_{11}x-x$ for $s\in D_A$.
\item $c$ is \emph{strict} at $A$, written $\mathrm{Str}(W,A,c)$, if there is $x\in\katoQlb^2$ with $c(s)=W(s)x-x$ for
  $s\in D_A$.
\item We write $H^1_{\mathrm{Gr}}(W,A)=0$, resp.\ $H^1_{\mathrm{str}}(W,A)=0$, if every $c\in\katoZ(G_\Q,W)$ with
  $\mathrm{Unr}(W,c)$ and $\mathrm{Gr}(W,A,c)$, resp.\ $\mathrm{Str}(W,A,c)$, is a coboundary (vanishing of the Greenberg
  Selmer group \cite{Greenberg1989}, resp.\ of the strict Selmer group).
\end{itemize}
\formal{GlobalTheory.Ord2.Unr}
\formal{GlobalTheory.Ord2.Gr}
\formal{GlobalTheory.Ord2.Triv}
\formal{GlobalTheory.Ord2.GrZero}
\formal{GlobalTheory.Ord2.StrZero}
\end{definition}

\begin{lemma}\label{pf:kato:placeindep}
Let $c\in\katoZ(G_\Q,W)$ and let $B,B'$ be places over the same prime. If $c(s)=W(s)x-x$ for $s\in I_B$ (resp.\
$s\in D_B$), then $c(s)=W(s)x'-x'$ for $s\in I_{B'}$ (resp.\ $s\in D_{B'}$) for some $x'$. In particular the conditions
$\mathrm{Unr}$ and $\mathrm{Str}$ do not depend on the choice of the places.
\formallabel{GlobalTheory.T.localConditions_independent_of_place_over}
\end{lemma}

\begin{proof}
Write $B'=gB$, so $I_{B'}=gI_Bg^{-1}$. For $s\in I_B$, the cocycle relation and $c(g^{-1})=-W(g)^{-1}c(g)$ give
$c(gsg^{-1})=c(g)+W(g)(W(s)x-x)-W(gsg^{-1})c(g)=W(gsg^{-1})x'-x'$ with $x'=W(g)x-c(g)$. The same computation applies to
decomposition groups.
\end{proof}

\begin{definition}[Congruence cochains]\label{pf:kato:def:congr}
Let $\katoO$ be a coefficient ring, $n\in\N$ and $H\subset G_\Q$ a subgroup.
\begin{itemize}
\item $a\colon G_\Q\to\katoO^2$ is an \emph{$(\katoO,n)$-cocycle of $W$ on $H$} if it is continuous on $H$ and
  $a(st)\katoeq na(s)+W(s)a(t)$ for $s,t\in H$; it is an \emph{$(\katoO,n)$-coboundary on $H$} if there is
  $x\in\katoO^2$ with $a(s)\katoeq nW(s)x-x$ for $s\in H$.
\item A map $\xi\colon G_\Q\times G_\Q\to\katoQlb$ is \emph{$(\katoO,n)$-trivial on $H$} if there is $u\colon
  G_\Q\to\katoO$, continuous on $H$, with $\xi(s,t)\katoeq n\varepsilon(s)u(t)-u(st)+u(s)$ for $s,t\in H$.
\item $\xi$ is an \emph{$(\katoO,n,S)$-$2$-cocycle} if it is continuous with values in $\katoO$,
  $\varepsilon(r)\xi(s,t)+\xi(r,st)\katoeq n\xi(rs,t)+\xi(r,s)$ for all $r,s,t$, and $\xi(sa,tb)\katoeq n\xi(s,t)$ for
  all $s,t$ and $a,b\in N_S$.
\end{itemize}
\formal{GlobalTheory.Ord2.Coc1}
\formal{GlobalTheory.Ord2.Cob1}
\formal{GlobalTheory.Ord2.Cob2}
\formal{GlobalTheory.Ord2.Coc2}
\end{definition}

An $(\katoO,n,S)$-$2$-cocycle defines a class in $H^2(G_S,\katoO/\ell^n(1))$, and $(\katoO,n)$-triviality on $D_B$ is the
vanishing of its restriction to the decomposition group.

\begin{definition}[Tame relations and Euler systems]\label{pf:kato:def:es}
Let $\katoO$ be a coefficient ring containing the entries of $W$, $y\colon G_\Q\to\katoO^2$, $n\ge1$, $w$ a prime
and $\psi\colon G_\Q\to\Z/\ell^n$ a homomorphism. The \emph{tame relation} $\mathcal{T}(\katoO,W,y,n,w,\psi)$ holds if
there are a continuous
$Y\colon G_\Q\times\Z/\ell^n\to\katoO^2$, $P\in\katoO$, $u\in\katoQlb$ with $|u|=1$ and $b\in\katoO^2$ such that
\begin{enumerate}[label=(T\arabic*)]
\item\label{pf:kato:T1} $Y(st,x)=Y(s,x)+W(s)\,Y(t,x-\psi(s))$ for all $s,t,x$;
\item\label{pf:kato:T2} $\sum_{x}Y(s,x)=P\,y(s)+W(s)b-b$ for all $s$;
\item\label{pf:kato:T3} $P\katoeq nu\det(W(\sigma)-1)$ for every place $B$ over $w$ and every Frobenius $\sigma$ at $B$
  for $w$;
\item\label{pf:kato:T4} for every prime $q\notin\{w,\ell\}$ there are a place $B$ over $q$ and $t\colon\Z/\ell^n\to\katoO^2$
  with $Y(s,x)=W(s)t(x)-t(x)$ for $s\in I_B$ and all $x$.
\end{enumerate}
$W$ \emph{has an Euler system} at $A$, written $\mathrm{ES}(W,A)$, if there are a coefficient ring $\katoO$
containing the entries of $W$, a cocycle
$y\in\katoZ(G_\Q,W)$ with values in $\katoO^2$ such that $\mathrm{Unr}(W,y)$ holds and $\mathrm{Gr}(W,A,y)$ fails,
and a finite set $S'$ of primes such that $\mathcal{T}(\katoO,W,y,n,w,\psi)$ holds for all $n\ge1$, all primes
$w\notin S'$ with $\ell^n\mid w-1$, and all surjective $\psi\colon G_\Q\to\Z/\ell^n$ killing $\mu_w$.
\formallabel{GlobalTheory.Ord2.ES}
\end{definition}

By Shapiro's lemma, \ref{pf:kato:T1} says that $Y$ is a cocycle of the induced module
$W\otimes\katoO[\Z/\ell^n]$, i.e.\ a class over the cyclic extension of degree $\ell^n$ cut out by $\psi$; \ref{pf:kato:T2}
says that its corestriction is $P\cdot y$, and \ref{pf:kato:T3} that $P$ is, up to a unit, the Euler factor
$\det(1-W(\mathrm{Frob}_w))$ modulo $\ell^n$. This is an Euler system ``of rank one, one prime at a time'' in the
sense of \cite{Kolyvagin1990,Rubin2000}.

\begin{definition}[Unramified sources and the existence statement]\label{pf:kato:def:zetaur}
A newform $g$ is an \emph{unramified source} for $W$ at $A$, written $\mathrm{Src}^{\mathrm{ur}}(\iota,e,A,W,g)$, if $g$
is a newform of weight $2e+1$ whose level is prime to $\ell$,
$|\iota^{-1}(a_\ell(g))|=1$, and for some character $\chi$ killing $\mu_\ell$, $g$ is attached to
$W\otimes\varepsilon^{-e}\delta^{-1}\chi$ via $\iota$ and $W(s)_{11}\varepsilon(s)^e\delta(s)^{-1}\chi(s)=1$ for $s\in I_A$.
The \emph{existence statement} \textup{(Z)} is the proposition: for all $\ell,\Omega,\iota,e,S,A,W,g$, $\mathrm{NF}(e,S,A,W)$ and
$\mathrm{Src}^{\mathrm{ur}}(\iota,e,A,W,g)$ imply $\mathrm{ES}(W,A)$.
\formallabel{GlobalTheory.Ord2.ZetaUr}
\end{definition}

The rest of this subsection and the next five prove:

\begin{theorem}\label{pf:kato:vanishordl}
The statements of \cref{pf:kato:hyp:litcoh} and \textup{(Z)} imply \textup{(V)}.
\formallabel{GlobalTheory.Ord2.vanishOrdL_of_zetaUr}
\end{theorem}

\begin{theorem}\label{pf:kato:esvanish}
Assume the statements of \cref{pf:kato:hyp:litcoh}. If $\mathrm{NF}(e,S,A,W)$ and $\mathrm{ES}(W,A)$ hold, then
$\katoHf(\Q,W)=0$.
\formallabel{GlobalTheory.Ord2.es_vanish}
\end{theorem}

\begin{lemma}\label{pf:kato:nf}
Let $\ell\ne2$, $e\ge1$, $W$ admitted with $\mathrm{Src}(\iota,e,W)$. There are a finite set $S$, a place $A$ over
$\ell$, a representation $W'$ isomorphic to $W$ and a newform $g$ with $\mathrm{NF}(e,S,A,W')$ and
$\mathrm{Src}^{\mathrm{ur}}(\iota,e,A,W',g)$.
\formallabel{GlobalTheory.Ord2.exists_nf_srcUr}
\end{lemma}

\begin{proof}
Take $g$, $\chi$, $A$, $\kappa$ from clause (1) of $\mathrm{Src}$ and $h_1\in\GL_2(\katoQlb)$ triangularizing $W$ on $D_A$
with diagonal $(\kappa_0,\kappa_1)$. By \cref{pf:kato:def:admitted}(e), $W$ has level $M\ge1$ away from $\ell$; let $S$ be
the set of primes dividing $\ell M$. For a prime $q\notin S$ there is a place $B$ over $q$ on whose inertia group
some conjugate of $W$ is $\mathrm{diag}(\chi_0,\chi_1)$ with $\chi_i$ killing $\mu_M$; as $q\nmid M$, $I_B$ fixes $\mu_M$, so
$W$ is trivial on $I_B$, hence on all inertia groups over $q$ (they are conjugate).

The continuous representation $r=h_1Wh_1^{-1}$ of the compact group $G_\Q$ has entries in a finite extension $E$ of
$\Q_\ell$ (\lref{GlobalTheory.CompactRepField.exists_finiteDimensional_entries_mem}), and there is an upper
triangular $h_2\in\GL_2(\katoQlb)$ such that $h_2rh_2^{-1}$ has entries in the ring of integers $\katoO_E$ of $E$
(\lref{GlobalTheory.T.exists_upperTriangular_conj_integral_compact}: the image of $r$ stabilises a lattice, which
can be reached by an upper triangular change of basis). Put $h=h_2h_1$ and $W'=hWh^{-1}$; $\katoO_E$ is a coefficient
ring. Then $W'\cong W$, $\det W'=\delta$, $W'$ is continuous, integral, trivial on inertia outside $S$, odd with
$\delta^2=1$, and irreducible on $H_0(W')=H_0(W)$ (conjugation preserves irreducibility). On $D_A$,
$r$ is upper triangular with diagonal $(\kappa_0,\kappa_1)$, and conjugation by the upper triangular $h_2$ preserves this:
$W'(s)_{10}=0$, $W'(s)_{00}=\kappa_0(s)$, $W'(s)_{11}=\kappa_1(s)$. For $s\in I_A$, clause (1) of $\mathrm{Src}$ gives
$\kappa_1\varepsilon^e\delta^{-1}\chi=1$, and $\kappa_0\kappa_1=\delta$; hence $W'(s)_{00}=\varepsilon(s)^e\chi(s)$. As $\chi$ kills $\mu_\ell$,
$\chi^{\ell-1}=1$, so $W'(s)_{00}^{\ell-1}=\varepsilon(s)^{e(\ell-1)}$ on $I_A$, with $\ell-1\ge1$. Clause (2) of $\mathrm{Src}$,
applied to $A$ and $\kappa$, gives $s\in D_A$ with $W'(s)_{00}=\kappa_0(s)\ne\varepsilon(s)$. This proves $\mathrm{NF}(e,S,A,W')$.
Finally $g$ is attached to $W'\otimes\varepsilon^{-e}\delta^{-1}\chi$, which is conjugate to $W\otimes\varepsilon^{-e}\delta^{-1}\chi$
(attachment depends only on characteristic polynomials and semisimplicity), and
$W'(s)_{11}\varepsilon(s)^e\delta(s)^{-1}\chi(s)=1$ on $I_A$; this is $\mathrm{Src}^{\mathrm{ur}}(\iota,e,A,W',g)$.
\end{proof}

\begin{lemma}\label{pf:kato:conj}
Let $r,r'\colon G_F\to\GL_n(\katoQlb)$ be isomorphic. If $\katoHf(F,r)=0$ then $\katoHf(F,r')=0$.
\formallabel{UpperTheory.V.X.Kato.conj}
\end{lemma}

\begin{proof}
Let $r'=PrP^{-1}$. The bijection $f(x)=P^{-1}x$ satisfies $f(r'(s)x)=r(s)f(x)$; it maps cocycles of $r'$ to cocycles
of $r$, Bloch--Kato classes to Bloch--Kato classes (both local conditions are transported by $f$, resp.\ $f\otimes1$), and
$f(c)$ is a coboundary of $r$ if and only if $c$ is a coboundary of $r'$.
\end{proof}

% ---------------------------------------------------------------------------
\subsection{Algebraic lemmas}\label{pf:kato:sec:alg}

In this subsection $\mathrm{NF}(e,S,A,W)$ holds, $H_0=H_0(W)$ and $N_0=N_0(W)$. On $H_0$ we have $W^*=W$, and on
$N_0$ we have $W=W^*=1$; in particular a cocycle of $W$ or $W^*$ restricts to a homomorphism on $N_0$.

\begin{lemma}[Values of cocycles]\label{pf:kato:finext}
Let $K_0\subset\katoQlb$ be a finite extension of $\Q_\ell$ containing the entries of $W$, and $z\in\katoZ(G_\Q,W)$. There is a finite
extension $E\supset K_0$ of $\Q_\ell$ in $\katoQlb$ such that $z$ has values in $E^2$.
\formallabel{GlobalTheory.CocInt.values_finExt}
\end{lemma}

\begin{proof}
There is a countable family $(K_i)_{i\in I}$ of finite extensions of $\Q_\ell$ in $\katoQlb$ such that every finite subset of $\katoQlb$
lies in some $K_i$ (\lref{GlobalTheory.CompactRepField.exists_countable_family}). For $i\in I$ let
$H_i=\{g\in G_\Q:z(g)\in(K_0K_i)^2\}$. By the cocycle relations $z(gh)=z(g)+W(g)z(h)$ and $z(g^{-1})=-W(g^{-1})z(g)$, $H_i$ is a subgroup;
it is closed, as $z$ is continuous and finite-dimensional $\Q_\ell$-subspaces of $\katoQlb$ are closed; and the $H_i$ cover $G_\Q$. By
the Baire category theorem for the compact Hausdorff group $G_\Q$, some $H_i$ has non-empty interior, hence is open and of finite
index. Let $g_1,\dots,g_m$ represent $G_\Q/H_i$ and let $E$ be generated by $K_0K_i$ and the entries of the $z(g_j)$. Every $g\in G_\Q$
is $g_jh^{-1}$ with $h\in H_i$, and $z(g)=z(g_j)+W(g_j)z(h^{-1})\in E^2$.
\end{proof}

\begin{lemma}[Integral scaling]\label{pf:kato:scaling}
Let $\katoO$ be a coefficient ring containing the entries of $W$, $z\in\katoZ(G_\Q,W)$ with $\mathrm{Unr}(W,z)$, and $A$ a
place over $\ell$. There are a coefficient ring $\katoO'\supset\katoO$ and $k\ge0$ such that: $\ell^kz$ has values in
$\katoO'^2$; for every prime $q\ne\ell$ and \emph{every} place $B$ over $q$ there is $x\in\katoO'^2$ with
$\ell^kz(s)=W(s)x-x$ for $s\in I_B$; if $\mathrm{Str}(W,A,z)$ there is $x\in\katoO'^2$ with $\ell^kz(s)=W(s)x-x$ for
$s\in D_A$; and if $\mathrm{Gr}(W,A,z)$ there is $x\in\katoO'$ with $\ell^kz(s)_1=W(s)_{11}x-x$ for $s\in D_A$. The same holds
for $W^*$ in place of $W$.
\formallabel{GlobalTheory.T.galois_cocycle_integral_after_scaling}
\end{lemma}

\begin{proof}
For each $q\in S\setminus\{\ell\}$ choose a place $B_q$ and $x_q$ as in $\mathrm{Unr}$, and choose the vectors $x_T$, $x_G$
given by $\mathrm{Str}$ and $\mathrm{Gr}$ when these hold. For $q\notin S$ take $x_q=0$: there $W$ is trivial on $I_{B_q}$,
so $z=W(\cdot)x-x=0$ on $I_{B_q}$. By \cref{pf:kato:finext}, $z$ takes values in a finite extension of
$\Q_\ell$; enlarging it, let $E$ be a finite extension containing the values of
$z$, the entries of $W$, $\katoO$, and the finitely many vectors chosen, and $\katoO'$ its ring of integers. By
compactness $z$ is bounded, say by $C$, and so are the chosen vectors. For a place $B=gB_q$ over $q$, \cref{pf:kato:placeindep}
gives the vector $W(g)x_q-z(g)$, of norm $\le\max(C,|x_q|)$. Choose $k$ with $|\ell^k|\max(C,|x_q|,|x_T|,|x_G|)\le1$; then
all the scaled vectors lie in $\katoO'^2$. For $W^*$ the same argument applies, since $W^*$ is also continuous,
integral and unramified outside $S$.
\end{proof}

\begin{lemma}[Cup products]\label{pf:kato:cup}
Let $\katoO$ contain the entries of $W$, and $n\in\N$.
\begin{enumerate}[label=(\alph*)]
\item If $a\in\katoZ(G_\Q,W)$ and $b\in\katoZ(G_\Q,W^*)$ take values in $\katoO^2$ and satisfy $\mathrm{Unr}$, then $a$ and
  $b$ vanish on $N_S$, and $a\cup b$ is an $(\katoO,n,S)$-$2$-cocycle.
\item Let $H$ be a subgroup, $a$ an $(\katoO,n)$-cocycle of $W$ on $H$ and $b$ an $(\katoO,n)$-cocycle of $W^*$ on $H$. If $a$
  or $b$ is an $(\katoO,n)$-coboundary on $H$, then $a\cup b$ is $(\katoO,n)$-trivial on $H$.
\end{enumerate}
\formallabel{GlobalTheory.T.cupProductDualTwist_isTwoCocycle_unramified}
\formallabel{GlobalTheory.T.cup_product_coboundary_mod}
\end{lemma}

\begin{proof}
(a) Let $q\notin S$ and $B$ a place over $q$. By $\mathrm{Unr}$ and \cref{pf:kato:placeindep}, $a=W(\cdot)x-x$ on $I_B$, and $W$
is trivial on $I_B$, so $a$ vanishes on $I_B$. On $\ker W$ the cocycle $a$ is a homomorphism, and
$\{s\in\ker W:a(s)=0\}$ is a closed subgroup, normal in $G_\Q$: for $s$ in it,
$a(gsg^{-1})=a(g)+W(g)a(s)-W(gsg^{-1})a(g)=0$. It contains all inertia groups over primes not in $S$, hence $N_S$. The
same holds for $b$, since $W^*$ is trivial on these inertia groups ($\varepsilon$ is unramified at $q\ne\ell$). The identity
$\varepsilon(r)(a\cup b)(s,t)+(a\cup b)(r,st)=(a\cup b)(rs,t)+(a\cup b)(r,s)$ holds exactly, by the cocycle relations and
\eqref{pf:kato:eq:omega}; continuity and integrality are clear ($W^*$ has entries in $\katoO$ as $\delta=\pm1$). For $a',b'\in
N_S$, $a(sa')=a(s)$, $W^*(sa')=W^*(s)$ and $b(tb')=b(t)$, so $(a\cup b)(sa',tb')=(a\cup b)(s,t)$.

(b) If $a\katoeq nW(\cdot)x-x$ on $H$ with $x\in\katoO^2$, then by \eqref{pf:kato:eq:omega}, for $s,t\in H$,
\[
(a\cup b)(s,t)\katoeq n\varepsilon(s)\omega(x,b(t))-\omega\bigl(x,W^*(s)b(t)\bigr)\katoeq n\varepsilon(s)u(t)-u(st)+u(s),\qquad
u=\omega(x,b(\cdot)),
\]
using $W^*(s)b(t)\katoeq nb(st)-b(s)$. If $b\katoeq nW^*(\cdot)x-x$ on $H$, the same computation with
$u(s)=-\omega\bigl(a(s),W^*(s)x\bigr)$ gives $(a\cup b)(s,t)\katoeq n\varepsilon(s)u(t)-u(st)+u(s)$.
\end{proof}

\begin{lemma}\label{pf:kato:commdet}
There is $\phi\in[H_0,H_0]$ with $\det(W(\phi)-1)\ne0$.
\formallabel{GlobalTheory.T.exists_commutator_det_nonzero}
\end{lemma}

\begin{proof}
We prove more generally: if $k$ is an algebraically closed field, $\rho\colon\Gamma\to\GL_2(k)$ a homomorphism and $H\subset\Gamma$
a subgroup with $\det\rho=1$ on $H$ and $\rho$ irreducible on $H$, then $\det(\rho(\phi)-1)\ne0$ for some $\phi\in[H,H]$; for $W$
apply this to $H=H_0$.

First, if $X,Y\in\SL_2(k)$ satisfy $\det(X-1)=\det(Y-1)=\det(XY-1)=0$, then $(X-1)(Y-1)=0$. Indeed $\det(X-1)=\det X-\tr X+1$,
so $\tr X=\tr Y=\tr XY=2$, and $N=X-1$, $M=Y-1$ have trace and determinant $0$, so $N^2=M^2=0$; moreover
$\tr(NM)=\tr(XY-1-N-M)=0$. Suppose $NM\ne0$. Then $N,M\ne0$ and $\ker N=\operatorname{im}N=L_N$, $\ker M=\operatorname{im}M=L_M$ are lines. If
$L_N=L_M$ then $NM=0$; otherwise, for $v\in L_N\setminus0$, $Mv\in L_M\setminus0$ and $NMv\in L_N\setminus0$, so $NMv=cv$
with $c\ne0$ while $NM$ kills $L_M$; then $\tr(NM)=c\ne0$, a contradiction.

Now suppose $\det(\rho(\phi)-1)=0$ for all $\phi\in[H,H]$. As $[H,H]$ is a group, $(\rho(\phi)-1)(\rho(\psi)-1)=0$ for all
$\phi,\psi\in[H,H]$. Let $F$ be the space of vectors fixed by $\rho([H,H])$; it is $H$-stable since $[H,H]$ is normal in $H$. For
every $\psi\in[H,H]$ and $w\in k^2$, $(\rho(\psi)-1)w\in F$. If $F=0$ this forces $\rho(\psi)=1$ for all $\psi$, i.e.\ $F=k^2$; so
$F\ne0$, and by irreducibility $F=k^2$: $\rho([H,H])=1$ and $\rho(H)$ is commutative. For $g\in H$ an eigenspace of $\rho(g)$ is
$H$-stable, hence equal to $k^2$; so $\rho(H)$ consists of scalars, and every line is $H$-stable, contradicting
irreducibility.
\end{proof}

\begin{lemma}\label{pf:kato:pair}
Let $y\in\katoZ(G_\Q,W)$ and $z\in\katoZ(G_\Q,W^*)$ be both non-zero on $N_0$. There are $g,h\in N_0$ with
$\omega(y(h),z(g))\ne\omega(y(g),z(h))$.
\formallabel{GlobalTheory.T.cocycle_pair_alternating_ne_on_cycloKernel}
\end{lemma}

\begin{proof}
For $s\in H_0$ and $g\in N_0$ we have $sgs^{-1}\in N_0$, $y(sgs^{-1})=W(s)y(g)$ and $z(sgs^{-1})=W^*(s)z(g)=W(s)z(g)$
(cocycle relation, $W(g)=W^*(g)=1$, $W^*=W$ on $H_0$). Suppose $\omega(y(h),z(g))=\omega(y(g),z(h))$ for all $g,h\in N_0$. Let
$\Sigma\subset\katoQlb^2\times\katoQlb^2$ be the span of the pairs $(y(g),z(g))$, $g\in N_0$; it is stable under
$(x,x')\mapsto(W(s)x,W(s)x')$ for $s\in H_0$, and by bilinearity $\omega(y(g),\sigma_2)=\omega(\sigma_1,z(g))$ for $g\in N_0$ and
$\sigma=(\sigma_1,\sigma_2)\in\Sigma$. The space $K'=\{x':(0,x')\in\Sigma\}$ is $H_0$-stable; it is $0$, for otherwise it is everything,
and choosing $g$ with $y(g)\ne0$ and $x'$ with $\omega(y(g),x')\ne0$ gives $\omega(y(g),x')=\omega(0,z(g))=0$. The first projection
of $\Sigma$ is $H_0$-stable and non-zero, hence everything. So $\Sigma$ is the graph of a linear map $\varphi$ commuting with
$W(H_0)$; by irreducibility of $W$ on $H_0$ and Schur's lemma ($\katoQlb$ is algebraically closed), $\varphi=c$ is a scalar.
Then $z(g)=cy(g)$ on $N_0$, and $c\ne0$ since $z\ne0$ on $N_0$. For $g_0$ with $y(g_0)\ne0$ the hypothesis gives
$c\,\omega(y(h),y(g_0))=c\,\omega(y(g_0),y(h))$, i.e.\ $2c\,\omega(y(g_0),y(h))=0$, for all $h\in N_0$, hence
$\omega(y(g_0),\sigma_1)=0$ for all $\sigma\in\Sigma$. As the first projection of $\Sigma$ is everything, $y(g_0)=0$, a contradiction.
\end{proof}

\begin{lemma}\label{pf:kato:grkill}
A cocycle $y\in\katoZ(G_\Q,W)$ vanishing on $N_0$ satisfies $\mathrm{Gr}(W,A,y)$.
\formallabel{GlobalTheory.T.greenbergOrdinary_of_vanishing_on_cycloKernel}
\end{lemma}

\begin{proof}
Write $D=D_A$ and, for $s\in D$, $W(s)=\left(\begin{smallmatrix}\alpha(s)&\beta(s)\\0&\gamma(s)\end{smallmatrix}\right)$; then
$\alpha,\gamma$ are characters of $D$ and $y(st)_1=y(s)_1+\gamma(s)y(t)_1$ for $s\in D$. Let $U$ be the set of $u\in D$ with
$\alpha(u)=\gamma(u)=\varepsilon(u)=1$; it is a subgroup, and $\beta(uv)=\beta(u)+\beta(v)$, $y(uv)_1=y(u)_1+y(v)_1$ for $u,v\in U$. If
$u,v\in U$ and $\beta(u)=\beta(v)$, then $W(u)=W(v)$, so $u^{-1}v\in N_0$ and $y(v)=y(u)+W(u)y(u^{-1}v)=y(u)$.

As $\varepsilon$ maps $I_A$ onto $\Z_\ell^\times$ (\lref{GlobalTheory.CycInertia.cycZ_Ine_surjective}), choose $t\in I_A$ with
$\varepsilon(t)=1+\ell$. Let $m$ be as in $\mathrm{NF}$ and $T=t^{2m}$. Then $\alpha(T)=\varepsilon(t)^{2em}=(1+\ell)^{2em}=:q_0$, an integer
$>1$, and $\delta(T)=\delta(t)^{2m}=1$, so $\gamma(T)=q_0^{-1}$. For $u\in U$, conjugation gives $TuT^{-1}\in U$ with
$\beta(TuT^{-1})=\alpha(T)\gamma(T)^{-1}\beta(u)=q_0^2\beta(u)$ and $y(TuT^{-1})_1=\gamma(T)y(u)_1=q_0^{-1}y(u)_1$ (as
$y(T)_1+\gamma(T)y(T^{-1})_1=y(1)_1=0$). The element $u^{q_0^2}\in U$ has the same $\beta$, so $q_0^{-1}y(u)_1=q_0^2y(u)_1$;
since $q_0^3\ne1$, $y(u)_1=0$ for all $u\in U$.

For $s\in D$ the element $c=(sT)^{-1}(Ts)$ lies in $U$ (its diagonal entries and its cyclotomic character are $1$). Hence
$y(Ts)_1=y(sT)_1+\gamma(sT)y(c)_1=y(sT)_1$, i.e.\ $y(T)_1+\gamma(T)y(s)_1=y(s)_1+\gamma(s)y(T)_1$. As $\gamma(T)\ne1$, this gives
$y(s)_1=\gamma(s)x-x$ with $x=y(T)_1/(\gamma(T)-1)$.
\end{proof}

\begin{lemma}\label{pf:kato:hompadic}
Let $\Gamma$ be a compact group, $\katoO$ a coefficient ring, $r\colon\Gamma\to\GL_2(\katoO)$ a continuous irreducible representation
with $(\det r)^2=1$, and $\lambda\colon\Gamma\to\Z_\ell$ a continuous homomorphism with $\ker r\subset\ker\lambda$. Then $\lambda=0$.
\formallabel{GlobalTheory.T.hom_padicInt_trivial_of_irreducible_compact_rep}
\end{lemma}

\begin{proof}
For a $2\times2$ matrix $M$ write $|M|\le c$ if all entries have norm $\le c$, and put $\rho=|\ell|$. We use two statements of the
formal development: (A) (\lref{GlobalTheory.T.exists_augmentation_combination_eq_id_traceless_irreducible}) as $r$ is
irreducible, there are $\gamma_1,\dots,\gamma_N\in\Gamma$ and $a_1,\dots,a_N\in\Q_\ell$ with $X=\sum_ia_i\bigl(r(\gamma_i)Xr(\gamma_i)^{-1}-X\bigr)$
for every matrix $X$ of trace $0$; (B) (\lref{GlobalTheory.T.galoisHom_zl_lipschitz_of_ker_le_linearGroup}) as
$\ker r\subset\ker\lambda$, there is $C$ such that $|\lambda(x)|\le C|\ell^m|$ whenever $|r(x)-1|\le|\ell^m|$. We also use the elementary
estimates, valid for $|Z|\le\epsilon\le1$, $|Z'+Z|\le\epsilon^2$, integral $A$ with integral inverse and $m,\beta\in\N$:
$|(1+Z)^m-1-mZ|\le\epsilon^2$, $|\prod_i(1+Z_i)-1-\sum_iZ_i|\le\epsilon^2$ if all $|Z_i|\le\epsilon$, and
$|A(1+Z)^\beta A^{-1}(1+Z')^\beta-1-\beta(AZA^{-1}-Z)|\le\epsilon^2$.

Suppose $\lambda(g)\ne0$. As $\GL_2(\katoO/\ell\katoO)$ is finite and $(\det r)^2=1$, there is $k\ge1$ with $\det r(g^k)=1$ and
$|r(g^k)-1|\le\rho$; put $h_0=g^k$ and $v_0=\lambda(h_0)=k\lambda(g)\ne0$. For $n\ge0$ let $h=h_0^{\ell^n}$; by induction, using the first
estimate, $\det r(h)=1$ and $Z=r(h)-1$ satisfies $|Z|\le\epsilon=\rho^{n+1}$. As $\det(1+Z)=1$, $t=\tr Z=-\det Z$ has $|t|\le\epsilon^2$, and
$Z'=r(h^{-1})-1=t-Z$. Choose $s$ with $b_i=\ell^sa_i\in\Z_\ell$ for all $i$, and integers $\beta_i\ge0$ with
$|\beta_i-b_i|\le\rho^{2n+1}$. Put
\[
  P=h^{-\ell^s}\prod_i\gamma_ih^{\beta_i}\gamma_i^{-1}h^{-\beta_i}.
\]
Since $\lambda$ takes values in an abelian group, $\lambda(P)=-\ell^{s+n}v_0$, so $|\lambda(P)|=\rho^{s+n}|v_0|$. On the other hand, with
$D_i=r(\gamma_i)Zr(\gamma_i)^{-1}-Z$ (so $|D_i|\le\epsilon$) and $X=Z-\frac t2$, of trace $0$, (A) gives $\sum_ib_iD_i=\ell^sX$, and the three
estimates give
\[
  r(P)-1\equiv\ell^sZ'+\sum_i\beta_iD_i\equiv\ell^s(Z'+X)=\ell^s\tfrac t2\qquad\text{up to matrices of size}\le\rho^{2n+1},
\]
while $|\ell^s\frac t2|\le\epsilon^2/\rho=\rho^{2n+1}$. Hence $|r(P)-1|\le\rho^{2n+1}$, and (B) gives $\rho^{s+n}|v_0|\le C\rho^{2n+1}$ for all
$n$, i.e.\ $|v_0|\le C\rho^{n+1-s}$. Letting $n\to\infty$, $v_0=0$, a contradiction.
\end{proof}

\begin{lemma}\label{pf:kato:kereps}
There is $a\in\ker W$ with $\varepsilon(a)\ne1$.
\formallabel{GlobalTheory.T.exists_ker_rep_cyclotomic_ne_one_normalForm}
\end{lemma}

\begin{proof}
Suppose $\ker W\subset\ker\varepsilon$. Let $\lambda\colon G_\Q\to\Z_\ell$ be the composite of $\varepsilon$ with the homomorphism
$\Z_\ell^\times\to\Z_\ell$ given by the $\ell$-adic logarithm; it is continuous and $\ker W\subset\ker\lambda$. By \cref{pf:kato:hompadic}, applied to $W$ (irreducible on $G_\Q$, since irreducible
on $H_0$), $\lambda=0$. But for
$\ell$ odd $\lambda$ is non-zero, as $\varepsilon(G_\Q)=\Z_\ell^\times$ contains elements whose logarithm is non-zero.
\end{proof}

\begin{lemma}\label{pf:kato:dualkill}
A cocycle $z\in\katoZ(G_\Q,W^*)$ vanishing on $N_0$ is a coboundary.
\formallabel{GlobalTheory.T.tateDual_cocycle_coboundary_of_vanish_on_fixing}
\end{lemma}

\begin{proof}
Take $a$ as in \cref{pf:kato:kereps}. Then $\delta(a)=1$, so $W^*(a)=\varepsilon(a)$ is a scalar. For $g\in G_\Q$, the element
$(ga)^{-1}(ag)$ lies in $N_0$ ($W(a)=1$, and $\varepsilon$ is abelian), so $z(ag)=z(ga)$. Expanding,
$z(a)+\varepsilon(a)z(g)=z(g)+W^*(g)z(a)$, i.e.\ $z(g)=W^*(g)x-x$ with $x=(\varepsilon(a)-1)^{-1}z(a)$.
\end{proof}

\begin{lemma}[Heights and commutators]\label{pf:kato:htcomm}
Let $\katoO$ contain the entries of $W$, $y\in\katoZ(G_\Q,W)$, $z\in\katoZ(G_\Q,W^*)$, $u\colon G_\Q\to\katoQlb$ and $n\in\N$ with
$(y\cup z)(s,t)\katoeq n\varepsilon(s)u(t)-u(st)+u(s)$ for all $s,t$. For $\phi\in H_0$ and $g,h\in N_0$, writing $\katoht=\katoht_{y,z,u}$,
\[
  \katoht\bigl(\phi\,[g,h]\bigr)-\katoht(\phi)\katoeq n\det\bigl(W(\phi)-1\bigr)\bigl(\omega(y(h),z(g))-\omega(y(g),z(h))\bigr),
\]
where $[g,h]=ghg^{-1}h^{-1}$.
\formallabel{GlobalTheory.T.height_commutator_invariance}
\end{lemma}

\begin{proof}
Put $k=[g,h]\in N_0$. On $N_0$, $y$ and $z$ are homomorphisms, so $y(k)=z(k)=0$; hence $W(\phi k)=W(\phi)$, $y(\phi k)=y(\phi)$,
$z(\phi k)=z(\phi)$ and $\katoht(\phi k)-\katoht(\phi)=\det(W(\phi)-1)(u(\phi k)-u(\phi))$. Let
$E(s,t)=(y\cup z)(s,t)-\varepsilon(s)u(t)+u(st)-u(s)$, so $E\katoeq n0$. Using $z(k)=y(k)=0$, $W^*=1$ on $N_0$, $\varepsilon(\phi)=1$ and
$khg=gh$, a direct computation gives
\[
  u(\phi k)-u(\phi)-\bigl(\omega(y(h),z(g))-\omega(y(g),z(h))\bigr)=E(\phi,k)+E(g,h)-E(h,g)-E(k,hg)\katoeq n0.
\]
Multiplying by $\det(W(\phi)-1)$, which has norm $\le1$, gives the claim.
\end{proof}

\begin{lemma}[Depth of a non-coboundary]\label{pf:kato:depth}
Let $H$ be a set and $k,a\colon H\to\katoQlb$ maps such that $k(s_0)\ne1$ for some $s_0\in H$ and there is no $X\in\katoQlb$ with
$a(s)=k(s)X-X$ for all $s\in H$. Then there is $r\in\N$ such that for all $n\in\N$ and $x,y\in\katoQlb$ with $|x|\le1$: if
$|xa(s)-(k(s)y-y)|\le|\ell^n|$ for all $s\in H$, then $|x|\le|\ell|^{\max(n-r,0)}$.
\formallabel{GlobalTheory.T.exists_annihilatorDepth_of_not_coboundary}
\end{lemma}

\begin{proof}
Put $q=|\ell|<1$ and suppose that no $r$ has the property. Then for every $r$ there are $n,x,y$ with $|x|\le1$,
$|xa(s)-(k(s)y-y)|\le q^n$ on $H$ and $|x|>q^{\max(n-r,0)}$. Then $n\ge r$ (otherwise $|x|>1$), so $x\ne0$, and $X_r=y/x$ satisfies
$|a(s)-(k(s)X_r-X_r)|\le q^n/|x|<q^r$ for $s\in H$. Let $X=a(s_0)/(k(s_0)-1)$. The case $s=s_0$ gives $|X-X_r|\le q^r/|k(s_0)-1|$, and
for $s\in H$ the identity $a(s)-(k(s)X-X)=\bigl(a(s)-(k(s)X_r-X_r)\bigr)-(k(s)-1)(X-X_r)$ gives
$|a(s)-(k(s)X-X)|\le(1+|k(s)-1|/|k(s_0)-1|)\,q^r$. Letting $r\to\infty$, $a(s)=k(s)X-X$ for all $s\in H$, a contradiction.
\end{proof}

\begin{lemma}[Bounded multipliers]\label{pf:kato:boundedtors}
Let $\katoO$ be a coefficient ring containing the entries of $W$ and $Y\colon G_\Q\to\katoO^2$ a cocycle of $W$ which is not a
coboundary. There is $k_0\in\N$ such that for all $n\in\N$ and $a\in\katoO$: if $aY$ is an $(\katoO,n)$-coboundary on $G_\Q$, then
$a\ell^{k_0}\katoeq n0$.
\formallabel{GlobalTheory.T.nonzero_class_bounded_torsion}
\end{lemma}

\begin{proof}
For $s\in G_\Q$ and $i\in\{0,1\}$ consider the linear form $\varphi_{s,i}(a,x)=aY(s)_i-(W(s)x-x)_i$ on $\katoQlb\times\katoQlb^2$. If
$\varphi_{s,i}(a,x)=0$ for all $s,i$ and $a\ne0$, then $Y(s)=W(s)(a^{-1}x)-a^{-1}x$ for all $s$, which is excluded. So the common
kernel of the $\varphi_{s,i}$ is contained in the kernel of the linear form $(a,x)\mapsto a$, which is therefore (the space being
finite-dimensional) a finite linear combination $a=\sum_jc_j\varphi_{s_j,i_j}(a,x)$. With $C=\sum_j|c_j|$, $|a|\le Cr$ whenever
$|\varphi_{s,i}(a,x)|\le r$ for all $s,i$. Choose $k_0$ with $C|\ell|^{k_0}\le1$. If $aY(s)\katoeq nW(s)x-x$ for all $s$, then
$|a|\le C|\ell^n|$, so $|a\ell^{k_0}|\le|\ell^n|$.
\end{proof}

\begin{lemma}[Divisible chains]\label{pf:kato:divchain}
Let $X$ be an abelian group in which every element is killed by a power of $\ell$, whose $\ell$-torsion subgroup $X[\ell]$ is finite,
and such that for every $m$ there is $x\in X$ with $\ell^mx\ne0$. Then there are $x_0,x_1,\ldots\in X$ with $x_0\ne0$ and
$\ell x_{n+1}=x_n$ for all $n$.
\formallabel{GlobalTheory.T.exists_divisible_chain_primary_torsion_group}
\end{lemma}

\begin{proof}
Call $y\in X$ \emph{infinitely divisible} if $y\in\ell^mX$ for every $m$.

(1) Some nonzero $y\in X[\ell]$ is infinitely divisible. Otherwise choose, for each nonzero $y\in X[\ell]$, an integer $m(y)$ with
$y\notin\ell^{m(y)}X$, and $M\ge m(y)$ for all $y$ ($X[\ell]$ is finite). Choose $x$ with $\ell^Mx\ne0$ and $k\ge0$ with
$\ell^{k+M}x\ne0=\ell^{k+1+M}x$. Then $y=\ell^{k+M}x$ is a nonzero element of $X[\ell]$ and $y=\ell^{m(y)}\bigl(\ell^{M-m(y)+k}x\bigr)$, a
contradiction.

(2) If $y$ is infinitely divisible, then $y=\ell z$ for some infinitely divisible $z$. The set $F=\{z:\ell z=y\}$ is non-empty and is
a coset of $X[\ell]$, hence finite. Otherwise choose, for each $z\in F$, an integer $m(z)$ with $z\notin\ell^{m(z)}X$, and $M\ge m(z)$
for all $z\in F$. Write $y=\ell^{M+1}w$; then $z=\ell^Mw$ lies in $F$ and in $\ell^{m(z)}X$, a contradiction.

Let $x_0$ be given by (1), and choose $x_{n+1}$ by (2) applied to $x_n$.
\end{proof}

% ---------------------------------------------------------------------------
\subsection{Local lemmas}\label{pf:kato:sec:local}

In this subsection $\katoO$ is a coefficient ring, $n\ge1$, and $W$ is continuous with entries in $\katoO$ and
$\delta^2=1$, so that $W^*$ has entries in $\katoO$.

\begin{lemma}[Unramified cup products]\label{pf:kato:unrcup}
Let $q\ne\ell$ be a prime and $B$ a place over $q$. Let $a$ be an $(\katoO,n)$-cocycle of $W$ on $D_B$ and $b$ an
$(\katoO,n)$-cocycle of $W^*$ on $D_B$, both $(\katoO,n)$-coboundaries on $I_B$. Then $a\cup b$ is $(\katoO,n)$-trivial
on $D_B$.
\formallabel{GlobalTheory.T.cupProduct_unramifiedClasses_vanishes}
\end{lemma}

\begin{proof}
Write $a=a_0+a_1$ with $a_0(s)=W(s)x-x$ and $a_1\katoeq n0$ on $I_B$, and similarly $b=b_0+b_1$. By bilinearity
$a\cup b=a_0\cup b+a_1\cup b_0+a_1\cup b_1$, and the first two terms are $(\katoO,n)$-trivial on $D_B$ by
\cref{pf:kato:cup}(b). It remains to treat $\xi=a_1\cup b_1$, where $a_1,b_1$ vanish modulo $\ell^n$ on $I_B$.

Since $I_B$ is normal in $D_B$, for $s\in D_B$ and $i\in I_B$ we have $a_1(si)\katoeq na_1(s)$ and, writing $is=si'$
with $i'\in I_B$, $W(i)a_1(s)\katoeq na_1(is)=a_1(si')\katoeq na_1(s)$; similarly for $b_1$. Hence
$\xi(si,t)\katoeq n\xi(s,t)\katoeq n\xi(s,ti)$ for $s,t\in D_B$, $i\in I_B$; also $\varepsilon=1$ on $I_B$ as $q\ne\ell$. By
continuity of $\xi$ on the compact group $D_B\times D_B$, there is an open subgroup $K\subset D_B$ containing $I_B$ such
that $\varepsilon\katoeq n1$ on $K$ and $\xi(sk,t)\katoeq n\xi(s,t)\katoeq n\xi(s,tk)$ for $s,t\in D_B$, $k\in K$
(uniform continuity on a compact group). As $D_B/I_B$ is procyclic, topologically generated by
the image of a Frobenius element, there are $\phi\in D_B$ and $m\ge1$ with $\phi^m\in K$ and a continuous homomorphism
$d\colon D_B\to\Z/m\ell^n$ with $\phi^{-\katorep{d(s)}}s\in K$ for all $s\in D_B$, where $\katorep{\cdot}$ is the representative
in $[0,m\ell^n)$ (take a Frobenius element $\phi$, $m$ the order of its image in
the finite cyclic group $D_B/K$, and $d$ the composite of $D_B\to D_B/I_B\cong\widehat\Z$, $\phi\mapsto1$, with reduction modulo $m\ell^n$). Define
$T(0)=\xi(1,1)$ and $T(j+1)=T(j)-\xi(\phi^j,\phi)$. The cocycle relation gives, by induction on $j$,
\[
  \xi(\phi^i,\phi^j)\katoeq n\varepsilon(\phi)^iT(j)-T(i+j)+T(i)\qquad(i,j\ge0),
\]
and $T(j+m\ell^n)\katoeq nT(j)$: by the $K$-invariance of $\xi$ and $\phi^m\in K$, the increments $T(j+1)-T(j)=-\xi(\phi^j,\phi)$ are
$m$-periodic modulo $\ell^n$, so $T(j+m\ell^n)-T(j)$ is congruent to $\ell^n$ times an element of $\katoO$.
Put $u(s)=T(\katorep{d(s)})$ for $s\in D_B$; it is locally constant. For $s=\phi^ik_1$, $t=\phi^jk_2$ with
$i=\katorep{d(s)}$, $j=\katorep{d(t)}$ and $k_1,k_2\in K$, the invariance of $\xi$ gives $\xi(s,t)\katoeq n\xi(\phi^i,\phi^j)$, and
$\katorep{d(st)}\equiv i+j$ modulo $m\ell^n$; combining with the two displayed properties of $T$ and $\varepsilon(k_1)\katoeq n1$ gives
$\xi(s,t)\katoeq n\varepsilon(s)u(t)-u(st)+u(s)$.
\end{proof}

\begin{remark}
The formal statement of \cref{pf:kato:unrcup} takes the statements of \cref{pf:kato:hyp:litcoh} as a hypothesis, but
its proof does not use them: it only uses that $H^2$ of a procyclic group with torsion coefficients vanishes, in the
explicit form above.
\end{remark}

\begin{lemma}\label{pf:kato:epsdec}
Let $w$ be a prime with $\ell^n\mid w-1$ and $B$ a place over $w$. Then $\varepsilon(s)\katoeq n1$ for all $s\in D_B$, and
$\varepsilon=1$ on $I_B$.
\formallabel{GlobalTheory.T.cyclotomicChar_congruent_one_on_decompositionGroup}
\end{lemma}

\begin{proof}
As $w\ne\ell$, reduction modulo the maximal ideal of $B$ is injective on $\ell$-power roots of unity, inertia acts
trivially on them, and an element of $D_B$ acts on the finite group $\mu_{\ell^n}$ as some power of $x\mapsto x^w$, which is
the identity on $\mu_{\ell^n}$ since $w\equiv1\bmod\ell^n$.
\end{proof}

\begin{lemma}[Wild inertia]\label{pf:kato:wild}
Let $B$ be a place of $\Omega$ at which $\ell$ is a unit, $M$ an abelian group killed by $\ell^n$, and $F\colon G_\Q\to M$ a map whose
restriction to $I_B$ is locally constant and satisfies $F(st)=F(s)+F(t)$ for $s,t\in I_B$. Let $\sigma\in G_\Q$ be \emph{wild} at $B$,
i.e.\ $\sigma(y)/y-1$ lies in the maximal ideal of $B$ for every $y\in\Omega^\times$. Then $F(\sigma)=0$.
\formallabel{GlobalTheory.TameConjInv.eq_zero_of_wild}
\end{lemma}

\begin{proof}
As $\sigma(y)/y\in B^\times$ for all $y\ne0$, $\sigma(B)=B$, and $\sigma(y)-y=y(\sigma(y)/y-1)$ lies in the maximal ideal for $y\in B$;
so $\sigma\in I_B$. By additivity $F(1)=0$; as $F$ is locally constant on $I_B$, there is a finite extension $L\subset\Omega$ of $\Q$ with
$F=0$ on $I_B\cap G_L$. As $G_L$ has finite index in $G_\Q$, $\sigma^N\in G_L$ for some $N\ge1$; write $N=\ell^am$ with $\ell\nmid m$.
Powers of wild elements are wild (\lref{GlobalTheory.WildInertiaChar.wild_pow}), so $\tau=\sigma^m$ is wild, and $\tau^{\ell^a}$
fixes $L$. If $\tau$ is wild, $k\in\N$ is a unit at $B$ and $\tau^k$ fixes $x$, then $\tau$ fixes $x$
(\lref{GlobalTheory.WildInertiaChar.fixed_of_pow_fixed}; indeed, if $d=\tau x-x\ne0$, then
$0=(\tau^kx-x)/d=\sum_{i<k}\tau^i(d)/d\equiv k$ modulo the maximal ideal). Hence $\tau$ fixes $L$, $\tau\in I_B\cap G_L$, and
$mF(\sigma)=F(\sigma^m)=0$. Since also $\ell^nF(\sigma)=0$ and $\gcd(m,\ell^n)=1$, $F(\sigma)=0$.
\end{proof}

\begin{lemma}[Tame inertia]\label{pf:kato:tameconj}
Let $w$ be a prime with $\ell^n\mid w-1$, $B$ a place over $w$, $M$ an abelian group killed by $\ell^n$ and $f\colon
G_\Q\to M$ a map whose restriction to $I_B$ is a locally constant homomorphism. Then $f(gsg^{-1})=f(s)$ for
$g\in D_B$ and $s\in I_B$.
\formallabel{GlobalTheory.T.inertiaTorsionQuotient_conj_invariant_decomposition}
\end{lemma}

\begin{proof}
Let $P_B\subset I_B$ be the wild inertia group (the $\sigma$ with $\sigma(x)/x\equiv1$ modulo the maximal ideal of $B$ for all
$x\ne0$). We use two standard facts of local class field theory and ramification theory \cite[Ch.~IV]{Serre1979}, in the
form of the formal lemmas \lref{GlobalTheory.TameInertiaGen.exists_tameGenAt} and
\lref{GlobalTheory.WeilDeligne.isTameChar_tameL}, and \cref{pf:kato:wild}:
(i) there is $\gamma\in I_B$ such that $I_B$ is the closure of the subgroup generated by $\gamma$ and $P_B$;
(ii) let $t\colon G_\Q\to\Z/\ell^n$ be the Kummer character of $w$ (for a compatible system of $\ell$-power roots of $w$,
reduced modulo $\ell^n$); its restriction to $I_B$ is a locally constant homomorphism, $t(gsg^{-1})=\varepsilon(g)t(s)$ for
$g\in G_\Q$ and $s\in I_B$, and $t(\tau)=1$ for some $\tau\in I_B$;
(iii) a locally constant homomorphism from $I_B$ to an abelian group killed by $\ell^n$ vanishes on $P_B$ (\cref{pf:kato:wild};
$\ell$ is a unit at $B$ as $w\ne\ell$).

Let $F=(f,t)\colon I_B\to M\times\Z/\ell^n$. By (iii) $F$ vanishes on $P_B$; as $F$ is a locally constant homomorphism,
(i) shows that for every $s\in I_B$ there is $k\in\Z$ with $F(s)=kF(\gamma)$. Applied to $\tau$, this gives $1=k_0t(\gamma)$,
so $t(\gamma)$ is a unit of $\Z/\ell^n$. Let $g\in D_B$ and $s\in I_B$; then $gsg^{-1}\in I_B$ and, by (ii) and
\cref{pf:kato:epsdec}, $t(gsg^{-1})=\varepsilon(g)t(s)=t(s)$. Writing $F(s)=k_1F(\gamma)$ and $F(gsg^{-1})=k_2F(\gamma)$, the
second components give $(k_2-k_1)t(\gamma)=0$, so $\ell^n\mid k_2-k_1$, and $f(gsg^{-1})-f(s)=(k_2-k_1)f(\gamma)=0$.
\end{proof}

\begin{lemma}[Tame cup products]\label{pf:kato:brtame}
Let $w$ be a prime with $\ell^n\mid w-1$, $B$ a place over $w$, and $\psi\colon G_\Q\to\Z/\ell^n$ a homomorphism with
$\psi(\tau)=1$ for some $\tau\in I_B$. Let $h\colon G_\Q\to\katoO$ satisfy $h(st)\katoeq nh(s)+h(t)$ for $s,t\in D_B$ and
$h\katoeq n0$ on $I_B$. If $(s,t)\mapsto\katorep{\psi(s)}h(t)$ is $(\katoO,n)$-trivial on $D_B$, then $h(\sigma)\katoeq n0$ for every
$\sigma\in D_B$.
\formallabel{GlobalTheory.T.tame_cochain_cocycle_extension}
\end{lemma}

\begin{proof}
Let $u$ be the cochain with $\katorep{\psi(s)}h(t)\katoeq n\varepsilon(s)u(t)-u(st)+u(s)$ for $s,t\in D_B$. For $s,t\in I_B$, $h(t)\katoeq n0$
and $\varepsilon(s)=1$, so $u(st)\katoeq nu(s)+u(t)$: the reduction of $u$ modulo the ball of radius $|\ell^n|$ is a locally
constant homomorphism on $I_B$ (to the quotient of the closed unit ball by this ball, which is killed by $\ell^n$). Let
$\sigma\in D_B$ and $\tau'=\sigma\tau\sigma^{-1}\in I_B$. By \cref{pf:kato:tameconj}, $u(\tau')\katoeq nu(\tau)$. As $\psi$ is a homomorphism
to an abelian group, $\psi(\tau')=\psi(\tau)=1$, and $\tau'\sigma=\sigma\tau$. The relation at $(\tau',\sigma)$ and at $(\sigma,\tau)$ reads
\[
  h(\sigma)\katoeq n\varepsilon(\tau')u(\sigma)-u(\sigma\tau)+u(\tau'),\qquad
  \katorep{\psi(\sigma)}h(\tau)\katoeq n\varepsilon(\sigma)u(\tau)-u(\sigma\tau)+u(\sigma),
\]
and $h(\tau)\katoeq n0$. Subtracting, and using $\varepsilon\katoeq n1$ on $D_B$ (\cref{pf:kato:epsdec}),
$h(\sigma)\katoeq nu(\tau')-u(\tau)\katoeq n0$.
\end{proof}

\begin{proposition}[Local height at a Frobenius]\label{pf:kato:locw}
Let $w$ be a prime with $\ell^n\mid w-1$, $B$ a place over $w$ and $\tau\in D_B$ a Frobenius at $B$ for $w$ with
$\det W(\tau)=1$; assume $W$ trivial on $I_B$. Let $\psi\colon G_\Q\to\Z/\ell^n$ be a locally constant homomorphism with
$\psi(t_0)=1$ for some $t_0\in I_B$, and $P\in\katoO$. Let $y\in\katoZ(G_\Q,W)$ and $z\in\katoZ(G_\Q,W^*)$ take values in
$\katoO^2$ and vanish on $I_B$; let $u\colon G_\Q\to\katoO$ be continuous with $u(sa)\katoeq nu(s)$ for $a\in I_B$ and
\begin{equation}\label{pf:kato:eq:du}
  (y\cup z)(s,t)\katoeq n\varepsilon(s)u(t)-u(st)+u(s)\qquad(s,t\in G_\Q);
\end{equation}
let $d\colon G_\Q\to\katoO^2$ be continuous with
\begin{equation}\label{pf:kato:eq:dertw}
  d(st)\katoeq nd(s)+W(s)d(t)+\katorep{\psi(s)}P\,W(s)y(t)\qquad(s,t\in G_\Q),
\end{equation}
and assume that $\xi(s,t)=(d\cup z)(s,t)-\katorep{\psi(s)}P\varepsilon(s)u(t)$ is $(\katoO,n)$-trivial on $D_B$. Then
$P\,\katoht_{y,z,u}(\tau)\katoeq n0$.
\formallabel{GlobalTheory.T.local_height_frobenius_place}
\end{proposition}

\begin{proof}
\emph{Preliminaries.} $D_B$ is the closure of $\{\tau^ki:k\ge0,\ i\in I_B\}$, because $D_B/I_B$ is procyclic,
topologically generated by the image of $\tau$, and the non-negative powers of a topological generator of a profinite
group are dense. As $\det W$ is continuous, trivial on $I_B$ and $1$ at $\tau$, it is $1$ on $D_B$; hence $W^*=\varepsilon W$ on
$D_B$, and $W^*=1$ on $I_B$ since $\varepsilon=1$ there. By \cref{pf:kato:epsdec}, $\varepsilon\katoeq n1$ on $D_B$. Taking $s=t=1$ in
\eqref{pf:kato:eq:du} gives $u(1)\katoeq n0$ (as $y(1)=0$), so $u\katoeq n0$ on $I_B$.

\emph{Step 1.} Let $\Delta=\det(W(\tau)-1)$ and $b=\katoadj(W(\tau)-1)z(\tau)\in\katoO^2$. We claim
\begin{equation}\label{pf:kato:eq:step1}
  \Delta\,z(s)\katoeq nW^*(s)b-b\qquad(s\in D_B).
\end{equation}
The map $X(s)=\Delta z(s)-(W^*(s)b-b)$ is a continuous cocycle of $W^*$ on $D_B$, vanishing on $I_B$. Since
$W(\tau)\katoadj(W(\tau)-1)=\Delta+\katoadj(W(\tau)-1)$, we have $W(\tau)b=\Delta z(\tau)+b$, hence
$X(\tau)=\Delta z(\tau)+b-\varepsilon(\tau)W(\tau)b=(1-\varepsilon(\tau))W(\tau)b\katoeq n0$. By the cocycle relation and induction,
$X(\tau^ki)\katoeq n0$ for $k\ge0$, $i\in I_B$, and by continuity $X\katoeq n0$ on $D_B$.

\emph{Step 2.} Put $h(t)=P\bigl(\Delta u(t)+\omega(y(t),W^*(t)b)\bigr)$. For $s,t\in D_B$, using $y(st)=y(s)+W(s)y(t)$,
\eqref{pf:kato:eq:omega}, \eqref{pf:kato:eq:step1} and \eqref{pf:kato:eq:du},
\[
\begin{aligned}
&\Delta u(st)+\omega\bigl(y(st),W^*(st)b\bigr)\\
&\quad=\Delta u(st)+\omega\bigl(y(s),W^*(s)(W^*(t)b-b)\bigr)+\omega\bigl(y(s),W^*(s)b\bigr)
  +\varepsilon(s)\,\omega\bigl(y(t),W^*(t)b\bigr)\\
&\quad\katoeq n\Delta u(st)+\Delta(y\cup z)(s,t)+\omega\bigl(y(s),W^*(s)b\bigr)+\varepsilon(s)\,\omega\bigl(y(t),W^*(t)b\bigr)\\
&\quad\katoeq n\bigl(\Delta u(s)+\omega(y(s),W^*(s)b)\bigr)+\varepsilon(s)\bigl(\Delta u(t)+\omega(y(t),W^*(t)b)\bigr),
\end{aligned}
\]
so $h(st)\katoeq nh(s)+h(t)$ (as $\varepsilon(s)\katoeq n1$). For $a\in I_B$, $y(a)=0$ and $u(a)\katoeq n0$, so $h(a)\katoeq n0$.

\emph{Step 3.} Let $u'$ be a cochain with $\xi(s,t)\katoeq n\varepsilon(s)u'(t)-u'(st)+u'(s)$ on $D_B$, and
$v(s)=-\bigl(\Delta u'(s)+\omega(d(s),W^*(s)b)\bigr)$, continuous on $D_B$. For $s,t\in D_B$, expanding $d(st)$ by
\eqref{pf:kato:eq:dertw} and using \eqref{pf:kato:eq:omega} and \eqref{pf:kato:eq:step1},
\[
\begin{aligned}
  &\omega\bigl(d(st),W^*(st)b\bigr)-\varepsilon(s)\omega\bigl(d(t),W^*(t)b\bigr)-\omega\bigl(d(s),W^*(s)b\bigr)\\
  &\qquad\katoeq n\Delta(d\cup z)(s,t)+\katorep{\psi(s)}P\varepsilon(s)\,\omega\bigl(y(t),W^*(t)b\bigr),
\end{aligned}
\]
and therefore
\[
\begin{aligned}
  \varepsilon(s)v(t)-v(st)+v(s)&\katoeq n-\Delta\xi(s,t)+\Delta(d\cup z)(s,t)+\katorep{\psi(s)}P\varepsilon(s)\,\omega\bigl(y(t),W^*(t)b\bigr)\\
  &=\katorep{\psi(s)}\varepsilon(s)h(t)\katoeq n\katorep{\psi(s)}h(t).
\end{aligned}
\]
So $(s,t)\mapsto\katorep{\psi(s)}h(t)$ is $(\katoO,n)$-trivial on $D_B$.

\emph{Step 4.} By \cref{pf:kato:brtame}, $h(\tau)\katoeq n0$.

\emph{Step 5.} For a $2\times2$ matrix $M$ with $\det M=1$ and vectors $x,x'$,
$\omega\bigl(x,M\katoadj(M-1)x'\bigr)=-\omega\bigl(\katoadj(M-1)x,x'\bigr)$: indeed $M\katoadj(M-1)=\det(M-1)+\katoadj(M-1)$,
$\omega(x,Nx')+\omega(Nx,x')=\tr(N)\omega(x,x')$ for every $N$, $\tr\katoadj(N)=\tr N$, and $\det(M-1)+\tr(M-1)=\det M-1=0$.
With $M=W(\tau)$ and $W^*(\tau)=\varepsilon(\tau)W(\tau)$ this gives
\[
  h(\tau)=P\bigl(\Delta u(\tau)-\varepsilon(\tau)\,\omega(\katoadj(W(\tau)-1)y(\tau),z(\tau))\bigr)
  \katoeq nP\,\katoht_{y,z,u}(\tau),
\]
using $\varepsilon(\tau)\katoeq n1$. Hence $P\,\katoht_{y,z,u}(\tau)\katoeq n0$.
\end{proof}

% ---------------------------------------------------------------------------
\subsection{Global lemmas}\label{pf:kato:sec:global}

In this subsection $\ell$ is odd, $\katoO$ is a coefficient ring and $S$ a finite set of primes containing $\ell$.

\begin{lemma}[Gluing]\label{pf:kato:glue}
Assume (PT) of \cref{pf:kato:hyp:litcoh}. Let $\xi$ be an $(\katoO,n,S)$-$2$-cocycle such that for every $v\in S$ there is a
place $B_v$ over $v$ on whose decomposition group $\xi$ is $(\katoO,n)$-trivial. Then there is a continuous
$u\colon G_\Q\to\katoO$ with $u(sa)\katoeq nu(s)$ for $a\in N_S$ and $\xi(s,t)\katoeq n\varepsilon(s)u(t)-u(st)+u(s)$ for all
$s,t\in G_\Q$.
\formallabel{GlobalTheory.T.local_cochains_glue_global}
\end{lemma}

\begin{proof}
For $n=0$ take $u=0$. Let $n\ge1$ and $R=\katoO/\ell^n\katoO$, a finite $\Z/\ell^n$-module on which $G_S$ acts by multiplication
by $\varepsilon$ modulo $\ell^n$ ($\varepsilon$ is trivial on $N_S$ since $\ell\in S$). Then $G_S$ acts trivially on
$R^\vee(1)=\Hom(R,\mu_{\ell^n})$, and $\Sha^1(G_S,R^\vee(1))$ consists of continuous homomorphisms $f\colon G_S\to R^\vee(1)$
vanishing on the decomposition groups at all $v\in S$. Such an $f$, viewed on $G_\Q$, vanishes on $N_S$ and on $D_{B_v}$, hence (being a
homomorphism to an abelian group, invariant under conjugation) on every inertia group of $G_\Q$; its kernel is an open
subgroup containing all inertia groups, hence equal to $G_\Q$ by Minkowski's theorem that $\Q$ has no non-trivial unramified
extension (\lref{UpperTheory.V.CF.minkAbs}, \cite[III.2.18]{Neukirch1999}). So $\Sha^1(G_S,R^\vee(1))=0$ and, by
(PT), $\Sha^2(G_S,R)=0$. By the defining properties of an $(\katoO,n,S)$-$2$-cocycle, $\xi$ modulo $\ell^n$ descends to a
continuous $2$-cocycle of $G_S$ with values in $R$, whose restriction to the decomposition groups at all $v\in S$ is trivial;
hence it is the coboundary of a continuous cochain $G_S\to R$. Lifting this cochain to $\katoO$ through a section of
$\katoO\to R$ (which is locally constant, $R$ being discrete) and composing with $G_\Q\to G_S$ gives $u$.
\end{proof}

\begin{lemma}[Reciprocity]\label{pf:kato:rec}
Assume (PT). Let $\xi$ be an $(\katoO,n,S)$-$2$-cocycle and $v_0\in S$, such that for every $v\in S\setminus\{v_0\}$, $\xi$ is
$(\katoO,n)$-trivial on $D_{B_v}$ for some place $B_v$ over $v$. Then $\xi$ is $(\katoO,n)$-trivial on $D_B$ for every place $B$
over $v_0$.
\formallabel{GlobalTheory.T.integralTwoCocycle_reciprocity_trivial_at_last_place}
\end{lemma}

\begin{proof}
For $n=0$ there is nothing to prove. Let $n\ge1$. Choose a basis of $\katoO$ over $\Z_\ell$; the coordinates $\xi_i$ of $\xi$
modulo $\ell^n$ are continuous $2$-cocycles of $G_S$ with values in $\Z/\ell^n(1)\cong\mu_{\ell^n}$. By (PT),
$\sum_{v\in S}\mathrm{inv}_v(\mathrm{res}_v\xi_i)=0$, where the restriction at $v$ is taken to the decomposition group of $B_v$
for $v\ne v_0$ and of $B$ for $v=v_0$. For $v\ne v_0$ the restriction is zero by hypothesis, hence
$\mathrm{inv}_{v_0}(\mathrm{res}_{v_0}\xi_i)=0$; as $\mathrm{inv}_{v_0}$ is injective, $\mathrm{res}_{v_0}\xi_i=0$ for all $i$, which
means that $\xi$ is $(\katoO,n)$-trivial on $D_B$.
\end{proof}

\begin{lemma}[Auxiliary primes]\label{pf:kato:auxprime}
Let $n\in\N$, $S,S'$ finite sets of primes, $U\subset G_\Q$ an open subgroup and $\sigma\in[\ker\varepsilon,\ker\varepsilon]$. There are a prime
$w\notin S\cup S'$, a place $B$ over $w$ and a Frobenius $\tau$ at $B$ for $w$ such that $\sigma^{-1}\tau\in U$,
$\ell^n\mid w-1$, and every $v\in S$ is an $\ell^n$-th power modulo $w$.
\formallabel{GlobalTheory.T.exists_auxiliary_prime_frobenius_near_commutator}
\end{lemma}

\begin{proof}
Let $\zeta\in\Omega$ be a primitive $\ell^n$-th root of unity and, for $v\in S$, $\alpha_v\in\Omega$ with $\alpha_v^{\ell^n}=v$. Every
$g\in\ker\varepsilon$ fixes $\zeta$ and maps $\alpha_v$ to $\zeta^{k}\alpha_v$ for some $k$; so $\ker\varepsilon$ acts on each $\alpha_v$ through
an abelian group, and $\sigma\in[\ker\varepsilon,\ker\varepsilon]$ fixes $\zeta$ and all $\alpha_v$. The set $O$ of $\tau$ with $\sigma^{-1}\tau\in U$
fixing $\zeta$ and all $\alpha_v$ is open and contains $\sigma$. By Chebotarev's density theorem
\cite{Chebotarev1926}, in the form that Frobenius elements at primes outside a finite set are dense in $G_\Q$
(\lref{GlobalTheory.Published.chebotarevDegreeOne'}), $O$ contains a Frobenius $\tau$ at a place $B$ over a prime
$w\notin S\cup S'\cup\{\ell\}$. As $\tau$ fixes $\zeta$ and acts as $x\mapsto x^w$ on the residue field, $\zeta^w=\zeta$, i.e.\
$\ell^n\mid w-1$. As $\tau$ fixes $\alpha_v$, $\alpha_v^w\equiv\alpha_v$ modulo the maximal ideal of $B$; since $v\ne w$, $\alpha_v$ is a unit
at $B$, so $\alpha_v^{w-1}\equiv1$ and the residue of $\alpha_v$ lies in $\F_w$; thus $v=\alpha_v^{\ell^n}$ is an $\ell^n$-th power
modulo $w$.
\end{proof}

\begin{lemma}[Characters of conductor $w$]\label{pf:kato:primechar}
Let $n\ge1$ and $w$ a prime with $\ell^n\mid w-1$. There is a surjective homomorphism $\psi\colon G_\Q\to\Z/\ell^n$ killing
$\mu_w$ such that: for every place $B$ over $w$ there is $\tau\in I_B$ with $\psi(\tau)=1$; and for every prime $v\ne w$ which is
an $\ell^n$-th power modulo $w$ and every place $B$ over $v$, $\psi$ is trivial on $D_B$.
\formallabel{GlobalTheory.T.exists_character_primeConductor_degree_prime_pow}
\end{lemma}

\begin{proof}
Let $f\colon(\Z/w)^\times\to\Z/\ell^n$ be a surjective homomorphism (it exists as $(\Z/w)^\times$ is cyclic of order divisible by
$\ell^n$) and $\psi=f\circ c_w$, where $c_w\colon G_\Q\to(\Z/w)^\times$ is the action on $\mu_w$. Then $\psi$ kills $\mu_w$. The
inertia group of a place over $w$ maps onto $(\Z/w)^\times$ ($\Q(\mu_w)/\Q$ is totally ramified at $w$), so some $\tau\in I_B$ has
$f(c_w(\tau))=1$; in particular $\psi$ is surjective. If $v\ne w$, the decomposition group of a place over $v$ acts on $\mu_w$
through the powers of $v$ in $(\Z/w)^\times$; if $v\equiv x^{\ell^n}\bmod w$, then $f(v)=\ell^nf(x)=0$, so $\psi$ is trivial there.
\end{proof}

\begin{lemma}[Kolyvagin derivative]\label{pf:kato:der}
Let $W$ be continuous with entries in $\katoO$, trivial on inertia outside $S$, and $y\colon G_\Q\to\katoO^2$. Let $w\notin S$ and
$\psi\colon G_\Q\to\Z/\ell^n$ a homomorphism killing $\mu_w$ with $\mathcal{T}(\katoO,W,y,n,w,\psi)$, with data $Y,P,u,b$. Put
\[
  d(s)=\sum_{x\in\Z/\ell^n}\katorep{x}\,Y(s,x)-\katorep{\psi(s)}\,W(s)b .
\]
Then $d$ is continuous with values in $\katoO^2$;
\eqref{pf:kato:eq:dertw} holds, i.e.\ $d(st)\katoeq nd(s)+W(s)d(t)+\katorep{\psi(s)}P\,W(s)y(t)$ for all $s,t$;
$d(a)\katoeq n0$ for $a\in N_{S\cup\{w\}}$; and for every $q\in S\setminus\{\ell\}$ there is a place $B$ over $q$ on whose inertia
group $d$ is an $(\katoO,n)$-coboundary. Moreover $P\katoeq nu\det(W(\sigma)-1)$ for every Frobenius $\sigma$ at a place over $w$.
\formallabel{GlobalTheory.T.unramified_cocycle_kills_roots}
\end{lemma}

\begin{proof}
The last assertion is \ref{pf:kato:T3}. For $x,a\in\Z/\ell^n$ we have $\katorep{x+a}=\katorep x+\katorep a-\ell^nc(x,a)$ with
$c(x,a)\in\{0,1\}$. By \ref{pf:kato:T1}, reindexing,
\[
\textstyle\sum_x\katorep xY(st,x)=\sum_x\katorep xY(s,x)+W(s)\sum_x\katorep{x+\psi(s)}Y(t,x),
\]
and by \ref{pf:kato:T2}, $\sum_x\katorep{x+\psi(s)}Y(t,x)=\sum_x\katorep xY(t,x)+\katorep{\psi(s)}\bigl(Py(t)+W(t)b-b\bigr)-\ell^nK$ with
$K\in\katoO^2$. Also $\katorep{\psi(st)}=\katorep{\psi(s)}+\katorep{\psi(t)}-\ell^nk'$ with $k'\in\{0,1\}$. Substituting into the definition
of $d(st)$, the terms involving $b$ cancel and
\[
d(st)=d(s)+W(s)d(t)+\katorep{\psi(s)}P\,W(s)y(t)+\ell^n\bigl(k'W(st)b-W(s)K\bigr),
\]
which gives \eqref{pf:kato:eq:dertw}. Continuity follows from that of $Y$ and $W$ and the local constancy of $\psi$ (its kernel
is open, as $\psi$ kills $\mu_w$). Let $\mathcal K$ be the set of $s$ with $W(s)=1$, $\psi(s)=0$ and $Y(s,x)=0$ for all $x$; by
\ref{pf:kato:T1} it is a subgroup, closed, and normal in $G_\Q$ (as in the proof of \cref{pf:kato:cup}(a)). For a prime
$q\notin S\cup\{w\}$ and $B$ given by \ref{pf:kato:T4}, $W$ is trivial on $I_B$, $\psi$ is trivial on $I_B$ ($I_B$ fixes $\mu_w$) and
$Y(s,x)=W(s)t(x)-t(x)=0$ on $I_B$; by conjugation $\mathcal K$ contains all inertia groups over such $q$, hence
$N_{S\cup\{w\}}\subset\mathcal K$, and $d$ vanishes on $\mathcal K$. Finally, for $q\in S\setminus\{\ell\}$ ($q\ne w$) and $B$, $t$ as in
\ref{pf:kato:T4}, $\psi=0$ on $I_B$ and $d(s)=W(s)x'-x'$ on $I_B$ with $x'=\sum_x\katorep xt(x)\in\katoO^2$.
\end{proof}

\begin{lemma}\label{pf:kato:xi}
Let $W$ be continuous with entries in $\katoO$ and $\delta^2=1$, and $S_2$ a finite set of primes with
$N_{S_2}\subset N_0(W)\cap\ker\psi$, where $\psi\colon G_\Q\to\Z/\ell^n$ is a locally constant homomorphism. Let $P\in\katoO$,
$y\in\katoZ(G_\Q,W)$, $z\in\katoZ(G_\Q,W^*)$ with values in $\katoO^2$, $u\colon G_\Q\to\katoO$ continuous and
$d\colon G_\Q\to\katoO^2$ continuous, such that $y$ and $z$ vanish and $d\katoeq n0$ on $N_{S_2}$, $u(sa)\katoeq nu(s)$ for
$a\in N_{S_2}$, and \eqref{pf:kato:eq:du} and \eqref{pf:kato:eq:dertw} hold. Then
\[
  \xi(s,t)=(d\cup z)(s,t)-\katorep{\psi(s)}P\varepsilon(s)u(t)
\]
is an $(\katoO,n,S_2)$-$2$-cocycle.
\formallabel{GlobalTheory.T.cocycle_pairing_locally_constant}
\end{lemma}

\begin{proof}
Continuity and integrality are clear. Put $D(r,s)=d(rs)-d(r)-W(r)d(s)\katoeq n\katorep{\psi(r)}PW(r)y(s)$. Using $z(st)=z(s)+W^*(s)z(t)$ and
\eqref{pf:kato:eq:omega}, a direct computation gives
\[
\begin{aligned}
  \varepsilon(r)(d\cup z)(s,t)+(d\cup z)(r,st)-(d\cup z)(rs,t)-(d\cup z)(r,s)&=-\omega\bigl(D(r,s),W^*(rs)z(t)\bigr)\\
  &\katoeq n-\katorep{\psi(r)}P\varepsilon(r)\,(y\cup z)(s,t).
\end{aligned}
\]
On the other hand, by $\katorep{\psi(rs)}\equiv\katorep{\psi(r)}+\katorep{\psi(s)}$ modulo $\ell^n$, the same alternating sum of the
terms $\katorep{\psi(\cdot)}P\varepsilon(\cdot)u(\cdot)$ is congruent to $-\katorep{\psi(r)}P\varepsilon(r)\bigl(\varepsilon(s)u(t)-u(st)+u(s)\bigr)$. By
\eqref{pf:kato:eq:du} the two expressions agree modulo $\ell^n$, so $\varepsilon(r)\xi(s,t)+\xi(r,st)\katoeq n\xi(rs,t)+\xi(r,s)$. For
$a,b\in N_{S_2}$: $W(sa)=W(s)$, $\varepsilon(sa)=\varepsilon(s)$, $\psi(sa)=\psi(s)$, $z(tb)=z(t)$, $u(tb)\katoeq nu(t)$, and by
\eqref{pf:kato:eq:dertw} $d(sa)\katoeq nd(s)+W(s)d(a)\katoeq nd(s)$; hence $\xi(sa,tb)\katoeq n\xi(s,t)$.
\end{proof}

% ---------------------------------------------------------------------------
\subsection{Heights on commutators and the strict dual Selmer group}\label{pf:kato:sec:heights}

The following theorem and corollary are the two parts of \cref{lem:window:heights}. Recall
$\katoht_{y,z,u}(s)=\det(W(s)-1)u(s)-\omega(\katoadj(W(s)-1)y(s),z(s))$; if $y\cup z$ is the coboundary of $u$ modulo
$\ell^n$, then $\katoht_{y,z,u}(s)$ is, up to the factor $\det(W(s)-1)$, the value at $s$ of a local pairing of $y$ and $z$
(Step 5 of \cref{pf:kato:locw}).

\begin{theorem}[\cref{lem:window:heights}, first part]\label{pf:kato:heights}
Assume the statements of \cref{pf:kato:hyp:litcoh} and $\mathrm{NF}(e,S,A,W)$. Let $\katoO$ be a coefficient ring containing the
entries of $W$, $n\ge1$, and $S'$ a finite set of primes. Let $y\in\katoZ(G_\Q,W)$ and $z\in\katoZ(G_\Q,W^*)$ take values in
$\katoO^2$, with $\mathrm{Unr}(W,y)$, and assume that $z$ is an integral coboundary on $I_B$ for every place $B$ over every prime
$q\ne\ell$ and on $D_A$ (with vectors in $\katoO^2$). Assume $\mathcal{T}(\katoO,W,y,n,w,\psi)$ for all primes $w\notin S'$ with
$\ell^n\mid w-1$ and all surjective $\psi\colon G_\Q\to\Z/\ell^n$ killing $\mu_w$. Let $u\colon G_\Q\to\katoO$ be continuous, with
$u(sa)\katoeq nu(s)$ for $a\in N_S$, and $(y\cup z)(s,t)\katoeq n\varepsilon(s)u(t)-u(st)+u(s)$ for all $s,t$. Then
\[
  \det\bigl(W(\sigma)-1\bigr)\,\katoht_{y,z,u}(\sigma)\katoeq n0\qquad\text{for every }\sigma\in[H_0(W),H_0(W)].
\]
\formallabel{GlobalTheory.Ord2.ht_vanish_commutator}
\end{theorem}

\begin{proof}
Let $\sigma\in[H_0,H_0]$; then $\sigma\in H_0$, so $\det W(\sigma)=1$, and $\sigma\in[\ker\varepsilon,\ker\varepsilon]$. The function
$F(s)=\det(W(s)-1)\katoht_{y,z,u}(s)$ is continuous, so there is an open subgroup $U_1$ with $|F(\sigma a)-F(\sigma)|<|\ell^n|$ for
$a\in U_1$, and an open subgroup $U_2$ with $|\det W(\sigma a)-\det W(\sigma)|<|2|$ for $a\in U_2$.

\emph{Auxiliary prime.} By \cref{pf:kato:auxprime} (with $U=U_1\cap U_2$) there are a prime $w\notin S\cup S'$, a place $B$ over
$w$ and a Frobenius $\tau$ at $B$ for $w$ with $\sigma^{-1}\tau\in U_1\cap U_2$, $\ell^n\mid w-1$, and every $v\in S$ an $\ell^n$-th power
modulo $w$. As $\delta^2=1$, $\det W(\tau)=\pm1$; if it were $-1$, then $|{-1}-1|<|2|$, absurd; so $\det W(\tau)=1$.

\emph{Character and derivative.} Let $\psi$ be given by \cref{pf:kato:primechar}; it is surjective, kills $\mu_w$, takes the value $1$
on an element of $I_B$, and is trivial on the decomposition groups of all places over primes $v\in S$ (these are $\ne w$ and are
$\ell^n$-th powers modulo $w$). As $w\notin S'$, $\mathcal{T}(\katoO,W,y,n,w,\psi)$ holds, and \cref{pf:kato:der} gives $P$, a unit $u_0$
and $d$ with \eqref{pf:kato:eq:dertw}, $P\katoeq nu_0\det(W(\tau)-1)$, $d\katoeq n0$ on $N_{S\cup\{w\}}$, and $d$ an $(\katoO,n)$-coboundary
on the inertia group of some place over each $q\in S\setminus\{\ell\}$.

\emph{Vanishing on $N_S$.} The inertia groups over primes $v\notin S$ lie in $N_0=\ker\varepsilon\cap\ker W$, which is closed and normal;
so $N_S\subset N_0$. The cocycles $y$ and $z$ vanish on these inertia groups (they are coboundaries there and $W=W^*=1$ there),
hence on $N_S$ (\cref{pf:kato:cup}(a)). Put $S_2=S\cup\{w\}$. Then $N_{S_2}\subset N_S\subset N_0$, and $N_{S_2}\subset\ker\psi$ because
$\psi$ is trivial on inertia at primes $\ne w$.

\emph{The $2$-cocycle.} By \cref{pf:kato:xi},
$\xi(s,t)=(d\cup z)(s,t)-\katorep{\psi(s)}P\varepsilon(s)u(t)$ is an $(\katoO,n,S_2)$-$2$-cocycle. Let $v\in S$ and $B'$ a place over $v$.
On $D_{B'}$, $\psi=0$, so $\xi=d\cup z$ there, and $d$ is an $(\katoO,n)$-cocycle of $W$ on $D_{B'}$ by \eqref{pf:kato:eq:dertw}. If
$v=\ell$, take $B'=A$: $z$ is an integral coboundary on $D_A$, so $\xi$ is $(\katoO,n)$-trivial on $D_A$ by \cref{pf:kato:cup}(b). If
$v\ne\ell$, take for $B'$ the place given by \cref{pf:kato:der}: $d$ and $z$ are $(\katoO,n)$-coboundaries on $I_{B'}$, and
\cref{pf:kato:unrcup} shows that $\xi$ is $(\katoO,n)$-trivial on $D_{B'}$. By \cref{pf:kato:rec} (for $S_2$ and $v_0=w$), $\xi$ is
$(\katoO,n)$-trivial on $D_B$.

\emph{Local computation.} As $w\notin S$, $W$ is trivial on $I_B$, and $I_B\subset N_S$, so $y,z$ vanish on $I_B$ and
$u(sa)\katoeq nu(s)$ for $a\in I_B$. All hypotheses of \cref{pf:kato:locw} hold, and it gives $P\,\katoht_{y,z,u}(\tau)\katoeq n0$. Since
$P\katoeq nu_0\det(W(\tau)-1)$ with $|u_0|=1$ and $|\katoht_{y,z,u}(\tau)|\le1$, we get
$\det(W(\tau)-1)\katoht_{y,z,u}(\tau)=u_0^{-1}\bigl(P\katoht(\tau)-(P-u_0\det(W(\tau)-1))\katoht(\tau)\bigr)\katoeq n0$, i.e.\ $F(\tau)\katoeq n0$.
Finally $\tau=\sigma(\sigma^{-1}\tau)$ with $\sigma^{-1}\tau\in U_1$, so $F(\sigma)\katoeq nF(\tau)\katoeq n0$.
\end{proof}

\begin{corollary}[\cref{lem:window:heights}, second part]\label{pf:kato:strict}
Assume the statements of \cref{pf:kato:hyp:litcoh}. If $\mathrm{NF}(e,S,A,W)$ and $\mathrm{ES}(W,A)$ hold, then
$H^1_{\mathrm{str}}(W^*,A)=0$: every cocycle of $W^*$ which is unramified away from $\ell$ and is a coboundary on $D_A$
is a coboundary.
\formallabel{GlobalTheory.Ord2.StrictKill.main}
\end{corollary}

\begin{proof}
Let $z\in\katoZ(G_\Q,W^*)$ satisfy $\mathrm{Unr}(W^*,z)$ and $\mathrm{Str}(W^*,A,z)$, and suppose that $z$ is not a coboundary. Let
$\katoO,y,S'$ be given by $\mathrm{ES}(W,A)$. By \cref{pf:kato:scaling}, applied first to $y$ and then to $z$, there are a coefficient
ring $\katoO_2\supset\katoO$ and $k_1,k_2$ such that $y'=\ell^{k_1}y$ and $z'=\ell^{k_2}z$ have values in $\katoO_2^2$, $y'$ is an integral
coboundary on all inertia groups at primes $\ne\ell$, and $z'$ is an integral coboundary on all inertia groups at primes $\ne\ell$ and on
$D_A$. If $Y,P,u,b$ witness $\mathcal{T}(\katoO,W,y,n,w,\psi)$, then $\ell^{k_1}Y,P,u,\ell^{k_1}b$ witness
$\mathcal{T}(\katoO_2,W,y',n,w,\psi)$, as one checks directly on \ref{pf:kato:T1}--\ref{pf:kato:T4}.

\emph{Gluing.} For every $n$, $y'\cup z'$ is an $(\katoO_2,n,S)$-$2$-cocycle (\cref{pf:kato:cup}(a)); it is $(\katoO_2,n)$-trivial on $D_A$
(\cref{pf:kato:cup}(b), as $z'$ is a coboundary there) and on the decomposition group of a place over each $v\in S\setminus\{\ell\}$
(\cref{pf:kato:unrcup}). By \cref{pf:kato:glue} there is a continuous $u_n\colon G_\Q\to\katoO_2$, with $u_n(sa)\katoeq nu_n(s)$ for
$a\in N_S$, such that $(y'\cup z')(s,t)\katoeq n\varepsilon(s)u_n(t)-u_n(st)+u_n(s)$.

\emph{Choice of $g,h,\phi$.} If $y'$ vanished on $N_0$, so would $y$, and $y$ would satisfy $\mathrm{Gr}(W,A,y)$ by \cref{pf:kato:grkill},
contradicting $\mathrm{ES}$. If $z'$ vanished on $N_0$, $z$ would be a coboundary by \cref{pf:kato:dualkill}. By \cref{pf:kato:pair}
there are $g,h\in N_0$ with $\delta'=\omega(y'(h),z'(g))-\omega(y'(g),z'(h))\ne0$, and by \cref{pf:kato:commdet} there is
$\phi\in[H_0,H_0]$ with $\Delta=\det(W(\phi)-1)\ne0$. As $N_0\subset H_0$, $\phi[g,h]\in[H_0,H_0]$, and $W(\phi[g,h])=W(\phi)$.

\emph{Conclusion.} Fix $n\ge1$ and write $\katoht=\katoht_{y',z',u_n}$. \Cref{pf:kato:heights} at $\phi$ and at $\phi[g,h]$ gives
$\Delta\katoht(\phi)\katoeq n0$ and $\Delta\katoht(\phi[g,h])\katoeq n0$, and \cref{pf:kato:htcomm} gives
$\katoht(\phi[g,h])-\katoht(\phi)\katoeq n\Delta\delta'$. Hence
\[
  \Delta^2\delta'=\Delta\bigl(\Delta\delta'-(\katoht(\phi[g,h])-\katoht(\phi))\bigr)+\Delta\katoht(\phi[g,h])-\Delta\katoht(\phi)\katoeq n0,
\]
using $|\Delta|\le1$. This holds for every $n$, so $\Delta^2\delta'=0$, a contradiction.
\end{proof}

\begin{remark}
In the classical theory \cite{Kolyvagin1990,Rubin2000} one bounds a dual Selmer group by pairing derivative classes with dual
classes at many auxiliary primes and invoking global duality. Here a single auxiliary prime is used at a time: the height is the
local pairing at a Frobenius element, the reciprocity law replaces global duality, and the commutator identity of
\cref{pf:kato:htcomm} turns the vanishing of heights on $[H_0,H_0]$ into a contradiction. Only the non-vanishing of the bottom class
on $N_0$ is used, which follows from its non-triviality in the Greenberg quotient (\cref{pf:kato:grkill}).
\end{remark}

% ---------------------------------------------------------------------------
\subsection{From the strict to the Bloch--Kato Selmer group}\label{pf:kato:sec:upgrade}

\begin{lemma}[Local duality modulo $\ell^n$]\label{pf:kato:ltd}
Assume (LT) of \cref{pf:kato:hyp:litcoh}. Let $A$ be a place over $\ell$, $\katoO$ a coefficient ring, $\theta,\theta'\colon
D_A\to\katoO^\times$ continuous characters with $\theta\theta'=\varepsilon$ on $D_A$, and $n\ge1$. Let $b\colon D_A\to\katoO$ be continuous
with $b(st)\katoeq nb(s)+\theta(s)b(t)$ for $s,t\in D_A$, and assume that for every continuous $a\colon D_A\to\katoO$ with
$a(st)\katoeq na(s)+\theta'(s)a(t)$ there is a continuous $u\colon D_A\to\katoO$ with
$b(s)\theta'(s)a(t)\katoeq n\varepsilon(s)u(t)-u(st)+u(s)$ for $s,t\in D_A$. Then there is $x\in\katoO$ with $b(s)\katoeq n\theta(s)x-x$
for all $s\in D_A$.
\formallabel{GlobalTheory.T.coboundary_of_orthogonal_dualCharacter_modPrimePow}
\end{lemma}

\begin{proof}
Let $M=\katoO/\ell^n\katoO$, a finite discrete $\Z/\ell^n$-module on which $D_A$ acts through $\theta'$, and let
$M^\vee(1)=\Hom(M,\Z/\ell^n)$, on which $s\in D_A$ acts by $(s\varphi)(m)=\varepsilon(s)\varphi(\theta'(s)^{-1}m)$; this is the Tate dual
of $M$. By (LT) for $D_A\cong G_{\Q_\ell}$, the cup product induces an injective map from $H^1(D_A,M^\vee(1))$ to the group of
homomorphisms $H^1(D_A,M)\to H^2(D_A,\Z/\ell^n(1))$. There is $\varphi_0\in\Hom(M,\Z/\ell^n)$ such that every element of
$\Hom(M,\Z/\ell^n)$ is $\varphi_0(x\,\cdot)$ for some $x\in\katoO$, and $\varphi_0(xM)=0$ only if $\ell^n\mid x$
(\lref{UpperTheory.V.X.So.dualGen}: this holds for the quotient of a discrete valuation ring finite over $\Z_\ell$ by any non-zero
element). As $\theta\theta'=\varepsilon$,
\begin{equation}\label{pf:kato:eq:key}
  \varepsilon(s)\,\varphi_0\bigl(x\theta'(s)^{-1}m\bigr)=\varphi_0\bigl(\theta(s)xm\bigr)\qquad(s\in D_A,\ x\in\katoO,\ m\in M).
\end{equation}
Put $\beta(s)=\varphi_0(b(s)\,\cdot)\in M^\vee(1)$. It is continuous, as it depends only on $b(s)$ modulo $\ell^n$, and by the congruence
for $b$ and \eqref{pf:kato:eq:key}, $\beta(st)=\beta(s)+s\beta(t)$. Let $y$ be a continuous cocycle of $D_A$ in $M$ and $a\colon
D_A\to\katoO$ the composite of $y$ with a set-theoretic section $M\to\katoO$; it is continuous ($M$ is discrete) and satisfies
$a(st)\katoeq na(s)+\theta'(s)a(t)$. Let $u$ be given by the hypothesis. Reducing modulo $\ell^n$ and applying $\varphi_0$, the
cup product $(\beta\cup y)(s,t)=\beta(s)\bigl(sy(t)\bigr)=\varphi_0\bigl(b(s)\theta'(s)y(t)\bigr)$ is the coboundary of the continuous
cochain $t\mapsto\varphi_0(u(t)\bmod\ell^n)$ with values in $\Z/\ell^n(1)$. Hence the class of $\beta$ pairs trivially with
$H^1(D_A,M)$, and $\beta$ is a coboundary: $\beta(s)=s\psi-\psi$ for some $\psi\in M^\vee(1)$. Write $\psi=\varphi_0(x\,\cdot)$ with
$x\in\katoO$. By \eqref{pf:kato:eq:key}, $\varphi_0\bigl((b(s)-\theta(s)x+x)m\bigr)=0$ for all $m\in M$, so
$\ell^n\mid b(s)-(\theta(s)x-x)$ for all $s\in D_A$.
\end{proof}

The formal statement of \cref{pf:kato:ltd} also assumes the finiteness statement (LF); the proof does not use it.

\begin{lemma}[Local duality for characters]\label{pf:kato:locdual}
Assume (LF), (LT), (LE) of \cref{pf:kato:hyp:litcoh}. Let $A$ be a place over $\ell$, $\katoO$ a coefficient ring, and
$k,k',a,b\colon G_\Q\to\katoQlb$ continuous on $D_A$ with values in $\katoO$ there, such that for $s,t\in D_A$:
$k(s)k'(s)=\varepsilon(s)$, $k$ and $k'$ are multiplicative, $b(st)=b(s)+k(s)b(t)$ and $a(st)=a(s)+k'(s)a(t)$. Assume $k\ne1$ and $k'\ne1$
on $D_A$, and that for every $n$ the map $(s,t)\mapsto b(s)k'(s)a(t)$ is $(\katoO,n)$-trivial on $D_A$. Then there is $x\in\katoQlb$ with
$b(s)=k(s)x-x$ for all $s\in D_A$, or there is $x$ with $a(s)=k'(s)x-x$ for all $s\in D_A$.
\formallabel{GlobalTheory.T.local_duality_character_orthogonal_classes}
\end{lemma}

\begin{proof}
Assume that $a$ is not a coboundary on $D_A$. As $kk'=\varepsilon$ is a unit and $k,k'$ are integral, $k$ and $k'$ are continuous characters
$\theta,\theta'\colon D_A\to\katoO^\times$ with $\theta\theta'=\varepsilon$; moreover $\theta'\ne1$, and $\theta'\ne\varepsilon$ because $\theta\ne1$. We use three
statements:
\begin{enumerate}[label=(\roman*)]
\item (\cref{pf:kato:depth}, with $H=D_A$ and $k=\theta'$.) There is $r$ such that for all $n$ and all $x,y$ with $|x|\le1$:
  if $xa(s)\katoeq n\theta'(s)y-y$ for all $s\in D_A$, then $|x|\le|\ell|^{\max(n-r,0)}$.
\item (\lref{GlobalTheory.T.exists_mul_scalarClass_eq_smul_add_coboundary_modPrimePow}, from (LF), (LT), (LE) and $\theta'\ne1,\varepsilon$,
  i.e.\ the local Euler characteristic formula for the rank-one module $\katoO(\theta')$.) For every $r$ there is an integer $m>0$ such that
  for every $n\ge1$: if $a$ satisfies the conclusion of (i) for this $r$, then for every continuous $\katoO$-valued $a'$ with
  $a'(st)\katoeq na'(s)+\theta'(s)a'(t)$ there are $x,y\in\katoO$ with $ma'(s)\katoeq nxa(s)+\theta'(s)y-y$ on $D_A$.
\item (\cref{pf:kato:ltd}, from (LT).) If $b'$ is a continuous
  $\katoO$-valued $(\katoO,n)$-cocycle of $\theta$ on $D_A$ such that $(s,t)\mapsto b'(s)\theta'(s)a'(t)$ is $(\katoO,n)$-trivial for every
  $(\katoO,n)$-cocycle $a'$ of $\theta'$, then $b'$ is an $(\katoO,n)$-coboundary (local Tate duality).
\end{enumerate}
Let $r$ be given by (i), $m$ by (ii), and $n\ge1$. For an $(\katoO,n)$-cocycle $a'$ of $\theta'$, (ii) gives $x,y$ with
$ma'\katoeq nxa+\theta'y-y$; if $u_n$ trivializes $(s,t)\mapsto b(s)k'(s)a(t)$ modulo $\ell^n$, then
$s\mapsto xu_n(s)-b(s)\theta'(s)y$ trivializes $(s,t)\mapsto b(s)\theta'(s)\,ma'(t)$ modulo $\ell^n$, by the cocycle relation for $b$.
By (iii), $mb$ is an $(\katoO,n)$-coboundary: $mb(s)\katoeq n\theta(s)x_n-x_n$. Hence $|b(s)-(k(s)X_n-X_n)|\le|\ell^n|/|m|$ with
$X_n=x_n/m$, for all $s\in D_A$. Choose $s_0\in D_A$ with $k(s_0)\ne1$; then $X_n\to X=b(s_0)/(k(s_0)-1)$, and $b(s)=k(s)X-X$ for all
$s\in D_A$, by passing to the limit.
\end{proof}

\begin{lemma}\label{pf:kato:nontriv}
If $\mathrm{NF}(e,S,A,W)$ holds, then $W(s)_{11}\ne1$ and $W^*(s')_{00}\ne1$ for some $s,s'\in D_A$.
\formallabel{GlobalTheory.Ord2.Upgrade.nontriv}
\end{lemma}

\begin{proof}
Let $t\in I_A$ with $\varepsilon(t)=1+\ell$, $m$ as in $\mathrm{NF}$ and $T=t^{2m}$. As in the proof of \cref{pf:kato:grkill},
$W(T)_{00}=q_0=(1+\ell)^{2em}>1$ and $\det W(T)=1$, so $W(T)_{11}=q_0^{-1}\ne1$ and
$W^*(T)_{00}=\varepsilon(T)\delta(T)W(T)_{00}=(1+\ell)^{2m}q_0\ne1$.
\end{proof}

\begin{theorem}[Strict to Greenberg]\label{pf:kato:upgrade}
Assume the statements of \cref{pf:kato:hyp:litcoh}, $\mathrm{NF}(e,S,A,W)$ and $\mathrm{ES}(W,A)$. If
$H^1_{\mathrm{str}}(W^*,A)=0$, then $H^1_{\mathrm{Gr}}(W^*,A)=0$.
\formallabel{GlobalTheory.Ord2.Upgrade.upgrade}
\end{theorem}

\begin{proof}
Let $z\in\katoZ(G_\Q,W^*)$ with $\mathrm{Unr}(W^*,z)$ and $\mathrm{Gr}(W^*,A,z)$, and let $\katoO,y$ be given by $\mathrm{ES}(W,A)$. By
\cref{pf:kato:scaling} there are a coefficient ring $\katoO_2$ and $j,k$ such that $y_1=\ell^jy$ and $\ell^kz$ have values in $\katoO_2^2$,
are integral coboundaries on all inertia groups at primes $\ne\ell$, and $\ell^kz(s)_1=W^*(s)_{11}x-x$ on $D_A$ for some $x\in\katoO_2$. Put
$z_1=\ell^kz-(W^*(\cdot)v_0-v_0)$ with $v_0=(0,x)$. Since $W^*$ is upper triangular on $D_A$, $z_1(s)_1=0$ for $s\in D_A$, and $z_1$ is
still an integral coboundary on all inertia groups at primes $\ne\ell$.

For every $n$, $y_1\cup z_1$ is an $(\katoO_2,n,S)$-$2$-cocycle (\cref{pf:kato:cup}(a)), $(\katoO_2,n)$-trivial on the decomposition groups
over all $v\in S\setminus\{\ell\}$ (\cref{pf:kato:unrcup}), hence on $D_A$ by \cref{pf:kato:rec}. For $s,t\in D_A$, as $z_1(t)_1=0$ and
$W^*(s)_{10}=0$,
\[
  (y_1\cup z_1)(s,t)=y_1(s)_0\bigl(W^*(s)z_1(t)\bigr)_1-y_1(s)_1\bigl(W^*(s)z_1(t)\bigr)_0=b(s)\,k'(s)\,a(t),
\]
with $b=-(y_1)_1$, $a=(z_1)_0$, $k=W_{11}$ and $k'=W^*_{00}$ on $D_A$. These satisfy the hypotheses of \cref{pf:kato:locdual}: $k$ and $k'$
are characters of $D_A$ with $kk'=W_{11}\varepsilon\delta W_{00}=\varepsilon\delta^2=\varepsilon$; $b(st)=b(s)+k(s)b(t)$ and $a(st)=a(s)+k'(s)a(t)$ by
the upper triangular shape; $k\ne1$ and $k'\ne1$ by \cref{pf:kato:nontriv}. Hence $b$ or $a$ is a coboundary on $D_A$. If $b$ is, then
$y(s)_1=W(s)_{11}x'-x'$ on $D_A$ for some $x'$, i.e.\ $\mathrm{Gr}(W,A,y)$, contradicting $\mathrm{ES}(W,A)$. So $a(s)=k'(s)x'-x'$ on $D_A$, and
then $z_1(s)=W^*(s)(x',0)-(x',0)$ on $D_A$: $z_1$ satisfies $\mathrm{Str}(W^*,A,z_1)$ and $\mathrm{Unr}$. As $H^1_{\mathrm{str}}(W^*,A)=0$, $z_1$ is
a coboundary, and so is $z=\ell^{-k}\bigl(z_1+W^*(\cdot)v_0-v_0\bigr)$.
\end{proof}

\begin{definition}[Selmer cocycles modulo $\ell^n$]\label{pf:kato:def:seln}
Assume $\mathrm{NF}(e,S,A,W)$. Let $\katoO$ be a coefficient ring containing the entries of $W$, $n\in\N$ and $V\in\{W,W^*\}$. A map
$z\colon G_\Q\to\katoO^2$ is a \emph{Selmer cocycle of $V$ modulo $\ell^n$} if it is an $(\katoO,n)$-cocycle of $V$ on $G_\Q$,
$z\katoeq n0$ on $N_S$, $z$ is an $(\katoO,n)$-coboundary on the inertia group of some place over each $q\in S\setminus\{\ell\}$, and
$z(s)_1\katoeq nV(s)_{11}x-x$ for all $s\in D_A$ and some $x\in\katoO$. Two of them are \emph{distinct} if their difference is not an
$(\katoO,n)$-coboundary on $G_\Q$.
\formal{GlobalTheory.Ord2.GwCount.SelN}
\end{definition}

For $n\ge1$ let $M_n=(\katoO/\ell^n)^2$ with $G_\Q$ acting through $W$. As $W$ is trivial on $N_S$, $M_n$ is a finite module of
$G_S=G_\Q/N_S$; as $\ell\in S$, so is $\mu_{\ell^n}$, and $M_n^\vee(1)=\Hom(M_n,\mu_{\ell^n})$. Let $\mathcal L=(\mathcal L_v)_{v\in S}$ be the
local conditions $\mathcal L_v=\ker\bigl(H^1(D_{B_v},M_n)\to H^1(I_{B_v},M_n)\bigr)$ for $v\ne\ell$ (places $B_v$ over $v$ being fixed), and
$\mathcal L_\ell$ the classes of cocycles $c$ with $c(s)_1=W(s)_{11}x-x$ on $D_A$ for some $x\in\katoO/\ell^n$ (the Greenberg condition).
Let $H^1_{\mathcal L}(G_S,M_n)$ be the classes whose restrictions lie in the $\mathcal L_v$, and $H^1_{\mathcal L^\perp}(G_S,M_n^\vee(1))$ the
dual Selmer group defined by the orthogonal complements under local duality.

\begin{theorem}[Greenberg--Wiles inequality]\label{pf:kato:gwineq}
Assume the statements of \cref{pf:kato:hyp:litcoh} and $\mathrm{NF}(e,S,A,W)$, and let $\katoO$ be a coefficient ring containing the
entries of $W$. There is $C\in\N$ such that for every $n\ge1$: if there are $N$ pairwise distinct Selmer cocycles of $W$ modulo $\ell^n$,
then there are $N'$ pairwise distinct Selmer cocycles of $W^*$ modulo $\ell^n$ with $N\le CN'$.
\formallabel{GlobalTheory.T.greenbergWiles_dualSelmer_lower_bound_mod_prime_pow}
\end{theorem}

\begin{proof}
We use four statements of the formal development:
(a) (\lref{GlobalTheory.T.card_galoisInvariants_ordinaryRankTwo_mod_pow_bounded}) there is $C_0$ with $\#H^0(G_S,M_n)\le C_0$
for all $n\ge1$;
(b) (\lref{GlobalTheory.T.card_ordinaryLocalConditionModPrimePow_le_const_mul}, from the local finiteness and Euler
characteristic statements (LF), (LE)) there is $C_A$ with $\#\mathcal L_\ell\le C_A\cdot\#(\katoO/\ell^n)$ for all $n\ge1$;
(c) (\lref{GlobalTheory.T.ordinarySelmerModPrimePow_injects_residualSelmer}) a Selmer cocycle of $W$ modulo $\ell^n$ defines a class in
$H^1_{\mathcal L}(G_S,M_n)$, and distinct Selmer cocycles define distinct classes;
(d) (\lref{GlobalTheory.T.dualSelmer_classes_give_noncohomologous_cochains}, from local duality) every class of
$H^1_{\mathcal L^\perp}(G_S,M_n^\vee(1))$ is represented by a Selmer cocycle of $W^*$ modulo $\ell^n$, distinct classes by distinct
Selmer cocycles.
We also use two standard facts. For $v\in S\setminus\{\ell\}$, $\#\mathcal L_v=\#H^0(D_{B_v},M_n)$: by inflation--restriction $\mathcal L_v\cong
H^1(D_{B_v}/I_{B_v},X)$ with $X=M_n^{I_{B_v}}$, and as $D_{B_v}/I_{B_v}\cong\widehat\Z$ is topologically generated by a Frobenius element $F$,
$H^1(D_{B_v}/I_{B_v},X)\cong X/(F-1)X$, which has the same order as $\ker(F-1)=H^0(D_{B_v},M_n)$ since $X$ is finite. Next,
$\#H^0(G_\R,M_n)=\#(\katoO/\ell^n)$: a complex conjugation $c$ acts by $W(c)$ with $W(c)^2=1$ and $\det W(c)=-1$; as $\ell$ is odd,
$(1\pm W(c))/2$ are complementary idempotents of $M_2(\katoO)$, each of rank one, so the invariants of $c$ on $M_n$ are free of rank one
over $\katoO/\ell^n$. Finally $H^1(G_\R,M_n)=0$ as $\ell$ is odd.

Let $C=C_0C_A$ and $n\ge1$. The Greenberg--Wiles formula (GW), with the conditions $\mathcal L_v$ and $\mathcal L_\infty=0$, reads
\begin{multline*}
  \#H^1_{\mathcal L}(G_S,M_n)\cdot\#H^0(G_S,M_n^\vee(1))\cdot\prod_{v\in S}\#H^0(D_{B_v},M_n)\cdot\#H^0(G_\R,M_n)\\
  =\#H^1_{\mathcal L^\perp}(G_S,M_n^\vee(1))\cdot\#H^0(G_S,M_n)\cdot\prod_{v\in S}\#\mathcal L_v .
\end{multline*}
By the first standard fact the factors for $v\in S\setminus\{\ell\}$ cancel. Using $\#H^0(G_S,M_n^\vee(1))\ge1$, $\#H^0(D_A,M_n)\ge1$, the
second standard fact, (a) and (b), we get
$\#H^1_{\mathcal L}(G_S,M_n)\cdot\#(\katoO/\ell^n)\le\#H^1_{\mathcal L^\perp}(G_S,M_n^\vee(1))\cdot C_0\cdot C_A\#(\katoO/\ell^n)$, i.e.
$\#H^1_{\mathcal L}(G_S,M_n)\le C\,\#H^1_{\mathcal L^\perp}(G_S,M_n^\vee(1))$. If $y_1,\dots,y_N$ are pairwise distinct Selmer cocycles of $W$
modulo $\ell^n$, (c) gives $N\le\#H^1_{\mathcal L}(G_S,M_n)$, and (d) gives $N'=\#H^1_{\mathcal L^\perp}(G_S,M_n^\vee(1))$ pairwise
distinct Selmer cocycles of $W^*$ modulo $\ell^n$; and $N\le CN'$.
\end{proof}

\begin{lemma}[Finiteness of cocycles modulo $\ell^n$]\label{pf:kato:gwfin}
Assume the statements of \cref{pf:kato:hyp:litcoh} and $\mathrm{NF}(e,S,A,W)$, and let $\katoO$ be a coefficient ring containing the
entries of $W$. For every $n$ there is $B_n$ such that every family of pairwise distinct $(\katoO,n)$-cocycles of $W^*$ on $G_\Q$ which are
$\katoeq n0$ on $N_S$ has at most $B_n$ members.
\formallabel{GlobalTheory.T.card_dualCocyclesModPrimePow_unramifiedOutsideS_bounded_alt}
\end{lemma}

\begin{proof}
$W^*$ is continuous, has entries in $\katoO$, and is trivial on $N_S$, since $W$ and $\varepsilon$ are trivial on the inertia groups at the
primes outside $S$. Let $M_n^*=(\katoO/\ell^n)^2$ with $G_S$ acting through $W^*$. The group $Z^1(G_S,M_n^*)$ of continuous cocycles is finite
(\lref{GlobalTheory.Ord2.GwFin.finite_cocycles_restricted_rep}, from the finiteness statement (GF)); put $B_n=\#Z^1(G_S,M_n^*)$. Let
$z$ be an $(\katoO,n)$-cocycle of $W^*$ with $z\katoeq n0$ on $N_S$. Its reduction $\bar z$ modulo $\ell^n$ is locally constant, satisfies the
cocycle relation, and $\bar z(sa)=\bar z(s)+W^*(s)\bar z(a)=\bar z(s)$ for $a\in N_S$; so $\bar z\in Z^1(G_S,M_n^*)$. If $\bar z_i=\bar z_j$, then
$z_i-z_j\katoeq n0$ is an $(\katoO,n)$-coboundary. Hence pairwise distinct cocycles have pairwise distinct reductions, and a family of them
has at most $B_n$ members.
\end{proof}

\begin{lemma}[Torsion bound]\label{pf:kato:torsbound}
Assume $\mathrm{NF}(e,S,A,W)$, let $\katoO$ be a coefficient ring containing the entries of $W$, and assume that for every $n$ there is
$B_n$ as in \cref{pf:kato:gwfin}. Then for every $m$ there is $B'_m$ such that for all $n$, every family of pairwise distinct Selmer
cocycles $z_i$ of $W^*$ modulo $\ell^n$ such that each $\ell^mz_i$ is an $(\katoO,n)$-coboundary on $G_\Q$ has at most $B'_m$ members.
\formallabel{GlobalTheory.T.card_primePowTorsion_selmerClassesModPrimePow_bounded}
\end{lemma}

\begin{proof}
Put $B'_m=\max\bigl(\max_{n<m}B_n,\ \#(\katoO/\ell^m)^2\cdot B_m\bigr)$, and let $z_1,\dots,z_N$ be as in the statement. Selmer cocycles modulo
$\ell^n$ are $(\katoO,n)$-cocycles which are $\katoeq n0$ on $N_S$; so if $n<m$, then $N\le B_n\le B'_m$. Let $n=m+k$, and choose $x_i\in\katoO^2$
with $\ell^mz_i(s)\katoeq nW^*(s)x_i-x_i$ for all $s$. Put
\[
  z'_i(s)=\ell^{-k}\bigl(z_i(s)-\ell^{-m}(W^*(s)x_i-x_i)\bigr).
\]
As $\ell^mz_i-(W^*x_i-x_i)\katoeq{m+k}0$, $z'_i$ has values in $\katoO^2$; it is continuous; its cocycle defect
$z'_i(st)-z'_i(s)-W^*(s)z'_i(t)$ is $\ell^{-k}$ times that of $z_i$, hence $\katoeq m0$; and $z'_i\katoeq m0$ on $N_S$, where $W^*=1$. Suppose
$N>\#(\katoO/\ell^m)^2\cdot B_m$. Then more than $B_m$ of the $x_i$ have the same reduction modulo $\ell^m$. For two such indices $a\ne b$, if
$z'_a-z'_b\katoeq mW^*y-y$ for some $y\in\katoO^2$, then, with $u=\ell^{-m}(x_a-x_b)\in\katoO^2$,
$z_a-z_b=\ell^k(z'_a-z'_b)+W^*u-u\katoeq{m+k}W^*(\ell^ky+u)-(\ell^ky+u)$, contradicting the distinctness of $z_a$ and $z_b$. So these $z'_i$
form a family of more than $B_m$ pairwise distinct $(\katoO,m)$-cocycles which are $\katoeq m0$ on $N_S$, contradicting the choice of $B_m$.
Hence $N\le\#(\katoO/\ell^m)^2\cdot B_m\le B'_m$.
\end{proof}

\begin{theorem}[Dual classes of unbounded order]\label{pf:kato:gwstep1}
Assume the statements of \cref{pf:kato:hyp:litcoh}, $\mathrm{NF}(e,S,A,W)$ and $H^1_{\mathrm{Gr}}(W,A)\ne0$. There is a coefficient ring
$\katoO$ containing the entries of $W$ such that for every $m\in\N$ there are $n\in\N$ and a Selmer cocycle $z$ of $W^*$ modulo $\ell^n$ for
which $\ell^mz$ is not an $(\katoO,n)$-coboundary on $G_\Q$.
\formallabel{GlobalTheory.T.exists_dualSelmer_classes_unbounded_order}
\end{theorem}

\begin{proof}
Let $y$ be a cocycle of $W$ with $\mathrm{Unr}$ and $\mathrm{Gr}$ which is not a coboundary. By \cref{pf:kato:scaling} there are $\katoO$ and $k$
such that $Y=\ell^ky$ is integral with integral local conditions; $Y$ is not a coboundary, and it vanishes on $N_S$, being a
homomorphism on $\ker W$ which vanishes on the inertia groups outside $S$
(\lref{GlobalTheory.T.cocycle_vanishes_on_ker_toGS_of_inertia_coboundary}). We use three statements:
(a) (\cref{pf:kato:boundedtors}) there is $k_0$ such that, for $a\in\katoO$, if $aY$ is an $(\katoO,n)$-coboundary then $a\ell^{k_0}\katoeq n0$;
(b) (\cref{pf:kato:gwineq}) there is $C$ such that for all $n\ge1$, if there are $N$ pairwise distinct Selmer cocycles of $W$ modulo $\ell^n$,
then there are $N'$ pairwise distinct Selmer cocycles of $W^*$ modulo $\ell^n$ with $N\le CN'$;
(c) (\cref{pf:kato:torsbound} with \cref{pf:kato:gwfin}) for every $m$ there is $B_m$ such that for all $n$, a
family of pairwise distinct Selmer cocycles $z_i$ of $W^*$ modulo $\ell^n$ such that every $\ell^mz_i$ is an $(\katoO,n)$-coboundary has at
most $B_m$ members.
Given $m$, choose $t$ with $\ell^t>CB_m$ and put $n=t+k_0$. The cocycles $jY$, $0\le j<\ell^t$, are Selmer cocycles of $W$ modulo $\ell^n$,
pairwise distinct: if $(i-j)Y$ were an $(\katoO,n)$-coboundary for some $0\le j<i<\ell^t$, (a) would give $|(i-j)\ell^{k_0}|\le|\ell^{t+k_0}|$, i.e.
$\ell^t\mid i-j$. By (b) there are $N'\ge\ell^t/C>B_m$ pairwise distinct Selmer cocycles of $W^*$ modulo $\ell^n$, so by (c) one of them,
$z$, has $\ell^mz$ not an $(\katoO,n)$-coboundary.
\end{proof}

\begin{theorem}[Passage to the limit]\label{pf:kato:gwstep2}
Assume the statements of \cref{pf:kato:hyp:litcoh} and $\mathrm{NF}(e,S,A,W)$, and let $\katoO$ be a coefficient ring containing the entries
of $W$. If for every $m$ there are $n$ and a Selmer cocycle $z$ of $W^*$ modulo $\ell^n$ such that $\ell^mz$ is not an $(\katoO,n)$-coboundary
on $G_\Q$, then $H^1_{\mathrm{Gr}}(W^*,A)\ne0$.
\formallabel{GlobalTheory.T.dualSelmer_unbounded_order_nonzero_ordinaryClass}
\end{theorem}

\begin{proof}
Let $X$ be the discrete Selmer group of $W^*$: the group of continuous cocycles of $W^*$ with values in $(\katoO\otimes\Q_\ell/\Z_\ell)^2$
which vanish on $N_S$, are coboundaries on the inertia groups at the places over $q\in S\setminus\{\ell\}$ and satisfy the Greenberg
condition at $A$, modulo coboundaries; it is an $\ell$-primary torsion group. Note that $\det(W^*(g)-1)\ne0$ for some $g$, by
\cref{pf:kato:commdet}, since $W^*=W$ on $H_0$. We use four statements of the formal development, applied to $W^*$:
(a) (\lref{GlobalTheory.T.finite_lTorsion_selmer_rank_two_adjoint}, from the finiteness statements of \cref{pf:kato:hyp:litcoh})
the $\ell$-torsion of $X$ is finite;
(b) (\lref{GlobalTheory.T.exists_cocycle_local_conditions}) if $\det(W^*(g)-1)\ne0$ for some $g$ and the hypothesis of the theorem
holds, then for every $m$ there is $x\in X$ with $\ell^mx\ne0$;
(c) (\lref{GlobalTheory.T.divisibleOrdinarySelmerChain_exists_integral_representatives}) if $x_0,x_1,\ldots\in X$ satisfy $x_0\ne0$ and
$\ell x_{k+1}=x_k$, there are $N\in\N$ and maps $a_k\colon G_\Q\to\katoO^2$ such that each $a_k$ is a locally constant $(\katoO,k)$-cocycle of
$W^*$ on $G_\Q$, $a_{k+1}\katoeq ka_k$, $a_k\katoeq k0$ on $N_S$, $a_k$ is congruent modulo $\ell^k$ to a coboundary on $I_B$ for every
place $B$ over a prime $q\in S\setminus\{\ell\}$ and satisfies the Greenberg condition at $A$ modulo $\ell^k$ (with vectors in the fraction
field of $\katoO$), and $a_N$ is not congruent modulo $\ell^N$ to a coboundary on $G_\Q$;
(d) (\lref{GlobalTheory.T.nonzero_ordinary_class_of_compatible_cocycles_mod_pow}) if $\det(W^*(g)-1)\ne0$ for some $g$, a family as in (c)
yields a continuous cocycle of $W^*$ with $\mathrm{Unr}$ and $\mathrm{Gr}$ which is not a coboundary (its limit), i.e.\
$H^1_{\mathrm{Gr}}(W^*,A)\ne0$.
By (a) and (b), $X$ satisfies the hypotheses of \cref{pf:kato:divchain}, which gives $x_0\ne0$ and $x_k=\ell x_{k+1}$ in $X$; (c) and (d)
conclude.
\end{proof}

\begin{theorem}[Greenberg--Wiles comparison]\label{pf:kato:gw}
Assume the statements of \cref{pf:kato:hyp:litcoh} and $\mathrm{NF}(e,S,A,W)$. Then $H^1_{\mathrm{Gr}}(W^*,A)=0$ implies $H^1_{\mathrm{Gr}}(W,A)=0$.
\formallabel{GlobalTheory.Ord2.grZero_of_grZero_dual1}
\end{theorem}

\begin{proof}
If $H^1_{\mathrm{Gr}}(W,A)\ne0$, \cref{pf:kato:gwstep1} gives a coefficient ring $\katoO$ and Selmer cocycles of $W^*$ of unbounded order,
and \cref{pf:kato:gwstep2} gives $H^1_{\mathrm{Gr}}(W^*,A)\ne0$.
\end{proof}

\begin{theorem}[Bloch--Kato classes are Greenberg classes]\label{pf:kato:bkgr}
Assume (BK) of \cref{pf:kato:hyp:litcoh}. If $\mathrm{NF}(e,S,A,W)$ holds and $z$ is a Bloch--Kato class of $W$, then $\mathrm{Unr}(W,z)$ and
$\mathrm{Gr}(W,A,z)$ hold.
\formallabel{GlobalTheory.T.blochKato_cocycle_unramified_and_greenbergOrdinary}
\end{theorem}

\begin{proof}
$\mathrm{Unr}(W,z)$ is condition (i) of \cref{pf:kato:def:bk}. By condition (ii) for the embedding $\Q\to\katoQlb$, and after conjugation
(\lref{GlobalTheory.T.crystallineClass_restrict_decompositionGroup}), there is a continuous homomorphism $\iota_A\colon
G_{\Q_\ell}=\Aut(\katoQlb/\Q_\ell)\to G_\Q$ with image $D_A$, mapping the inertia group of $G_{\Q_\ell}$ into $I_A$ and compatible with the cyclotomic
characters, such that $z\circ\iota_A$ is a crystalline cocycle of $W\circ\iota_A$. As $W$ is upper triangular on $D_A$, $\kappa=W_{11}\circ\iota_A$ is a
character and the projection to the second coordinate is equivariant, so $z_1\circ\iota_A$ is a crystalline cocycle of $\kappa$. On inertia,
$W_{00}^m=\varepsilon^{em}$ and $(W_{00}W_{11})^2=\delta^2=1$, so $(\kappa\varepsilon^e)^{2m}=1$: $\kappa$ agrees with $\varepsilon^{-e}$ on the kernel of the
continuous character $\kappa\varepsilon^e$ of finite order on inertia, an open subgroup. As $-e<0$, (BK) shows that $z_1\circ\iota_A$ is a
coboundary: $z_1(\iota_A(\sigma))=\kappa(\sigma)x-x$. Since $\iota_A$ maps onto $D_A$, $z(s)_1=W(s)_{11}x-x$ for $s\in D_A$.
\end{proof}

\begin{proof}[Proof of \cref{pf:kato:esvanish}]
By \cref{pf:kato:strict}, $H^1_{\mathrm{str}}(W^*,A)=0$; by \cref{pf:kato:upgrade}, $H^1_{\mathrm{Gr}}(W^*,A)=0$; by \cref{pf:kato:gw},
$H^1_{\mathrm{Gr}}(W,A)=0$. A Bloch--Kato class of $W$ satisfies $\mathrm{Unr}$ and $\mathrm{Gr}$ by \cref{pf:kato:bkgr}, hence is a coboundary.
\end{proof}

\begin{proof}[Proof of \cref{pf:kato:vanishordl}]
Let $\ell\ne2$, $e\ge1$, and $W$ admitted with $\mathrm{Src}(\iota,e,W)$. By \cref{pf:kato:nf} there are $S$, $A$, $W'\cong W$ and $g$ with
$\mathrm{NF}(e,S,A,W')$ and $\mathrm{Src}^{\mathrm{ur}}(\iota,e,A,W',g)$. By \textup{(Z)}, $\mathrm{ES}(W',A)$ holds; by \cref{pf:kato:esvanish},
$\katoHf(\Q,W')=0$; by \cref{pf:kato:conj}, $\katoHf(\Q,W)=0$.
\end{proof}

% ---------------------------------------------------------------------------
\subsection{Construction of the Euler system}\label{pf:kato:sec:zeta}

This subsection proves the existence statement \textup{(Z)} of \cref{pf:kato:def:zetaur}; the formal theorem has no
hypotheses. The construction follows Kato \cite{Kato2004}: zeta elements built from Siegel units \cite{KubertLang1981},
an explicit reciprocity law in the cyclotomic tower \cite{CoatesWiles1977,Coleman1979}, and the non-vanishing of an
$L$-value. We first isolate the interface between the two halves of the proof.

Throughout, $W$, $S$, $A$ and $g$ satisfy $\mathrm{NF}(e,S,A,W)$ and $\mathrm{Src}^{\mathrm{ur}}(\iota,e,A,W,g)$. For
$r\ge0$ we consider families $Z=(Z_r)_{r\ge0}$ of maps $Z_r\colon G_\Q\times(\Z/\ell^r)^\times\to\katoQlb$, and for a power
series $L=\sum_nL_nT^n$ over $\katoQlb$ we put
\[
  L_{<N}(z)=\sum_{n<N}L_n(z-1)^n,\qquad
  \mathrm{m}_r(L,\xi,x)=\sum_{0\le k<\ell^r}\xi^{\ell^r-(kx\bmod\ell^r)}L_{<\ell^{3r}}(\xi^k),
\]
for $z\in\katoQlb$, a primitive $\ell^r$-th root of unity $\xi$ and $x\in\N$; up to the factor $\ell^r$, $\mathrm{m}_r(L,\xi,x)$ is the
mass of $x+\ell^r\Z_\ell$ for the measure whose Amice transform is $L$. The \emph{$(\varphi,e)$-moments} of $L$ are the sums
$\sum_{0\le x<\ell^r,\,\ell\nmid x}\varphi(x)x^e\,\mathrm{m}_r(L,\xi,x)$, for a Dirichlet character $\varphi$ modulo $\ell$.

\begin{definition}[Euler system data]\label{pf:kato:def:esdata}
Let $\katoO$ be a coefficient ring, $y\colon G_\Q\to\katoQlb^2$, $S'$ a finite set of primes, $\nu\colon G_\Q\to\katoQlb^\times$,
$\chi=(\chi_r)_r$ with $\chi_r\colon(\Z/\ell^r)^\times\to\katoQlb$, and $Z$ as above.
\begin{enumerate}[label=(\alph*)]
\item $(\katoO,y,S')$ is an \emph{integral Euler system} for $W$, written $\mathrm{Sys}(W,\katoO,y,S')$, if $W$ and $y$ have
  entries in $\katoO$, $y\in\katoZ(G_\Q,W)$, $\mathrm{Unr}(W,y)$ holds, and $\mathcal T(\katoO,W,y,n,w,\psi)$ holds for all $n\ge1$,
  all primes $w\notin S'$ with $\ell^n\mid w-1$ and all surjective $\psi\colon G_\Q\to\Z/\ell^n$ killing $\mu_w$.
\item $y$ is \emph{non-zero at $A$}, written $\mathrm{Nz}(W,A,y)$, if $\varepsilon(s)=1$, $W(s)_{11}=1$ and $y(s)_1\ne0$ for some
  $s\in I_A$.
\item $(\nu,Z)$ is an \emph{Iwasawa family at $A$}, written $\mathrm{Iw}(\katoO,A,\nu,Z)$, if $\nu$ is multiplicative on $D_A$,
  each $Z_r$ has values in $\katoO$ and satisfies $Z_r(st,x)=Z_r(s,x)+\nu(s)Z_r(t,x\varepsilon_r(s)^{-1})$ for $s,t\in D_A$, and
  $Z_r(s,x)=\sum_{x'}Z_{r+1}(s,x')$, the sum over $x'\in(\Z/\ell^{r+1})^\times$ reducing to $x$.
\item $Z$ is \emph{linked} to $y$, written $\mathrm{Sp}(A,e,\chi,Z,y)$, if for every $n$ there is $r_0$ such that for $r\ge r_0$
  and $s\in I_A$ with $\varepsilon(s)=1$,
  $y(s)_1\katoeq n\sum_{x\in(\Z/\ell^r)^\times}\katorep{x}^{-e}\chi_r(x)Z_r(s,x)$.
\item Let $j\colon\Omega\to\katoQlb$ be an embedding, $m_0,c_0\in\N$, and $R=(R_m)$ with $R_m(\psi,z)\in\katoQlb$ for Dirichlet
  characters $\psi$ modulo $\ell^m$ and $z\in\katoQlb$. $Z$ satisfies the \emph{explicit reciprocity law}
  $\mathrm{HT}(A,\katoO,Z,j,m_0,c_0,R)$ if there is a sequence $(\zeta_m)$ of primitive $\ell^m$-th roots of unity in $\Omega$ with
  $\zeta_{m+1}^\ell=\zeta_m$ such that for every $m\ge m_0$ there are finitely many polynomials $P_i\in\Omega[X]$, whose
  coefficients lie in $A$ and are fixed by $I_A$, and $o_i\in\katoO$, with $P_i(\zeta_m)\ne0$, such that:
  \begin{itemize}
  \item (Kummer presentation) the $u_i=P_i(\zeta_m)$ are units of $A$, and for every $s\in I_A$ with $\varepsilon_m(s)=1$ there are
    $k_i\in\Z/\ell^{m-c_0}$ and $\ell^{m-c_0}$-th roots $u_i^{1/\ell^{m-c_0}}$ with
    $s(u_i^{1/\ell^{m-c_0}})=\zeta_m^{\ell^{c_0}k_i}u_i^{1/\ell^{m-c_0}}$ and $Z_m(s,1)\katoeq{m-c_0}\sum_io_i\katorep{k_i}$;
  \item (Coates--Wiles derivatives) for every primitive Dirichlet character $\psi$ modulo $\ell^m$,
    \[
      \sum_{b\in(\Z/\ell^m)^\times}\psi(b)\sum_io_i\,j\Bigl(\frac{\zeta_m^bP_i'(\zeta_m^b)}{P_i(\zeta_m^b)}\Bigr)\katoeq{m-c_0}R_m\bigl(\psi,j(\zeta_m)\bigr).
    \]
  \end{itemize}
\item $(W,\katoO,\nu,\chi,Z)$ is \emph{pinned}, written $\mathrm{Pin}(e,A,W,\katoO,\nu,\chi,Z)$, if there are $a,a'\in\katoQlb$ with
  $|a|=|a'|=1$, $\kappa\ne0$, units $v_m$ with $v_{m+1}\katoeq mv_m$, an embedding $j\colon\Omega\to\katoQlb$ with $|j(x)|\le1$ for
  $x\in A$, $R$ as in (e), a power series $L$ with coefficients in $\katoO$, a Dirichlet character $\varphi$ modulo $\ell$ and
  $m_0,c_0,n_0\in\N$ such that: $\nu=1$ on $I_A$ and $\nu(\sigma)=a$ for some Frobenius $\sigma\in D_A$ for $\ell$;
  $\mathrm{HT}(A,\katoO,Z,j,m_0,c_0,R)$ holds; for $m\ge m_0$, primitive $\psi$ modulo $\ell^m$ and primitive $\ell^m$-th roots of unity
  $z$,
  \[
    R_m(\psi,z)\katoeq{m-c_0}v_m\,\kappa\,a'^m\sum_{b\in(\Z/\ell^m)^\times}\psi(b)\,L_{<\ell^{3m}}(z^b);
  \]
  $\chi_r(x)=\varphi(x\bmod\ell)^{-1}$ for $r\ge1$; and there is $r_1$ such that for all $r\ge\max(r_1,1)$ and all primitive
  $\ell^r$-th roots of unity $\xi$, the $(\varphi,e)$-moment of $L$ is not $\katoeq{n_0+r}0$.
\end{enumerate}
\formal{GlobalTheory.KZ.Sys}
\formal{GlobalTheory.KZ.NzAt}
\formal{GlobalTheory.KZ.IwSys2}
\formal{GlobalTheory.KZ.SpLink}
\formal{GlobalTheory.KZ.HTLawCsOfCoef}
\formal{GlobalTheory.KZ.Pin7}
\end{definition}

The left-hand side of the law in (e) is a sum of Gauss sums of Coates--Wiles logarithmic derivatives of the
Coleman power series of the local units $u_i$ \cite{CoatesWiles1977,Coleman1979}; the pinning says that these are,
up to units, the Gauss sums of a fixed power series $L$, whose moments (an $\ell$-adic $L$-value) do not vanish.

\begin{definition}[The two halves]\label{pf:kato:def:halves}
\textup{(Z1)} is the statement: whenever $\mathrm{NF}(e,S,A,W)$ and $\mathrm{Src}^{\mathrm{ur}}(\iota,e,A,W,g)$ hold, there are
$\katoO,y,S',\nu,\chi,Z$ with $\mathrm{Sys}(W,\katoO,y,S')$, $\mathrm{Iw}(\katoO,A,\nu,Z)$, $\mathrm{Sp}(A,e,\chi,Z,y)$ and
$\mathrm{Pin}(e,A,W,\katoO,\nu,\chi,Z)$. \textup{(Z2)} is the statement: whenever $\mathrm{NF}(e,S,A,W)$ and
$\mathrm{Src}^{\mathrm{ur}}(\iota,e,A,W,g)$ hold, any such data satisfy $\mathrm{Nz}(W,A,y)$.
\formal{GlobalTheory.KZ.ExistsZetaEulerSystemIwasawaData}
\formal{GlobalTheory.KZ.ZetaEulerSystemNonvanishing}
\end{definition}

\begin{lemma}\label{pf:kato:halves}
\textup{(Z1)} and \textup{(Z2)} imply \textup{(Z)}.
\formallabel{GlobalTheory.KZ.zetaUr_of_halves2}
\end{lemma}

\begin{proof}
The formal lemma holds for every predicate in place of the pinning. Given $\mathrm{NF}(e,S,A,W)$ and
$\mathrm{Src}^{\mathrm{ur}}(\iota,e,A,W,g)$, \textup{(Z1)} gives $\katoO,y,S'$ with $\mathrm{Sys}(W,\katoO,y,S')$, and \textup{(Z2)} gives
$s\in I_A$ with $\varepsilon(s)=1$, $W(s)_{11}=1$ and $y(s)_1\ne0$. If $y$ satisfied the Greenberg condition,
$y(s)_1=W(s)_{11}x-x=0$, a contradiction. So $y$ is unramified away from $\ell$, does not satisfy the Greenberg condition at
$A$, and satisfies the tame relations: $\mathrm{ES}(W,A)$ holds.
\end{proof}

\begin{lemma}\label{pf:kato:quotchar}
If $\mathrm{NF}(e,S,A,W)$ and $\mathrm{Src}^{\mathrm{ur}}(\iota,e,A,W,g)$ hold, then $W(s)_{11}=1$ for every $s\in I_A$ with $\varepsilon(s)=1$.
\formallabel{GlobalTheory.T.quotient_char_trivial_inertia_cyclotomic_tower}
\end{lemma}

\begin{proof}
Let $s\in I_A$ with $\varepsilon(s)=1$; then $s$ fixes all $\ell$-power roots of unity. Let $\chi$ be the character of
$\mathrm{Src}^{\mathrm{ur}}$; as $s$ fixes $\mu_\ell$, $\chi(s)=1$. Next, $s$ fixes the square roots of every rational number. It
suffices to treat a nonzero integer $q=\pm\ell^aq'$ with $\ell\nmid q'$. The fields $\Q(\sqrt{-1})$ and $\Q(\sqrt{q'})$ are unramified
at the odd prime $\ell$, so their elements are fixed by the inertia element $s$ (recall that $A$ lies over $\ell$), and
$\sqrt{(-1)^{(\ell-1)/2}\ell}\in\Q(\mu_\ell)$ is fixed by $s$; the square roots of $q$ are products of these numbers and a rational number.
The character $\delta=\det W$ is continuous with values $\pm1$, so its kernel is a closed subgroup of index at most $2$ whose fixed
field is $\Q(\sqrt q)$ for some $q\in\Q$; as $s$ fixes this field, $s\in\ker\delta$, i.e.\ $\delta(s)=1$. Finally $\mathrm{Src}^{\mathrm{ur}}$ gives
$W(s)_{11}\varepsilon(s)^e\delta(s)^{-1}\chi(s)=1$, i.e.\ $W(s)_{11}=1$.
\end{proof}

\begin{littheorem}[Coleman power series; {\cite{Coleman1979}}, see also {\cite[17.12]{Kato2004}}]\label{pf:kato:lit:coleman}
Let $K$ be a finite unramified extension of $\Q_\ell$ with ring of integers $\OO_K$ and Frobenius $\phi$, and let $(\zeta_{\ell^m})_{m\ge1}$ be a
compatible system of primitive $\ell^m$-th roots of unity. For every sequence $u=(u_m)_{m\ge1}$ of units $u_m\in\OO_K[\zeta_{\ell^m}]^\times$
which is compatible under the norm maps, there is a unique power series $f_u\in\OO_K[[T]]^\times$ with
$(\phi^{-m}f_u)(\zeta_{\ell^m}-1)=u_m$ for all $m\ge1$, where $\phi$ acts on the coefficients.
\formallabel{UpperTheory.Kd.C7.pCore}
\end{littheorem}

\begin{remark}[The formal statement]\label{pf:kato:rem:coleman}
The formal statement is (ii) in the proof of \cref{pf:kato:plocu}: from a family $Z$ of cocycles on $I_A$ which is norm compatible,
Frobenius-equivariant up to a unit $a$, and admits Kummer presentations by polynomials $P_i$ with coefficients fixed by $I_A$, it
produces a power series $G$ whose Gauss sums agree, up to the smoothing factor $1-d_0\psi^{-1}(d_0)$, units $v_m$ and $a^m$, with the
Gauss sums of the logarithmic derivatives $\zeta P_i'(\zeta)/P_i(\zeta)$ at the points $\zeta_m^b$, and it transfers the vanishing
of the $(\varphi,k)$-specialisations of $Z$ to the vanishing of the $(\varphi,k)$-moments of $G$. Thus it combines
\cref{pf:kato:lit:coleman} with the logarithmic derivative homomorphisms of Coates and Wiles \cite{CoatesWiles1977} and with the
Kummer pairing on the cyclotomic tower. It differs from the published results in its formulation: the units are given by Kummer
presentations at a finite level $m$ rather than as a norm-compatible system, all statements are congruences modulo $\ell^{m-c_1}$
with explicit losses, and the conclusion is stated for Gauss sums and moments of $G$ instead of for the Coleman measure.
\end{remark}

\begin{theorem}[Transfer of non-vanishing]\label{pf:kato:transfer}
Let $\ell$ be odd, $\katoO$ a coefficient ring, $G$ and $L$ power series with coefficients in $\katoO$, $a,a',\kappa\in\katoQlb$ with
$|a|=|a'|=1$ and $\kappa\ne0$, $(v_m)$ and $(v'_m)$ sequences with $|v_m|=|v'_m|=1$, $v_{m+1}\katoeq mv_m$ and $v'_{m+1}\katoeq mv'_m$, $c_0\ge2$
prime to $\ell$, $m_0\ge2$, $c_1,n_0\in\N$, $\varphi$ a Dirichlet character modulo $\ell$ and $e\ge1$. Assume:
\begin{enumerate}[label=(\roman*)]
\item for every $m\ge m_0$ there is a primitive $\ell^m$-th root of unity $z$ such that for every primitive Dirichlet character $\psi$
  modulo $\ell^m$, with sums over $b\in(\Z/\ell^m)^\times$,
  \[
    \bigl(1-c_0\psi^{-1}(c_0)\bigr)v'_m\kappa a'^m\sum_b\psi(b)L_{<\ell^{3m}}(z^b)\katoeq{m-c_1}v_ma^m\sum_b\psi(b)G_{<\ell^{3m}}(z^b);
  \]
\item there is $r_1$ such that for all $r\ge\max(r_1,1)$ and all primitive $\ell^r$-th roots of unity $\xi$, the $(\varphi,e)$-moment of $L$
  at $(r,\xi)$ is not $\katoeq{n_0+r}0$.
\end{enumerate}
Then there are $n_1,r_2$ such that for all $r\ge r_2$ and all primitive $\ell^r$-th roots of unity $\xi$, the $(\varphi,e)$-moment of $G$ at
$(r,\xi)$ is not $\katoeq{n_1+r}0$.
\formallabel{GlobalTheory.T.nonvanishing_transfer_characterSums_powerSeries}
\end{theorem}

\begin{proof}
Write $\mathrm{mom}_r(F,\xi)$ for the $(\varphi,e)$-moment of a power series $F$, put $q=|\ell|$, $y_0=(1+\ell)^e$, and for a Dirichlet
character $\psi$ modulo $\ell^m$ and $z\in\katoQlb$ let $\tau_m(z,\psi)=\sum_{b\in(\Z/\ell^m)^\times}\psi(b)z^b$. We use four statements of the
formal development.
\begin{enumerate}[label=(\Alph*)]
\item (Mellin transform; \lref{GlobalTheory.KZ.Cmp.mellin_comp}.) Let $\katoO$ contain the $(\ell-1)$-th roots of unity, $F$ a power
  series over $\katoO$ and $d\in\N$. There is a power series $A$ over $\katoO$ such that: for $r\ge1$ and every primitive $\ell^r$-th root
  $\xi$, $d^e\varphi(d)\,\mathrm{mom}_r(F,\xi)\katoeq{2r}\ell^rA_{<r}(y_0)$; and for $k\ge1$, every primitive $\ell^{k+1}$-th root $z$ and
  every primitive $\psi$ modulo $\ell^{k+1}$ with $|\psi^{-1}(b)-b^e\varphi(b)|<1$ whenever $b^{\ell-1}=1$, putting $\zeta=\psi^{-1}(1+\ell)$,
  \[
    \psi^{-1}(d)\sum_b\psi(b)F_{<\ell^{3(k+1)}}(z^b)\katoeq{k+1}\tau_{k+1}(z,\psi)\,A_{<\ell^{3k}}(\zeta).
  \]
\item (Twisting characters; \lref{GlobalTheory.T.exists_primitive_character_twist}.) For $k\ge1$ and every primitive $\ell^k$-th root
  $\zeta$ there is a primitive $\psi$ modulo $\ell^{k+1}$ with $|\psi^{-1}(b)-b^e\varphi(b)|<1$ whenever $b^{\ell-1}=1$, and
  $\psi^{-1}(1+\ell)=\zeta$.
\item (Gauss sums; \lref{GlobalTheory.T.gaussSumModPrimePower_mul_inv_eq_prime_pow},
  \lref{GlobalTheory.T.norm_gaussSumModPrimePower_le_half_power}.) For a primitive $\ell^m$-th root $z$ and a primitive $\psi$ modulo
  $\ell^m$, $\tau_m(z,\psi)\tau_m(z^{-1},\psi^{-1})=\ell^m$ and $|\tau_m(z^{-1},\psi^{-1})|\le q^{\lfloor m/2\rfloor}$.
\item (Rigidity; \lref{GlobalTheory.T.powerSeries_eq_of_congruent_values_at_primitiveRoots}.) Let $A,B$ be power series over $\katoO$
  with $B\ne0$, $|u|=1$, $(w_m)$ a sequence with $|w_m|=1$ and $w_{m+1}\katoeq mw_m$, and $c,m'\in\N$. If
  $A_{<\ell^{3m}}(z)\katoeq{\lfloor m/2\rfloor-c}w_mu^mB_{<\ell^{3m}}(z)$ for all $m\ge m'$ and all primitive $\ell^m$-th roots $z$, then for
  every $k$ there is $m_1$ such that $A_n\katoeq kw_mB_n$ for all $m\ge m_1$ and all $n$.
\end{enumerate}

\emph{Step 1.} Choose $N$ with $|\ell^N\kappa|\le1$ and put $p=\ell^N$; enlarge $\katoO$ so that it contains $p\kappa$ and the $(\ell-1)$-th
roots of unity. Apply (A) to $G$ with $d=1$, giving $A_G$, and to $L$ with $d=1$ and $d=c_0$, giving $A_1$ and $A_2$; put
$A_L=A_1-c_0A_2$. Since $|c_0|\le1$, (A) gives for $r\ge1$, $k\ge1$ and $z,\psi,\zeta$ as there:
\begin{gather}
  \mathrm{mom}_r(G,\xi)\katoeq{2r}\ell^rA_{G,<r}(y_0),\qquad f\,\mathrm{mom}_r(L,\xi)\katoeq{2r}\ell^rA_{L,<r}(y_0),\label{pf:kato:eq:mom}\\
  \sum_b\psi(b)G_{<\ell^{3(k+1)}}(z^b)\katoeq{k+1}\tau\,A_{G,<\ell^{3k}}(\zeta),\label{pf:kato:eq:gs}\\
  \bigl(1-c_0\psi^{-1}(c_0)\bigr)\sum_b\psi(b)L_{<\ell^{3(k+1)}}(z^b)\katoeq{k+1}\tau\,A_{L,<\ell^{3k}}(\zeta),\label{pf:kato:eq:gs2}
\end{gather}
with $f=1-c_0^{e+1}\varphi(c_0)$ and $\tau=\tau_{k+1}(z,\psi)$. Here $f\ne0$: otherwise the integer $c_0^{e+1}\ge2$ would be a root of
unity.

\emph{Step 2: a lower bound.} By (ii) and \eqref{pf:kato:eq:mom}, for $r\ge\max(r_1,1)$ we have $|f\,\mathrm{mom}_r(L,\xi)|>|f|q^{n_0+r}$ and
$|f\,\mathrm{mom}_r(L,\xi)-\ell^rA_{L,<r}(y_0)|\le q^{2r}$; hence, for $r$ large, $|A_{L,<r}(y_0)|\ge|f|q^{n_0}$. So there are $\delta>0$ and $R$
with $|A_{L,<r}(y_0)|>\delta$ for $r\ge R$.

\emph{Step 3: comparison of values.} Put $u=a'/a$ and $w_k=(v'_{k+1}/v_{k+1})u$; then $|u|=1$, $|w_k|=1$ and $w_{k+1}\katoeq kw_k$. Let
$k\ge m_0+2c_1$ and let $\zeta$ be a primitive $\ell^k$-th root of unity. Take $z$ as in (i) for $m=k+1$, $\psi$ as in (B) for $\zeta$, and
$\tau=\tau_{k+1}(z,\psi)$, $\tau'=\tau_{k+1}(z^{-1},\psi^{-1})$. Substituting \eqref{pf:kato:eq:gs} and \eqref{pf:kato:eq:gs2} into (i) and multiplying by
$p\,\tau'/(v_{k+1}a^{k+1})$, (C) and $|p|,|p\kappa|\le1$ give
\[
  p\,A_{G,<\ell^{3k}}(\zeta)\katoeq{\lfloor k/2\rfloor-c_1}w_ku^k\cdot p\kappa\,A_{L,<\ell^{3k}}(\zeta):
\]
indeed $\ell^{k+1}=\tau\tau'$, and all error terms, of norm at most $q^{k+1-c_1}$, are multiplied by $|\tau'|/q^{k+1}\le q^{\lfloor(k+1)/2\rfloor-k-1}$, giving the precision $\lfloor(k+1)/2\rfloor-c_1\ge\lfloor k/2\rfloor-c_1$.

\emph{Step 4: rigidity.} The power series $pA_G$ and $p\kappa A_L$ have coefficients in $\katoO$, and $p\kappa A_L\ne0$ since its truncated
values at $y_0$ do not vanish by Step 2. By Step 3 and (D), with $c=c_1$ and $m'=m_0+2c_1$, for every $k$ there is $m_1$ with
$(pA_G)_n\katoeq kw_m(p\kappa A_L)_n$ for all $m\ge m_1$ and all $n$. As $|w_m|=1$ and $|y_0-1|\le1$, this gives
$|(p\kappa A_L)_{<M}(y_0)|\le\max\bigl(|(pA_G)_{<M}(y_0)|,q^k\bigr)$ for every $k$, hence $|p\kappa|\,|A_{L,<M}(y_0)|\le|p|\,|A_{G,<M}(y_0)|$ for
all $M$.

\emph{Step 5.} For $r\ge R$, Steps 2 and 4 and $|p|\le1$ give $|A_{G,<r}(y_0)|>\delta':=|p\kappa|\delta$. Choose $n_1$ with $q^{n_1}<\delta'$.
For $r\ge\max(n_1,R,1)$, \eqref{pf:kato:eq:mom} gives $|\mathrm{mom}_r(G,\xi)-\ell^rA_{G,<r}(y_0)|\le q^{2r}<q^r\delta'<|\ell^rA_{G,<r}(y_0)|$, so
$|\mathrm{mom}_r(G,\xi)|=|\ell^rA_{G,<r}(y_0)|>q^{r}\delta'>q^{n_1+r}$, i.e.\ $\mathrm{mom}_r(G,\xi)$ is not $\katoeq{n_1+r}0$.
\end{proof}

\begin{theorem}[Non-vanishing of the bottom class]\label{pf:kato:plocu}
\textup{(Z2)} holds.
\formallabel{UpperTheory.Kd.C7.pLocu}
\end{theorem}

\begin{proof}
We use \cref{pf:kato:quotchar}, the formal statement of \cref{pf:kato:lit:coleman} (\cref{pf:kato:rem:coleman}) and
\cref{pf:kato:transfer}.
\begin{enumerate}[label=(\roman*)]
\item (\cref{pf:kato:quotchar}.) Under $\mathrm{NF}(e,S,A,W)$ and $\mathrm{Src}^{\mathrm{ur}}(\iota,e,A,W,g)$, $W(s)_{11}=1$ for every
  $s\in I_A$ with $\varepsilon(s)=1$.
\item (Coleman power series; the formal statement of \cref{pf:kato:lit:coleman}.) Let $\ell$ be odd, $Z$ a family such that on $I_A$ the
  maps $Z_r$ are $\katoO$-valued, satisfy $Z_r(st,x)=Z_r(s,x)+Z_r(t,x\varepsilon_r(s)^{-1})$ and are norm compatible; assume that
  $Z_m(\sigma s\sigma^{-1},1)=aZ_m(s,1)$, with $|a|=1$, for every Frobenius $\sigma\in D_A$ for $\ell$ with $\varepsilon(\sigma)=1$ and
  every $s\in I_A$ with $\varepsilon_m(s)=1$; and that for $m\ge m_0$ the $Z_m$ admit Kummer presentations by polynomials as in
  \cref{pf:kato:def:esdata}(e) (with some $c_0$). Let $d_0$ be prime to $\ell$. Then there are a power series $G$ with
  coefficients in $\katoO$, units $v_m$ with $v_{m+1}\katoeq mv_m$, and $c_1,m_1$, such that: for $m\ge m_1$, every such
  presentation $(P_i,o_i)$ of $Z_m$ and every $\psi$ modulo $\ell^m$,
  \[
    \bigl(1-d_0\psi^{-1}(d_0)\bigr)\sum_b\psi(b)\sum_io_i\,j\Bigl(\frac{\zeta_m^bP_i'(\zeta_m^b)}{P_i(\zeta_m^b)}\Bigr)
    \katoeq{m-c_1}v_ma^m\sum_b\psi(b)\,G_{<\ell^{3m}}\bigl(j(\zeta_m)^b\bigr);
  \]
  and for every Dirichlet character $\varphi$ modulo $\ell$ and $k\ge1$: if for every $n$ the specialisations
  $\sum_x\varphi(x)^{-1}\katorep x^{-k}Z_r(s,x)$ are $\katoeq n0$ for all large $r$ and all $s\in I_A$ with $\varepsilon(s)=1$, then for
  every $n$ the $(\varphi,k)$-moments of $G$ are $\katoeq{n+r}0$ for all large $r$ and all primitive $\ell^r$-th roots $\xi$.
\item (Transfer of non-vanishing, \cref{pf:kato:transfer}.)
\end{enumerate}
Now let $\katoO,y,S',\nu,\chi,Z$ satisfy $\mathrm{Sys}$, $\mathrm{Iw}$, $\mathrm{Sp}$ and $\mathrm{Pin}$, with the data $a,a',\kappa,v',j,R,L,
\varphi,m_0,c_0,n_0$ of the pinning, and suppose that $\mathrm{Nz}(W,A,y)$ fails. By (i), $W(s)_{11}=1$ on $I_A\cap\ker\varepsilon$, hence
$y(s)_1=0$ for all $s\in I_A\cap\ker\varepsilon$. As $\chi_r(x)=\varphi(x)^{-1}$, the link $\mathrm{Sp}$ shows that the
$(\varphi,e)$-specialisations of $Z$ are $\katoeq n0$ on $I_A\cap\ker\varepsilon$ for all large $r$, for every $n$. Since $\nu=1$ on $I_A$,
the restriction of $Z$ to $I_A$ satisfies the hypotheses of (ii): the cocycle and norm relations follow from $\mathrm{Iw}$; for a
Frobenius $\sigma$ with $\varepsilon(\sigma)=1$ and $s\in I_A\cap\ker\varepsilon_m$, the cocycle relation on $D_A$ and $\nu(s)=1$ give
$Z_m(\sigma s\sigma^{-1},1)=\nu(\sigma)Z_m(s,1)$, and $\nu(\sigma)=\nu(\sigma_0)=a$ for the Frobenius $\sigma_0$ of the pinning, as
$\sigma_0^{-1}\sigma\in I_A$ and $\nu=1$ on $I_A$; and the Kummer presentations are those of the explicit law. Apply (ii) with
$d_0=2$: it gives $G$. For $m$ large, combining the congruence of (ii), the explicit law of (e) and the congruence of the
pinning (multiplied by $1-2\psi^{-1}(2)$, of norm $\le1$) gives hypothesis (i) of \cref{pf:kato:transfer} with $c_0=2$ and $z=j(\zeta_m)$. By (iii)
and the last condition of the pinning, the $(\varphi,e)$-moments of $G$ are not $\katoeq{n_1+r}0$ for large $r$; by the second
part of (ii), applied with $k=e$ and $n=n_1$, they are. This contradiction proves $\mathrm{Nz}(W,A,y)$.
\end{proof}

% ---------------------------------------------------------------------------
\subsection{Kato's zeta elements}\label{pf:kato:sec:katozeta}

It remains to prove \textup{(Z1)}. Its formal proof combines the results below. They are forms of theorems of Shimura, of
Jacquet--Shalika and, above all, of Kato \cite{Kato2004} (see \cite{Scholl1998,Colmez2004} for expositions, and \cite{Fukaya2003} for
the theory of Coleman power series for symbols). For each of them we quote the published theorem and describe in a remark the formal
statement and how it differs. The formal statements are theorems of the development without hypotheses, except
\cref{pf:kato:lit:proj}; they are proved from first principles, and the formal construction never uses the \'etale cohomology of modular
curves (\cref{pf:kato:rem:towermodule}).

\begin{littheorem}[{\cite[\S\S6.1--6.2, \S\S6.4--6.6]{Shimura1971}}]\label{pf:kato:lit:shimura}
For $N\ge1$ let $\mathcal F_N$ be the field of modular functions of level $N$ (meromorphic functions on the upper half-plane, invariant
under $\Gamma(N)$ and meromorphic at the cusps) whose $q$-expansions at $\infty$ have coefficients in $\Q(e^{2\pi i/N})$. Then
$\mathcal F_1=\Q(j)$, and $\mathcal F_N$ is a Galois extension of $\Q(j)$ with group $\GL_2(\Z/N)/\{\pm1\}$, in which $\SL_2(\Z/N)$ acts through
$f\mapsto f\circ\gamma$ ($\gamma\in\SL_2(\Z)$) and the matrices $\mathrm{diag}(1,d)$ act on the $q$-expansion coefficients through
$\Gal(\Q(e^{2\pi i/N})/\Q)$. These actions extend to an action of $\GL_2(\A_{\Q,f})$ on $\bigcup_N\mathcal F_N$ (Shimura's reciprocity law).
\formallabel{GlobalTheory.T.exists_modular_function_tower_algebraic_coefficients}
\end{littheorem}

\begin{remark}[The formal statement]\label{pf:kato:rem:shimura}
The formal statement asserts: for every algebraic closure $\Omega$ of $\Q$ and every embedding $\Omega\to\C$ there is a \emph{standard
modular tower}, that is, a field $K_\infty\supset\Omega$ of meromorphic functions on the upper half-plane, each invariant under a principal
congruence subgroup, with $q$-expansions in $\Omega((q^{\Q}))$, an action of $\GL_2^+(\Q)$ by composition and an action of $G_\Q$ on
$q$-expansion coefficients satisfying Shimura's reciprocity law, containing the modular invariant $j$, with bounded denominators and a
saturation property. It is a model of $\Omega\otimes\bigcup_N\mathcal F_N$ in which only the properties used below are retained; the
adelic action of \cref{pf:kato:lit:shimura} is replaced by the actions of $\GL_2^+(\Q)$ and of $G_\Q$, and the coefficients are
algebraic numbers instead of elements of $\Q(e^{2\pi i/N})$. The formal development constructs the tower.
\end{remark}

\begin{corollary}\label{pf:kato:towerex}
For every algebraic closure $\Omega$ of $\Q$ and every embedding $\Omega\to\C$ there is a modular tower.
\formallabel{UpperTheory.V.Z.Kato.Joint.towerEx}
\end{corollary}

\begin{proof}
Forget the standardness in the formal statement of \cref{pf:kato:lit:shimura}.
\end{proof}

\begin{remark}[The tower module]\label{pf:kato:rem:towermodule}
In a modular tower, the Kummer classes of functions which are units up to $\ell^n$-th powers at all places form a module $\mathcal A$;
homogeneous polynomial maps $\Z_\ell^2\to\mathcal A$ of degree $2e-1$ invariant under $\Gamma_1(N)\cap\Gamma_0(\ell)$ form a module $B_{N,2e-1}$
with Hecke operators $T_q$, $U_\ell$, $\langle d\rangle$ and an action of $G_\Q$. This module replaces the \'etale cohomology of modular curves
with coefficients in $\mathrm{Sym}^{2e-1}$ used in \cite{Kato2004}.
\end{remark}

\begin{littheorem}[{\cite[Thm.~(1.3)]{JacquetShalika1976}}]\label{pf:kato:lit:js}
Let $\pi$ be an irreducible cuspidal automorphic representation of $\GL_n(\A_F)$, $F$ a number field, with unitary central character.
Then $L(s,\pi)\ne0$ for $\Re s=1$.
\formallabel{GlobalTheory.T.odd_chi_and_tendsto_twisted_LSeries_ne_zero_of_prim}
\end{littheorem}

\begin{remark}[The formal statement]\label{pf:kato:rem:js}
The formal statement asserts: let $k$ be odd, $N,M\ge1$, $g=\sum_ng(n)q^n$ a primitive form of weight $k$, level $N$ and character $\chi$,
$\varphi$ a Dirichlet character modulo $M$, and $w\ge(k+1)/2$. Then $\chi(-1)=-1$, and $\sum_n\varphi(n)g(n)n^{-s}$, which converges absolutely
for $\Re s>(k+1)/2$, tends to a non-zero limit as $s\to w$ along real $s>w$. The parity statement is elementary: the action of $-1\in\SL_2(\Z)$
multiplies $g$ by $\chi(-1)$ and by $(-1)^k$. For $w>(k+1)/2$ the limit is the value of an absolutely convergent Euler product. For
$w=(k+1)/2$ the statement is \cref{pf:kato:lit:js} for the cuspidal representation attached to $g$ twisted by the primitive character
inducing $\varphi$ (with $s$ shifted by $(k-1)/2$), up to finitely many Euler factors, which do not vanish on $\Re s=(k+1)/2$ by the
Ramanujan bound \cite{Deligne1974}. The formal proof does not use \cite{JacquetShalika1976}: it proves the non-vanishing by a positivity
argument for Dirichlet series (\lref{GlobalTheory.T.lseries_ne_zero_re_ge_one_of_positivity}) together with Rankin's mean square
estimate for the coefficients of cusp forms (\lref{GlobalTheory.T.rankin_mean_square_cuspForm}).
\end{remark}

\begin{littheorem}[Zeta value formula; {\cite[Thm.~6.6 and \S7]{Kato2004}}, {\cite{Shimura1976}}]\label{pf:kato:lit:eis}
Let $f$ be a normalised eigenform of weight $k\ge2$ and level $N$. Kato's zeta modular forms, products of two Eisenstein series, have the
following property: for a Dirichlet character $\chi$ modulo $m$, the sum over $b\in(\Z/m)^\times$ of $\chi(b)$ times the period of the
projection to $f$ of the $\sigma_b$-conjugate of the zeta modular form of level $m$ equals $(2\pi i)^{k-r-1}$ times the partial $L$-value
$L(f^*,\chi,r)$ (with the Euler factors at a finite set of primes removed), times explicit elementary factors and a period of $f$. The proof is
by the Rankin--Selberg method.
\formallabel{GlobalTheory.T.ordinaryEisensteinEigenfunctional_exists}
\end{littheorem}

\begin{remark}[The formal statement]\label{pf:kato:rem:eis}
The formal statement provides the auxiliary data for the linear form $\lambda$ of \cref{pf:kato:lit:proj}. For $\ell$ odd, $e\ge1$, a
newform $g$ of weight $2e+1$ with level prime to $\ell$ and $a_\ell(g)$ an $\ell$-adic unit, a character $\psi$ modulo $\ell$ and $N_0\ge1$ prime
to $\ell$, there are a level $N\ge5$ divisible by $N_0$ and by the level of $g$ and prime to $\ell$, Dirichlet characters $\chi,\theta$
modulo $N$ with $\chi$ inducing the character of $g$, $\chi(-1)=-1$ and $\theta(-1)=(-1)^e\psi(-1)$, a unit root $\alpha$ of
$X^2-a_\ell(g)X+\chi(\ell)\ell^{2e}$, an ordinary Hecke eigen-system, a constant $v\ne0$ satisfying a Rankin--Selberg type identity between
an eigen-functional and a measure built from Eisenstein series, and smoothing integers $c_1,d_1>1$ prime to $6N\ell$ avoiding finitely
many degenerate values. It differs from \cref{pf:kato:lit:eis} in that periods of modular symbols are replaced by an ordinary
eigen-functional on $q$-expansions, the $U_\ell$-stabilisation by the unit root $\alpha$ is built in, and only the non-vanishing of the
relevant constant is asserted; this non-vanishing comes from the formal statement of \cref{pf:kato:lit:js}. The formal proof chooses an
auxiliary prime $r\ge\ell+7$, passes to the form with coefficients $a_n(g)$ for $n$ prime to $MN_0r$ ($M$ the level of $g$), of level
$M(MN_0r)^2$ (\lref{GlobalTheory.EisM.depleted_newform_prim}), and constructs the data at that level
(\lref{GlobalTheory.EisM.chapterM_of_level}).
\end{remark}

\begin{littheorem}[{\cite{Deligne1971}}, {\cite[6.3 and 8.3]{Kato2004}}]\label{pf:kato:lit:hecke}
Let $f$ be a newform of weight $k\ge2$, level $N$ and character $\chi$, with coefficient field $F$, and $\lambda$ a place of $F$ over a prime $p$.
The quotient $V_{F_\lambda}(f)$ of $H^1_{\mathrm{et}}(Y_1(N)\otimes\Qbar,\mathrm{Sym}^{k-2})\otimes F_\lambda$ corresponding to $f$, on which
the Hecke operators act through the eigenvalues of $f$, is two-dimensional over $F_\lambda$, and for every prime $q\nmid Np$ the arithmetic Frobenius at $q$ acts on it with
characteristic polynomial $X^2-a_q(f)X+\chi(q)q^{k-1}$, up to the normalisation of the action.
\formallabel{GlobalTheory.T.heckeProjector_newform_galoisImage}
\end{littheorem}

\begin{remark}[The formal statement]\label{pf:kato:rem:hecke}
The formal statement concerns the tower module of weight $w=k-2$ and level $N$ (\cref{pf:kato:rem:towermodule}) for a newform $g$ of odd
weight $k\ge3$, level dividing $N$, $5\le N$, $\ell\nmid N$, and a root $\alpha$ of $X^2-a_\ell(g)X+\chi(\ell)\ell^{k-1}$ with
$\alpha^2\ne\chi(\ell)\ell^{k-1}$ and $\alpha^2\ne\chi(\ell)\ell^{k-2}$; let $\lambda$ be the eigen-system with $T_q\mapsto a_q(g)$, $U_\ell\mapsto\alpha$,
$\langle d\rangle\mapsto\chi(d)$. Given a continuous irreducible $\rho\colon G_\Q\to\GL_2(E)$ such that Frobenius elements at almost all
primes $q$ have trace $a_q(g)$ and determinant $q^{w+1}\chi(q)$, it asserts that there are $m\ge0$, Galois- and Hecke-equivariant maps
$\pi_i$ from the tower module to $E^2$ and $\sigma_i$ back ($1\le i\le m$) with $\pi_i\sigma_j=\delta_{ij}$, and a polynomial $P$ in the admissible
Hecke operators with $P(\lambda)=1$ which acts on the tower module as $\sum_i\sigma_i\pi_i$, the identity on the $\lambda$-eigenspace. So the
$\lambda$-part is a direct summand isomorphic to $\rho^m$, cut out by $P$. This replaces \cref{pf:kato:lit:hecke}: the module is different, the
multiplicity $m$ is not determined, and the eigen-system is $\ell$-stabilised. The formal proof shows the semisimplicity of the Hecke action
at $\lambda$ on a space of period polynomials (\lref{GlobalTheory.T.semisimple_potentials_at_newform_eigensystem}), transfers it to the
tower module (\lref{GlobalTheory.T.semisimple_tensor_of_semisimple_potentials}), and identifies the eigenspace with copies of $\rho$
through the traces of Frobenius elements (\lref{GlobalTheory.T.galois_representation_of_eigensystem}).
\end{remark}

\begin{littheorem}[{\cite[6.3, 8.3 and Thm.~9.7]{Kato2004}}]\label{pf:kato:lit:proj}
With the notation of \cref{pf:kato:lit:hecke}, there is a Hecke- and Galois-equivariant projection from the cohomology of modular curves
onto $V_{F_\lambda}(f)$, and the dual exponential maps of the images of Kato's $p$-adic zeta elements are the projections to $f$ of the zeta
modular forms of \cref{pf:kato:lit:eis}.
\formallabel{UpperTheory.V.Z.Kato.Joint.kato_global_admissibility_of_admissible}
\end{littheorem}

\begin{remark}[The formal statement]\label{pf:kato:rem:proj}
The formal statement is an implication. Its hypothesis is the existence of admissible Hecke projections, proved in the formal
development (\lref{GlobalTheory.T.exists_admissible_hecke_projection_nonvanishing}): from data as in the formal statement of
\cref{pf:kato:lit:hecke} and an approximate Hecke eigen-functional on $q$-expansions it produces a projection. Its conclusion is: under
$\mathrm{NF}(e,S,A,W)$ and $\mathrm{Src}^{\mathrm{ur}}(\iota,e,A,W,g)$, for an embedding $j\colon\Omega\to\katoQlb$ defining $A$, a modular
tower, and data $(N,\chi,\alpha)$ with an eigen-functional as in the formal statement of \cref{pf:kato:lit:eis}, there are a
$\Z_\ell$-linear map $\mathrm{pr}$ from the tower to $\katoQlb^2$, integral on $B_{N,2e-1}$, equivariant for $G_\Q$ acting through a twist of
$W$ by a power of $\varepsilon$ and a character of $(\Z/\ell)^\times$, and for the Hecke operators $T_q$ ($q\nmid\ell N$) acting through the
eigenvalues of $g$, $U_\ell$ through $\alpha$ and the diamond operators through a quadratic character; and a linear form $\lambda$ on
$q$-expansion sequences recovering the second coordinate of $\mathrm{pr}$ up to convergent unit periods, which on $q$-expansions of
modular forms of weight $2e+1$, level $N\ell$ and character $\chi$ equals $\kappa\ne0$ times the eigen-functional. The formal proof
combines the formal form of \cref{pf:kato:lit:hecke}, the ordinary triangular shape of the twist of $W$
(\lref{GlobalTheory.T.untwisted_representation_ordinary_triangular_shape}) and an eigen-functional built from the Hecke
polynomial $P$. The conditions on $\mathrm{pr}$ do not require it to be non-zero.
\end{remark}

\begin{littheorem}[{\cite[Prop.~8.12, Example~13.3, Thm.~12.5(4)]{Kato2004}}]\label{pf:kato:lit:zeta}
Let $f$ be a newform of weight $k\ge2$, $\lambda$ a place of its coefficient field over an odd prime $p$, $r\in\Z$, and $T=V_{\OO_\lambda}(f)(k-r)$.
For suitable auxiliary integers $c,d$ and $1\le j\le k-1$, Kato's zeta elements $z_m\in H^1(\Z[\zeta_m,1/p],T)$, for $m$ in the set of
\cite[13.3]{Kato2004}, form an Euler system for $T$: for $m\mid m'$ the norm map sends $z_{m'}$ to $\prod_q P_q(q^{-1}\sigma_q^{-1})z_m$, where
$q$ runs over the primes dividing $m'$ but not $m$, $P_q(t)=\det(1-\mathrm{Fr}_q^{-1}t\mid T)$ and $\sigma_q$ is the arithmetic Frobenius of $q$ in
$\Gal(\Q(\zeta_m)/\Q)$. For an arbitrary Galois-stable lattice $T'$ of $V_{F_\lambda}(f)(k-r)$ the analogous elements lie in $H^1(T')$ if the image
of $G_{\Q(\zeta_{p^\infty})}$ contains $\SL_2(\Z_p)$.
\formallabel{UpperTheory.Kd.K4.kato_eulerSystem_integrality}
\end{littheorem}

\begin{remark}[The formal statement]\label{pf:kato:rem:zeta}
The formal statement asserts: under $\mathrm{NF}(e,S,A,W)$ and $\mathrm{Src}^{\mathrm{ur}}(\iota,e,A,W,g)$, given a modular tower, a
projection $\mathrm{pr}$ and a weight function $c_w\colon\Z/N\to\katoO$, not identically zero, as in the formal statement of
\cref{pf:kato:lit:proj}, applying $\mathrm{pr}$ to averaged Kummer classes of products of Siegel units of level $\ell^{r+n}N^2$
\cite{KubertLang1981} weighted by $c_w$ gives a family $Y=(Y_r)$ of cocycles of $G_\Q$ with values in functions on $(\Z/\ell^r)^\times$ (the
Shapiro form of cocycles of $G_{\Q(\zeta_{\ell^r})}$), continuous, integral and norm compatible in $r$, and a cocycle $y$ with
$y\katoeq r\sum_x\katorep x^{-e}\psi(x)Y_r(\cdot,x)$; and for some $m$, $\ell^my$ is an integral Euler system ($\mathrm{Sys}$),
$(\nu_D,\ell^mY_{\bullet,1})$ is an Iwasawa family at $A$ for the unramified quotient character $\nu_D$, and $\ell^mY_{\bullet,1}$ is linked to
$\ell^my$ (\cref{pf:kato:def:esdata}). It differs from \cref{pf:kato:lit:zeta} as follows: the classes live in the tower module instead of
\'etale cohomology; the Euler system relations are asserted only in the cyclotomic $\ell$-power direction and, at an auxiliary prime $w$,
in the form of the tame relations \ref{pf:kato:T1}--\ref{pf:kato:T4} for the twist by $x^{-e}\psi(x)$; and integrality is obtained by an
explicit power $\ell^m$, without a hypothesis on the image of Galois.
\end{remark}

\begin{littheorem}[{\cite[Prop.~2.4 and Prop.~8.12]{Kato2004}}]\label{pf:kato:lit:norm}
Let $M,N\ge1$, $q$ a prime not dividing $MN$, and $c,d$ integers with $(c,6Mq)=(d,6Nq)=1$. The norm map from $K_2(Y(Mq,Nq))$ to
$K_2(Y(M,N))$ sends the element ${}_{c,d}z_{Mq,Nq}$ to ${}_{c,d}z_{M,N}$ multiplied by the Euler factor
$1-T'(q)\bigl(\begin{smallmatrix}1/q&0\\0&1\end{smallmatrix}\bigr)^*+q\bigl(\begin{smallmatrix}1/q&0\\0&1/q\end{smallmatrix}\bigr)^*$
(see \cite[\S2]{Kato2004} for the normalisation of the dual Hecke operator $T'(q)$), and the $p$-adic zeta elements satisfy the
corresponding relation.
\formallabel{UpperTheory.V.Z.Kato.Nb.kato_norm_relation_tame_prime}
\end{littheorem}

\begin{remark}[The formal statement]\label{pf:kato:rem:norm}
The formal statement concerns the twisted Kummer classes of the tower attached to Siegel units: for a prime $p\nmid\ell Ncd$, $p$ times the
sum over the level-$Np$ classes is congruent modulo $\ell^n$, up to a coboundary for the action twisted by $\varepsilon_r^{-1}$, to
$pY_0(s,y)-T_pY_0(s,yp^{-1})+p^w\langle p\rangle Y_0(s,yp^{-2})$, where $Y_0$ is the level-$N$ class and $y\in(\Z/\ell^r)^\times$ is the
cyclotomic variable. It is the tame case of \cref{pf:kato:lit:norm}, with the Frobenius at $p$ acting on the cyclotomic variable, in the
tower module and modulo $\ell^n$.
\end{remark}

\begin{littheorem}[Explicit reciprocity law; {\cite[Thm.~9.6, Thm.~9.7, Prop.~10.10]{Kato2004}}]\label{pf:kato:lit:erl}
With the notation of \cite[8.1]{Kato2004}, let $1\le r\le k-1$ with $r=k-1$ or $r'=k-1$. The dual exponential map sends the image of
Kato's $p$-adic zeta element ${}_{c,d}z^{(p)}_{1,N,m}(k,r,r',\xi,S)$ in the local Galois cohomology of $\Q(\zeta_m)\otimes\Q_p$ to the zeta
modular form ${}_{c,d}z_{1,N,m}(k,r,r',\xi,S)\in M_k(X_1(N))\otimes\Q(\zeta_m)$, a product of two Eisenstein series; after projection
to $f$, the same holds for the zeta elements of $f$.
\formallabel{UpperTheory.V.Z.Kato.K5.smoothedZetaClass_hodgeTateLaw}
\end{littheorem}

\begin{remark}[The formal statement]\label{pf:kato:rem:erl}
The formal statement asserts: in the situation of the formal statement of \cref{pf:kato:lit:zeta}, with a linear form $\lambda$ as in
\cref{pf:kato:lit:proj}, for $d_0\ge2$ prime to $\ell$ the smoothed classes $Z_r(s,x)=Y_r(s,x)_1-d_0Y_r(s,xd_0)_1$ satisfy the explicit
reciprocity law of \cref{pf:kato:def:esdata}(e) with
$R_m(\psi,z)=(1-d_0\psi^{-1}(d_0))v_m\kappa\bigl(\sum_a\psi(a)z^a\bigr)\psi^{-1}(N)\,\lambda(\mathcal E_{m,\psi})$, where $\kappa\ne0$, the $v_m$ are
convergent units and $\mathcal E_{m,\psi}$ is the $q$-expansion of a $\psi$-twisted product of Eisenstein series. It is the analogue of
\cref{pf:kato:lit:erl} in the form of Coates--Wiles and Coleman \cite{CoatesWiles1977,Coleman1979}: instead of the dual exponential
map and $B_{\mathrm{dR}}$, it describes the restriction to $D_A$ of the zeta classes in the cyclotomic tower by Kummer presentations
by local units, and the Gauss sums of their logarithmic derivatives by $\lambda$ applied to the $q$-expansions of the zeta modular forms.
\end{remark}

\begin{littheorem}[Siegel units; {\cite[Prop.~1.3, Prop.~3.11]{Kato2004}}, {\cite[Ch.~2]{KubertLang1981}}]\label{pf:kato:lit:siegel}
Let $E$ be an elliptic curve over a scheme $S$ and $c$ an integer prime to $6$. There is a unique ${}_c\theta_E\in\OO(E\setminus E[c])^\times$
with divisor $c^2(0)-E[c]$ which is invariant under the norm maps attached to multiplication by integers prime to $c$. Its pull-backs
along torsion sections of the universal elliptic curve are the Siegel units ${}_cg_{\alpha,\beta}$, modular units with explicit
$q$-expansions, and $\mathrm{dlog}\,g_{\alpha,\beta}$ is an Eisenstein series of weight $2$.
\formallabel{GlobalTheory.T.siegelUnit_edgeClass_kummer_dlog_congruence}
\end{littheorem}

\begin{remark}[The formal statement]\label{pf:kato:rem:siegel}
The formal statement is the computation at the cusp (on the Tate curve) behind the formal form of \cref{pf:kato:lit:erl}. For $p$ odd
there is a constant $c_p$ with the following property. Let $a={}_cg$ and $b={}_dg$ be Siegel units of level $Np^L$ in a modular tower
($c,d>1$ prime to $6Np$), $x$ and $y$ the Kummer characters of $p^n$-th roots of $a$ and $b$, and $\eta$ a locally constant cochain on the
geometric Galois group whose coboundary is $x\cup y$. Then there is $\xi\colon\Q\to A$ such that, for every $s$ in the
decomposition group of $A$ fixing $\zeta_{p^L}$ and acting on $\zeta_{p^{L+n}}$ through $t\in\Z/p^n$, and every element $w$ built from
$(x,y,\eta)$ whose logarithmic $q$-expansion is $\phi$, the coefficients satisfy
$p^{c_p}\phi_r\equiv s^{-1}(\xi_r)-\xi_r-p^{c_p}\Delta\,t\,\Phi_r$ modulo $p^n$ for all $r\in\Q$. Here $\Phi$ is the $q$-expansion of the product of the
two smoothed weight-one Eisenstein-type functions attached to $a$ and $b$, and $\Delta=\alpha_1\beta_2-\alpha_2\beta_1$ if the indices of the two
Siegel units are congruent to $(\alpha_1,N\alpha_2)$ and $(\beta_1,N\beta_2)$ modulo $p^L$. It identifies the cyclotomic Kummer class of the symbol $\{a,b\}$ at the cusp with the product of the logarithmic
derivatives of $a$ and $b$, up to a coboundary and the factor $p^{c_p}$; compare the proof of \cite[Prop.~10.10]{Kato2004} and
\cite{Fukaya2003}. The published statements are not formulated modulo $p^n$ with an explicit loss.
\end{remark}

\begin{theorem}[Existence of the Iwasawa data]\label{pf:kato:gul}
The formal forms of \cref{pf:kato:lit:zeta} and \cref{pf:kato:lit:erl}, together with the existence of admissible Hecke projections
(\lref{GlobalTheory.T.exists_admissible_hecke_projection_nonvanishing}), imply \textup{(Z1)}.
\formallabel{UpperTheory.V.Z.Kato.Joint.kato_global_unit_law_of_smoothed_classes}
\end{theorem}

\begin{proof}
The formal proof (\lref{UpperTheory.V.Z.Kato.Joint.glue2}) proceeds as follows. Let $\mathrm{NF}(e,S,A,W)$ and
$\mathrm{Src}^{\mathrm{ur}}(\iota,e,A,W,g)$ hold, with character $\chi_g$ killing $\mu_\ell$. Choose a character $\psi$ modulo $\ell$
with $\chi_g(s)=\iota^{-1}\psi(\varepsilon(s)\bmod\ell)$ (\lref{UpperTheory.V.Z.Kato.Joint.psiEx}) and an embedding
$j\colon\Omega\to\katoQlb$ defining $A$ (\lref{UpperTheory.V.Z.Kato.Joint.jEx}). Apply the formal form of \cref{pf:kato:lit:eis}
to $g$, $\psi$ and $N_0=1$, take a tower by \cref{pf:kato:towerex}, and put $c_w(x)=\iota^{-1}(\hat\theta(x)/2)$, where $\hat\theta$ is a
finite Fourier transform of $\theta$. The formal form of \cref{pf:kato:lit:proj} gives $\mathrm{pr}$ and $\lambda$, and that of
\cref{pf:kato:lit:zeta} gives $Y$, $y$ and $m$. Apply the formal form of \cref{pf:kato:lit:erl} with $d_0=\ell+1$. An explicit computation
of the right-hand side (\lref{UpperTheory.V.Z.Kato.Joint.rhsEq}, \lref{UpperTheory.V.Z.Kato.Joint.avgEq},
\lref{UpperTheory.V.Z.Kato.Joint.meas3}) identifies $\lambda(\mathcal E_{m,\psi})$ with Gauss sums of a power series $L$ attached to an
Eisenstein measure, which gives the congruence of the pinning, with $\nu=\nu_D$ (trivial on $I_A$, and of norm one at a Frobenius element).
The non-vanishing of the moments of $L$ (\lref{UpperTheory.V.Z.Kato.Joint.pinOf}) follows from $v\ne0$ and from the non-vanishing of
$\sum_i\iota(\beta_i)\psi^{-1}(t_i)t_i^e$, where the $t_i$ and $\beta_i$ are the smoothing integers and coefficients built from $c_1,d_1,\ell+1$;
this sum equals $1-(\ell+1)^{e+1}\ne0$ times a sum which is non-zero by the choice of $c_1,d_1$
(\lref{UpperTheory.V.Z.Kato.Joint.smNz}). Finally the smoothing is compatible with the other properties: if $\mathrm{Sys}$,
$\mathrm{Iw}$ and $\mathrm{Sp}$ hold for $(y,Z)$, they hold for $\bigl((1-d_0^{e+1})y,\,Z_r(s,x)-d_0Z_r(s,xd_0)\bigr)$ when $d_0\equiv1\bmod\ell$
(\lref{UpperTheory.V.Z.Kato.Joint.smSys}). This gives data as in \textup{(Z1)}.
\end{proof}

\begin{theorem}[Existence of the Euler system]\label{pf:kato:zetaur}
\textup{(Z)} holds.
\formallabel{UpperTheory.V.Z.Kato.Final.zetaUr}
\end{theorem}

\begin{proof}
The formal forms of \cref{pf:kato:lit:zeta} and \cref{pf:kato:lit:erl} are theorems of the development, and admissible Hecke projections
exist by \lref{GlobalTheory.T.exists_admissible_hecke_projection_nonvanishing}; so \textup{(Z1)} holds by \cref{pf:kato:gul}.
\textup{(Z2)} holds by \cref{pf:kato:plocu}, and \cref{pf:kato:halves} gives \textup{(Z)}.
\end{proof}

\begin{remark}
The non-triviality of the Euler system does not come from the projection $\mathrm{pr}$, which is not required to be non-zero, but from the
explicit reciprocity law and the non-vanishing of the moments of $L$ (\cref{pf:kato:plocu}), as in \cite{Kato2004}.
\end{remark}

% ---------------------------------------------------------------------------
\subsection{Proof of the main theorem}\label{pf:kato:sec:proof}

\begin{theorem}\label{pf:kato:main2}
$\sfname{HypothesesH'}$, $\sfname{Row0b}$ for the local parameter maps of \cref{sec:galois}, and \textup{(QBC)} imply
$\sfname{KatoTwoW}$.
\formallabel{UpperTheory.V.Z.Kato.WeilSphericalRecords.katoTwoW_of_printsH_of_rT_of_quadBC2}
\end{theorem}

\begin{proof}
Fix the data of $\sfname{KatoTwoW}$: $\sfname{Used}(\ell,K,e,N,\theta,\psi,W)$, $\iota$, a newform $f$ attached to $W\otimes\varepsilon^{-e}$,
and $\sfname{WitAQ}(\iota,W)$. By \cref{pf:kato:locq}, \textup{(LG)} holds. By \cref{pf:kato:usedadm}, $\det W=\theta$ and $W$ is admitted.
By \cref{pf:kato:twnew}, \textup{(TN)} holds; by \cref{pf:kato:frobroot}, \textup{(FR)} holds; hence \textup{(S)} holds by
\cref{pf:kato:srcs}. By \cref{pf:kato:twos}, $\sfname{KatoTwoW}^{\mathrm{sph}}$ holds, and it remains to check $\mathrm{Sph}(\iota,e,W)$. By
\cref{pf:kato:twwit}, \textup{(TW)} holds, so \textup{(SW)} holds by \cref{pf:kato:sphofwit}. Its hypotheses are satisfied: $W$ is
continuous; $W$ is irreducible, since it is irreducible on $H_0(W)$ (\lref{GlobalTheory.Ord2.IrrTransfer.isIrr_of_range}: a
representation whose image contains the image of an irreducible representation is irreducible); $\det W$ kills some $\mu_M$
(\cref{pf:kato:keropen}); and $\sfname{WitAQ}(\iota,W)$ holds. So $\mathrm{Sph}(\iota,e,W)$ holds.
\end{proof}

\begin{proof}[Proof of \cref{pf:kato:main}]
Apply \cref{pf:kato:main2}, with $\sfname{Row0b}$ taken from the statement on Galois representations contained in
$\sfname{HypothesesH'}$ (\lref{UpperTheory.Gal.SV.rowOf'}: that statement is $\sfname{Row0b}$ for the local parameter maps of
\cref{sec:galois}) and with \textup{(QBC)} given by \lref{UpperTheory.ConvolutionTheory.R2.quadBC2}, a theorem proved
elsewhere in the development by a theta lifting argument (it shows: for $F$ totally real, $E/F$ totally complex quadratic and
$\pi$ regular algebraic cuspidal on $\GL_2(\A_F)$, either $\pi\cong\pi\otimes\eta_{E/F}$ or $\pi$ has a cuspidal base change to $E$).
\end{proof}

\begin{remark}[On the formal statement]\label{pf:kato:rem:final}
(1) As noted in \cref{pf:kato:rem:unused}, many conditions of $\sfname{Used}$ are not used; the proof shows the vanishing for every
admitted $W$ for which $\sfname{WitAQ}(\iota,W)$ holds and the twists of $W$ have the stated automorphic properties. (2) The
conclusion concerns Selmer groups with $\katoQlb$-coefficients and the crystalline condition of \cref{pf:kato:rem:bcris}; it is
the vanishing of rational Selmer groups, not finiteness or a bound for integral ones (compare \cref{sec:window:kato}). (3) The
only inputs from the literature are the statements of $\sfname{HypothesesH'}$; Kato's theorem itself is not assumed, and the
Euler system, the explicit reciprocity law and the Euler-system argument are proved in the development. (4) The Euler-system
argument of \cref{pf:kato:sec:heights} uses one auxiliary prime at a time and is not a transcription of a published proof.
\end{remark}
